\documentclass[carta,firma,sinnotas]{lafrox}
\usepackage{fvextra}
\usepackage{pdflscape}

\datoslafrox{
  documento = {Subdocumento 8: Anexos},
  subtitulo = {Anexos 8.A--8.H},
  alcance = {Gestión de riesgos de la propuesta LafroX},
  formulario = {T-16: Plan de riesgos},
  fecha = {25 de noviembre de 2026},
  version = {0.1},
  nomenclatura = {LAFROX-Subdocumento8-Anexos},
}
% Soporte local de firma para el cuerpo y los anexos del Subdocumento 8.
% No altera la clase ni las plantillas compartidas.
\newcommand{\LFXportadasd}{%
  \AddToHookNext{shipout/foreground}{%
    \begin{tikzpicture}[remember picture,overlay]
      \node[anchor=north west,inner sep=0pt]
        at \LFXen{\LFXmi}{0.835}
        {\includegraphics[height=13mm]{logo/FIRMA_Alex_Aravena_transparente.png}};
    \end{tikzpicture}%
  }%
  \portadalafrox
}

% Firma abreviada sobre filete, en el mismo espacio del pie existente.
\renewcommand{\LFXpiederecho}{%
  \begingroup\LFXpiebase
  \begin{tabular}[b]{@{}r@{}}
    \raisebox{0pt}[0pt][0pt]{%
      \begin{tabular}[b]{@{}c@{}}
        \includegraphics[height=7mm]{logo/FIRMA_Alex_Aravena_transparente.png}\\[-1pt]
        \rule{3.4cm}{0.4pt}
      \end{tabular}%
    }\\[-1pt]
    {\LFXchico Media firma: \LFXrepresentante}\\
    {\intermedium\fontsize{9pt}{11pt}\selectfont\color{lafroxTinta}%
      \LFXfoliotexto}
  \end{tabular}%
  \endgroup
}


% Formato de portada común, basado en el Subdocumento 1.
\configurarportadalafrox{Anexos del Subdocumento 08}{Anexos del Subdocumento 08}{Anexos 8.A--8.H}{LAFROX-Subdocumento8-Anexos.pdf}{logo/FIRMA_Alex_Aravena_transparente.png}
\begin{document}
\portadalafrox
\preliminareslafrox
\setcounter{chapter}{7}

\chapter{Anexos del subdocumento 8}

Este documento reúne los anexos del plan de riesgos de LafroX.

\section{Anexo 8.A: Matriz de respuestas propuestas}

La Tabla 8.A.1 presenta una respuesta propuesta para cada uno de los 36 riesgos, con sus disparadores y responsables. La Tabla 8.A.2 especifica para los 29 riesgos con coordinación del CLIENTE la contraparte, solicitud, evidencia, plazo y contingencia por falta de respuesta. Sus objetivos de coordinación se proponen para acuerdo en el arranque; no acreditan una aceptación previa ni la eficacia de las respuestas. La estimación de esfuerzo por perfil, los plazos completos y la valoración residual proyectada están en 8.F.1--2; el costo-beneficio de las respuestas altas/críticas, en 8.F.3. La subcontratación no transfiere la obligación contractual final de LafroX.

\begingroup
\renewcommand{\thetable}{8.A.1}
\begin{landscape}
\setlength{\textwidth}{\linewidth}
\begin{tablalafrox}{Matriz de respuestas propuestas para los 36 riesgos vigentes. Fuente: riesgos, disparadores y responsables propuestos en el Formulario T-16; respuestas propuestas en este plan.}{tab:sd8-anexo-a-respuestas}{G{0.06} F{0.22} Y F{0.18}}{\cab{ID} & \cab{Disparador observable} & \cab{Respuesta propuesta y actuación ante el disparador} & \cab{Responsable propuesto en T-16}}
R-001 & Cadena del alcance sin certificación a D-5 del inicio de los 28 días finales de evidencia de E2 & Acordar y registrar con Comercial, TI y cada cadena el perfil de mensajes, el acceso a pruebas y la agenda de certificación en el levantamiento 2.1.5, previo al inicio de 5.3.1. Ventanas de certificación: cadena principal (5.3.1), del 01-10-2027 al 31-05-2028; demás cadenas (5.3.2), del 02-05-2028 al 14-07-2028. Reconfirmar la agenda en los meses 13--14 si se requiere. JP escala por CT ante falta de respuesta o de certificación a D-5 del inicio de las cuatro semanas finales de evidencia de E2. Continúan las pruebas independientes; la carga manual controlada mantiene la continuidad y la activación EDI requiere certificación. & Arquitecto de Solución \\
R-002 & Lotes o vínculos históricos sin evidencia de origen al perfilar la carga & Solicitar a la Jefa de Calidad, con TI y Bodega, evidencia de origen y decisión sobre lotes no reconstruibles; política V-07 al cierre del mes 2 y excepciones a D-5 de validar cada carga. Si no responde, JP escala por CT, se aísla el conjunto y queda pendiente su aceptación; no se inventa historia sanitaria. & Líder de Datos \\
R-003 & V-09 no acredita tres años utilizables o revela períodos incompletos & Solicitar al Jefe de Abastecimiento y a TI el historial de tres años y validación de V-09 al cierre del mes 3; corregir extractos a D-5 de validar la carga. Ante silencio, escalar por CT, separar períodos incompletos y continuar el perfilado disponible sin reducir la historia exigida. & Líder de Datos \\
R-004 & Diferencias sin causa ni ajuste aprobado en el acta previa al corte & Solicitar a Bodega, con Operaciones y Finanzas, protocolo de conciliación en mes 2, validado en mes 3. Confirmar conteo a D-5 del corte y cerrar acta y ajustes antes de D. Si no hay respuesta, escalar por CT, repetir conciliación y bloquear el corte afectado hasta contar con evidencia aprobada. & Líder de Datos \\
R-005 & No existe saldo inicial aceptado por cliente y transportista & Solicitar al Gerente de Operaciones y a Finanzas respaldos y conciliación V-08 de saldos por cliente y transportista; política en mes 2 y respuesta a D-5 de validar la carga M8 de cada ola. Ante silencio, escalar por CT y mantener M8 pendiente para esos saldos; no presumir saldo cero. & Líder de Datos \\
R-006 & V-10 no obtiene horas o montos históricos verificables & Solicitar a Comercial y Operaciones las horas históricas y a Finanzas los montos de rendición, con validación de V-10 al cierre del mes 1. Si no responden, escalar por CT y mantener pendiente la referencia de aceptación; la captura futura no reconstruye la historia anterior. & Líder Funcional en el mes 1 (V-10); después, Líder de Implantación y Gestión del Cambio \\
R-007 & Documentos todavía vigentes quedan fuera del volumen histórico previsto & Solicitar a Finanzas y TI el universo de deuda vigente, con respuesta inicial al cierre del mes 3 y conciliación a D-5 de validar la migración. Ante silencio, escalar por CT y mantener la cartera incompleta pendiente; preparar extracción complementaria sin omitir deuda obligatoria. & Líder de Datos \\
R-008 & Trámite tributario permanece pendiente de respuesta externa & Mantener trazabilidad de trámites pendientes, consulta y reintentos idempotentes hasta recibir confirmación del SII. Ante demora externa, conservar el estado pendiente y escalar según el procedimiento acordado; el silencio no equivale a aprobación. & Líder de Operación / SRE \\
R-009 & Calidad no aprueba parámetros por producto al preparar certificación & Solicitar a la Jefa de Calidad la matriz térmica versionada V-11 durante mes 1 y su aprobación al cierre del mes 2, antes de certificar. Si no responde, escalar por CT y mantener el tratamiento preventivo de lecturas fuera de ficha; no cerrar marcha blanca ni aprobar parámetros por el CLIENTE. & Líder de Calidad \\
R-010 & Equipo de terreno especificado no recibido antes de instalar o probar & Solicitar al Jefe de TI, con Operaciones, calendario de equipos en mes 1, entrega confirmada a D-5 y recepción técnica antes de la instalación o prueba D. Ante silencio o faltante, escalar por CT, continuar trabajo independiente y bloquear las actividades sin equipo. & Líder de Implantación y Gestión del Cambio \\
R-011 & Acuerdo operativo de un transportista pendiente antes de su ola & Solicitar al Gerente de Operaciones y al representante de cada transportista el acuerdo y la nómina desde mes 2, firmados a D-5 de su incorporación a la ola. Ante silencio, escalar por CT y no habilitar a esa empresa; conservar su cobertura pendiente sin retirarla del alcance. & Líder de Implantación y Gestión del Cambio \\
R-012 & Propietario niega acceso o autorización de instalación del termógrafo & Solicitar a Operaciones y al propietario permiso de instalación por vehículo con D-10 de la visita y obtenerlo a D-5. Ante silencio, escalar por CT, suspender la intervención y proponer otra visita o vehículo autorizado; mantener pendiente la evidencia térmica. & Líder de Implantación y Gestión del Cambio \\
R-013 & Transacciones se siguen registrando solo por el procedimiento anterior & Solicitar a Comercial y Operaciones, con Bodega por turno, responsables y agenda de acompañamiento a D-5 del inicio de cada ola, convocando a D-10. Si no responden o persiste el registro externo, escalar por CT, focalizar acompañamiento y mantener pendiente la evidencia de adopción. & Líder de Implantación y Gestión del Cambio \\
R-014 & Planificador deja de estar disponible antes de validar reglas pendientes & Solicitar al Gerente de Operaciones y al planificador confirmación de retiro y disponibilidad V-14 al cierre del mes 1; transferir y validar reglas con el ayudante en meses 1--3. Probar las rutas y excepciones sin el planificador en la marcha blanca de meses 13--15 (S7, 9.6.1). Si falta respuesta o se retira antes, escalar por CT y reconstruir con ayudante e historia verificable, sin dar por validadas reglas pendientes. & Líder de Implantación y Gestión del Cambio \\
R-015 & Usuarios convocados no completan formación y evaluación previstas & Solicitar a Operaciones, Bodega, Comercial y TI, según perfil, nómina y liberación horaria a D-5 de cada sesión convocada a D-10. Ante silencio o ausencia, escalar por CT y reprogramar evaluación; usar sólo suplentes certificados y no acreditar preparación ni cierre de ola sin evidencia. & Líder de Implantación y Gestión del Cambio \\
R-016 & Decisión requerida no llega en la fecha acordada y bloquea la tarea & Solicitar a CT área decisora y suplente, con agenda al cierre del mes 1; cursar insumos a D-10 y pedir respuesta a D-5. Las aceptaciones formales conservan diez días hábiles. Para promover la región de contingencia, avisar al suplente a los 15 min sin respuesta del titular y obtener decisión dentro de 30 min (S4, 4.3.2.5); no esperar el plazo ordinario. Mantener bloqueada la acción sin autorización; el silencio no permite cambios. & Jefe de Proyecto \\
R-017 & Corte del enlace del sitio o del respaldo restante; en Los Ángeles de madrugada, pérdida del camino Starlink disponible. & Mitigar. Verificar respaldos antes de aceptar la ola de Concepción; no considerar disponible un enlace sólo por estar ofertado. Mantener y ensayar los enlaces ya ofertados: en Talca y Concepción, fibra y LTE con Starlink en espera caliente; en cross-docking, Starlink con LTE Entel/Claro y conmutación menor a 30 s. En Los Ángeles de madrugada, probar pérdida del camino Starlink restante y continuidad local de su ventana de 3 h. Coordinar restitución y conciliación sin duplicar el mismo corte como falla de SII, POS o mapas. Referencia: S4, 4.2.5.3. & Líder de Operación / SRE \\
R-018 & Corte eléctrico excede autonomía del respaldo aplicable & Verificar autonomía, alarmas y procedimiento de recuperación de los componentes de LafroX, coordinando las dependencias de energía con el CLIENTE/proveedor. Si se excede la autonomía, activar recuperación y registrar la indisponibilidad e impacto. & Líder de Operación / SRE \\
R-019 & Cliente no otorga aprobación de tratamiento en las regiones propuestas & Solicitar por CT, con TI y la autoridad del CLIENTE encauzada por Gerencia General, aprobación escrita de residencia y transferencia en mes 1, resuelta al cierre del mes 2 o antes de habilitar datos, lo que ocurra primero. Ante silencio, escalar y preparar alternativa de arquitectura; mantener deshabilitada la transferencia no aprobada. & Arquitecto de Solución \\
R-020 & Intervención urgente requiere autorización y esta no se obtiene & Solicitar a CT, Operaciones y TI flujo y suplentes al cierre del mes 1; ante incidente, pedir autorización escrita antes del cambio. Sin respuesta, escalar en una hora o antes si lo exige el plazo de atención, manteniendo contención autorizada y estado seguro. El silencio no permite intervenir durante congelamiento. & Encargado de Seguridad de la Información \\
R-021 & Falla mecánica impide completar el recorrido de prueba & Registrar los casos TI pendientes y conservar la evidencia disponible. Coordinar con Operaciones del CLIENTE la repetición de los casos afectados cuando el transportista habilite otro recorrido. & Líder de Implantación y Gestión del Cambio \\
R-022 & Siniestro impide continuar los casos de prueba previstos & Suspender la prueba TI y registrar los casos interrumpidos. Con el recorrido alternativo autorizado por Operaciones del CLIENTE y el transportista, completar los casos pendientes. & Líder de Implantación y Gestión del Cambio \\
R-023 & Asalto interrumpe recorrido y evidencia prevista & Preservar la evidencia TI disponible e identificar los casos que requieren repetición. Reanudar la validación en una ventana habilitada por Operaciones del CLIENTE y el transportista. & Líder de Implantación y Gestión del Cambio \\
R-024 & Conductor sustituto no figura autorizado/certificado al presentarse & Solicitar a Operaciones y al transportista nómina y suplentes a D-5 de incorporar la empresa y confirmar conductor, vehículo y viaje antes de cada salida. Operaciones y el transportista buscan un conductor certificado para el mismo turno. Si no hay reemplazo, Operaciones redistribuye las entregas entre rutas habilitadas con capacidad o reprograma las restantes, avisa a los destinatarios y registra la decisión por entrega. El supervisor verifica al conductor antes de autorizar la salida. & Líder de Implantación y Gestión del Cambio \\
R-025 & Solicitud nueva excede requisitos o exclusiones aprobados & Solicitar por CT a la gerencia peticionaria solicitud escrita y decisión del cambio antes de su primera actividad dependiente; registrar y remitir el mismo día hábil. Ante silencio, mantener línea base y continuar trabajo independiente; si amenaza un hito, escalar al Comité Ejecutivo sin iniciar desarrollo no aprobado. & Jefe de Proyecto \\
R-026 & Cambio aplicable obliga a evaluar adecuación de controles o retención & Revisar cambios normativos aplicables, determinar su efecto con las áreas competentes y tramitar ajustes de controles, retención o eliminación mediante decisión y control de cambios documentados, sin posponer una obligación legal aplicable. & Arquitecto de Solución \\
R-027 & Diferencia de cálculo o reclamo en una liquidación mensual & Solicitar a Finanzas, con Comercial y CT, la aprobación de la fórmula inicial de 7.4.1 (fecha nominal: mes 17; a más tardar al cierre del mes 20, fin tardío del paquete) y tres liquidaciones en sombra reproducidas en meses 21--23 antes del pago variable del mes 24. Ante silencio, no activar el tramo; una disputa se escala y suspende sólo el devengo correspondiente según contrato, manteniendo el servicio. & Jefe de Proyecto \\
R-028 & Acuerdo de finalidades y controles pendiente a D-5 de la ola, u objeción formal al uso operacional propuesto. & Mitigar y escalar. JP coordina con Operaciones, TI y sindicato el acuerdo 9.7.2 antes de la ola: convocar a D-10 y resolver a D-5. Explicar finalidades, acceso y controles; mantener excluidas las cámaras en cabina y no usar GPS para jornada sin acuerdo previo. Ante desacuerdo, escalar por CT, mantener sólo el uso operacional autorizado y registrar la cobertura de la ola pendiente. & Jefe de Proyecto \\
R-029 & Preparador nuevo asignado sin tutor o evaluación, o aprendizaje que excede dos horas en el programa 9.5.2. & Mitigar. Mantener 9.5.2 en meses 13--56, turno 22:00--06:00, tutor y evaluación antes del trabajo autónomo, con aprendizaje máximo de dos horas. Bodega y Operaciones entregan altas por turno antes de asignarlas. IMP conduce hasta mes 21 y SRE desde 22, familia IMP. Ante sobrecarga, usar personal certificado, reforzar la inducción y escalar por CT; no duplicar las HH ordinarias de 9.5.2. & Líder de Implantación y Gestión del Cambio \\
R-030 & No hay contrato de interfaz o extracción verificado al cierre del mes 1, o el ensayo temprano contradice formato, acceso o comportamiento declarados. & Mitigar y escalar. ARQ solicita a TI acceso, diccionario y condiciones de extracción en mes 1 (2.1.5, 5.1.1 y V-16); DAT verifica la migración del WMS. Probar temprano sin modificar el ERP. Si falta una interfaz, evaluar intercambio de archivos controlado, sujeto a demostrar los requisitos; elevar por CT el bloqueo y sus impactos. La alternativa no acredita por sí sola latencia, integridad ni emisión tributaria. & Arquitecto de Solución \\
R-031 & Incidente alto/crítico, evidencia o certificación pendiente en las semanas finales del peak, sin ventana de intervención o validación disponible. & Mitigar y escalar. JP acuerda con Operaciones y CT disponibilidad y continuidad antes de septiembre; cerrar habilitaciones en agosto según S7. En meses 19--20 mantener evidencia diaria de las seis condiciones de marcha blanca. Un defecto durante 1--25 de septiembre no autoriza despliegue: activar continuidad aprobada y R-020 si requiere autorización. Si persiste una condición incumplida, no declarar cierre; elevar plan de recuperación sin mover hitos por cuenta propia. & Jefe de Proyecto \\
R-032 & Aviso de soporte/compatibilidad que deja una versión necesaria fuera de la política aprobada o hace fallar la regresión de la actualización prevista. & Mitigar. ARQ mantiene inventario de versiones y revisa soporte y compatibilidad trimestralmente desde mes 1, con DES y SRE. Preparar y probar adecuación antes del límite verificado de soporte, dentro de ventanas autorizadas. Si el cambio excede la actualización ordinaria 10.2.5, registrar el esfuerzo extraordinario y elevar impactos por CT; usar reversión ensayada y soporte permitido, sin sostener una versión insegura por conveniencia. & Arquitecto de Solución \\
R-033 & Ensayo de exportación/restauración incompleto, acceso o clave no disponible para el CLIENTE, o cambio de condiciones que impide cumplir el plan de salida. & Mitigar. ARQ verifica con TI propiedad de cuentas, acceso a datos, claves, exportaciones e infraestructura como código en la recepción inicial y en cada actualización anual del plan 1.6.3. Probar una restauración utilizable, no sólo descargar archivos. Ante restricción, documentar brecha, mantener continuidad autorizada y escalar por CT alternativas y decisión contractual; no declarar reversibilidad completa ni cambiar de proveedor sin prueba y aprobación. & Arquitecto de Solución \\
R-034 & Carga sin guía vigente que requiere nueva emisión mientras el ERP está indisponible; pérdida de sala con ERP aún no restituido. & Mitigar y escalar. ARQ/SRE acuerdan con TI y Finanzas restitución del ERP y ensayan AL-DTE-01/AL-DR-01 antes de H5/H10, incluida pérdida de sala. Mantener sólo salidas con guía preemitida aún válida y bloquear las que requieren nueva guía. Recuperar WMS en nube no restaura el ERP ni exceptúa el RTO requerido; registrar y escalar por CT cualquier incumplimiento, sin segundo emisor ni aprobación tributaria por silencio. & Arquitecto de Solución \\
R-035 & Conductor asignado y habilitado rechaza el flujo digital o entrega sin la evidencia requerida, sin que la causa sea firma pendiente o reemplazo sin aviso. & Mitigar y escalar. IMP coordina con Operaciones, transportista y conductor demostración y acompañamiento antes de su salida, dentro de 9.5.3. Confirmar disponibilidad a D-5 de la ola y aceptación práctica en andén. Ante rechazo, solicitar a la empresa conductor certificado y dispuesto, mantener acompañamiento y elevar por CT la cobertura pendiente; no atribuir a LafroX facultades laborales ni declarar completa la evidencia faltante. & Líder de Implantación y Gestión del Cambio \\
R-036 & Mediciones de concentración temporal o mezcla de transacciones distintas del perfil validado, o perfil sin confirmar al cierre del mes 1. & Mitigar. ARQ solicita el perfil a Operaciones/TI en los primeros cinco días hábiles del mes 1 y lo valida al cierre de ese mes; revalida antes de H5/H10, preparación del peak y cambios del perfil. SRE asume gestión operativa desde mes 22, familia ARQ. Ante divergencia, conservar mediciones, ejecutar CT-14, bloquear promoción afectada y elevar recuperación por CT antes del hito. El peak conocido, crecimiento comprometido y corrección del diseño dentro del alcance siguen a cargo de LafroX. & Arquitecto de Solución \\
\end{tablalafrox}
\end{landscape}
\endgroup

\subsection*{Dependencias del CLIENTE: compromisos y contingencias}
\addcontentsline{toc}{subsection}{Dependencias del CLIENTE: compromisos y contingencias}

Esta matriz establece los compromisos de coordinación de los 29 riesgos que requieren insumos o decisiones del CLIENTE. La contraparte aporta antecedentes y decide dentro de su autoridad; LafroX conserva la responsabilidad de diseño, capacidad, seguridad y mantenimiento. La gestión de servicios externos, suministros, siniestros y cambios normativos se realiza mediante las respuestas específicas de la Tabla 8.A.1. Los plazos de coordinación se proponen para acuerdo en el arranque; no acreditan respuestas recibidas ni modifican las obligaciones contractuales.

CT es la Contraparte Técnica. En ``ID / dueño'', JP, ARQ, DAT, CAL, IMP y SEG corresponden a las familias y responsables de la Tabla 8.1, que distingue los relevos operativos de las suplencias. D identifica la actividad o ventana indicada de la línea base; D-10 y D-5 son días hábiles anteriores. Salvo indicación específica, el dueño envía expediente completo a D-10 y solicita respuesta a D-5; si existe un plazo de mes más temprano, rige ese plazo. Para mes 1 solicita en sus primeros cinco días hábiles. Los plazos formales del art. 18.3 se conservan.

Si falta respuesta, el dueño informa a JP y CT el mismo día hábil del vencimiento. JP pide encauzamiento en un día hábil y, si no se resuelve en dos días hábiles del vencimiento, solicita instancia ejecutiva extraordinaria con la gerencia competente. Las amenazas a seguridad, continuidad o hitos se escalan de inmediato. La columna final añade la contingencia específica; no se espera la reunión mensual para actuar. Toda respuesta se cierra con evidencia validada, no con una promesa o el envío de la solicitud. Los titulares y suplentes por cargo se registran en mes 1 y se actualizan al cambiar su designación.

\begingroup
\renewcommand{\thetable}{8.A.2}
\begin{landscape}
\setlength{\textwidth}{\linewidth}
\begin{tablalafrox}{Contrapartes, solicitudes, plazos y contingencias de los 29 riesgos con coordinación del CLIENTE. Fuente: asignaciones de S2/S3/S6/S13 y objetivos de coordinación propuestos en S8, 8.1.2.}{tab:sd8-dependencias-cliente}{G{0.07} F{0.17} Y F{0.21} F{0.23}}{\cab{ID / dueño} & \cab{Contraparte del CLIENTE / apoyo} & \cab{Solicitud y evidencia de resolución} & \cab{Plazo de solicitud y respuesta} & \cab{Escalamiento y contingencia por falta de respuesta}}
R-001 / ARQ & Gerente Comercial; coordinación de TI con el responsable EDI de cada cadena. & Registrar cadenas del alcance, contacto técnico, perfil de mensajes, acceso a pruebas y agenda. Entregar resultado firmado del intercambio de pedidos y avisos por cadena. & Acordar perfil, acceso y agenda en 2.1.5, previo al 01-10-2027. Certificar a la principal en 5.3.1 (01-10-2027 al 31-05-2028) y a las demás en 5.3.2 (02-05-2028 al 14-07-2028). Reconfirmar en meses 13--14 si se requiere. D-5 es el umbral residual; D es el inicio de los 28 días finales de evidencia de E2. & JP escala al Gerente Comercial por CT ante falta de respuesta o de certificación a D-5. Continúan las pruebas independientes y la carga manual controlada sostiene la continuidad. La activación EDI de cada cadena requiere certificación. \\
R-002 / DAT & Jefa de Calidad; Jefe de TI y Jefa de Bodega aportan fuentes ERP/WMS y antecedentes de lotes. & Entregar extractos y evidencia de origen de los vínculos sanitarios; Calidad debe decidir por escrito la disposición de conjuntos cuya historia no pueda reconstruirse y la política ante lote desconocido (V-07). & Solicitar fuentes en mes 1; política y resolución inicial al cierre del mes 2. Para cada carga, resolver las excepciones a D-5 de su validación; nunca después del corte. & Escalar a CT con los lotes y evidencia faltante; aislar el conjunto y trabajar sobre datos verificables. Mantener la aceptación pendiente; no inventar lotes ni reconstruir historia sin fuente. \\
R-003 / DAT & Jefe de Abastecimiento; Jefe de TI entrega y explica los extractos del ERP. & Entregar tres años de ventas y pedidos, diccionario y cobertura temporal; validar el informe de perfilado V-09 y las excepciones identificadas. & Solicitar en mes 2; validación de V-09 al cierre del mes 3 y correcciones de cada extracto a D-5 de validar su carga. & Escalar por CT a Abastecimiento y TI; separar períodos incompletos y continuar perfilado de fuentes disponibles. No reducir unilateralmente los tres años exigidos ni dar por aceptado el historial faltante. \\
R-004 / DAT & Jefa de Bodega; Gerente de Finanzas aprueba ajustes de saldo y Gerente de Operaciones coordina conteos. & Entregar conciliación entre inventario legado y conteo físico, causas de diferencias y ajustes autorizados; presentar acta de validación previa al corte. & Solicitar protocolo en mes 2 y validarlo en mes 3. En cada corte, responsables y conteo confirmados a D-5; conciliación y ajustes firmados antes de D. & Escalar por CT a Bodega y Finanzas; repetir conteo y conciliación del conjunto afectado. No ejecutar su corte sin acta; continuar trabajo independiente y evaluar recuperación sin mover hitos por silencio. \\
R-005 / DAT & Gerente de Operaciones; Gerente de Finanzas valida saldos y representantes de transportistas aportan antecedentes. & Entregar saldos de envases por cliente y transportista y sus respaldos; validar V-08. Si se requiere conteo inicial, autorizar su ejecución y firmar conciliación de M8. & Solicitar política de saldos en mes 2 (S-10). Para cada ola de reparto: respuesta a D-5 de validar la carga M8 y conciliación aprobada antes de habilitarla. & Escalar a Operaciones y Finanzas por CT; preparar conteo o reconstrucción documental autorizados. Mantener pendiente M8 en los saldos afectados; no asignar cero ni presumir aceptación. \\
R-006 / IMP & Gerentes Comercial y de Operaciones: horas de entrega; Gerente de Finanzas: montos esperados de rendición. & Entregar registros históricos verificables y aprobar fuente, período y confiabilidad de la línea base V-10; identificar limitaciones y resolver formalmente una alternativa que preserve los criterios de aceptación. & Solicitar en los primeros cinco días hábiles del mes 1; validar V-10 al cierre de ese mes. Cualquier alternativa debe resolverse antes de iniciar la comparación de aceptación. & Escalar a las tres gerencias por CT. Conservar por separado la captura futura y la historia comprobada; suspender la comparación afectada hasta resolver la referencia, sin fabricar antecedentes ni cambiar la meta. \\
R-007 / DAT & Gerente de Finanzas; Jefe de TI ejecuta la extracción del ERP y documenta su cobertura. & Entregar listado completo de deuda vigente y conciliación con el extracto provisional; autorizar la inclusión de documentos omitidos y registrar volumen y diferencias. & Solicitar universo y criterio de vigencia en mes 2; respuesta inicial al cierre del mes 3. Cada extracto completo y conciliado a D-5 de validar la migración de cartera. & Escalar por CT a Finanzas y TI; medir y preparar la extracción complementaria. Mantener pendiente la cartera incompleta; tramitar impactos por el gobierno del proyecto sin tratar la integridad obligatoria como ampliación opcional. \\
R-009 / CAL & Jefa de Calidad, autoridad de aprobación sanitaria; Bodega y Operaciones validan aplicación en terreno. & Aprobar y versionar matriz de umbrales y duraciones por tipo de producto, fichas de almacenamiento y reglas de liberar, bloquear o rechazar; emitir aprobación de V-11. & Solicitar durante mes 1; matriz aprobada al cierre del mes 2, antes de certificar parametrización. Todo cambio posterior debe aprobarse antes de probarlo o usarlo. & Escalar inmediatamente por CT si amenaza certificación. Aplicar el tratamiento preventivo de V-11 a lecturas fuera de ficha y conservar el bloqueo; no cerrar marcha blanca sin matriz aprobada. \\
R-010 / IMP & Jefe de TI, como coordinador técnico de adquisición; Gerente de Operaciones confirma disponibilidad por sitio y ola. & Confirmar orden y fecha de entrega del equipamiento comprado por el CLIENTE, inventario por ola y evidencia de recepción técnica conforme a la especificación de LafroX. & Solicitar calendario en mes 1. En cada instalación, entrega y acceso confirmados a D-5 y recepción técnica antes de D; D es el inicio de instalación o prueba que requiere el equipo. & Escalar por CT al Jefe de TI y a Operaciones; ejecutar sólo trabajo independiente. Bloquear instalación y pruebas sin equipo; evaluar recuperación sin sustituir la compra del CLIENTE por una obligación de LafroX. \\
R-011 / IMP & Gerente de Operaciones; representante autorizado de cada empresa transportista firma su acuerdo. & Obtener acuerdos operacionales por empresa con nómina y coordinación de conductores, vehículos, acceso y capacitación; registrar cobertura de las diez empresas sin presumir firma. & Solicitar acuerdos en mes 2. Para cada ola de reparto, firma y nómina a D-5 del inicio de incorporación de la empresa; verificar antes de habilitar a sus conductores. & Escalar por CT a Operaciones; incorporar sólo empresas con acuerdo y registrar cobertura pendiente. Una ola sin esa empresa no acredita cobertura completa ni elimina obligaciones de aceptación. \\
R-012 / IMP & Gerente de Operaciones gestiona acceso; propietario o representante autorizado del camión otorga permiso. & Entregar autorización de acceso e instalación de termógrafo por vehículo, agenda de visita y condiciones del propietario; documentar cualquier alternativa sometida a aprobación. & Solicitar con D-10 y recibir autorización a D-5 de cada visita de instalación. Comprobar permiso y vehículo antes de intervenir en D. & Escalar a Operaciones por CT; cancelar la intervención no autorizada y proponer otra visita o vehículo autorizado con trazabilidad equivalente. Mantener pendiente la evidencia térmica; no instalar ni sustituir controles sin aprobación. \\
R-013 / IMP & Gerentes Comercial y de Operaciones; Jefa de Bodega para el turno nocturno y supervisores de cada ola. & Confirmar supervisores, turnos y usuarios que registrarán transacciones oficiales; aprobar agenda de acompañamiento y revisar evidencia diaria de adopción por usuario y ola. & Convocar con D-10 y confirmar responsables y agenda a D-5 del inicio de cada ola. Revisar transacciones diariamente durante marcha blanca y resolver ausencias el mismo día. & Escalar por CT a la gerencia del turno o canal; focalizar acompañamiento y registrar brechas. No acreditar adopción con registros sólo en el proceso anterior ni cerrar la ola sin evidencia. \\
R-014 / IMP & Gerente de Operaciones; Planificador de Rutas y su ayudante participan en transferencia y validación. & El antecedente ubica la jubilación dentro de dos años, sin fecha calendario especificada; confirmar con el Gerente de Operaciones y el planificador la fecha efectiva y disponibilidad en V-14, reservar talleres y entregar reglas y excepciones documentadas con prueba de transferencia al ayudante. La transferencia se comprueba con rutas y excepciones probadas sin el planificador durante la marcha blanca. & Solicitar disponibilidad en los primeros cinco días hábiles del mes 1 y confirmarla al cierre del mes 1. Ejecutar y validar transferencia en meses 1--3, antes de perder disponibilidad. Prueba sin el planificador en meses 13--15 (S7, 9.6.1). & Escalar por CT a Operaciones; activar talleres con ayudante y antecedentes verificables. Mantener reglas no corroboradas como incertidumbre y obtener validación antes de usarlas en rutas aceptadas. \\
R-015 / IMP & Gerente de Operaciones y Jefa de Bodega para turnos; Gerente Comercial para preventistas; Jefe de TI para administradores. & Confirmar nómina por perfil, liberación horaria, suplentes y acceso a sesiones; entregar disponibilidad para formación y evaluación y revisar evidencia individual de certificación. & Convocar con D-10 de cada sesión y confirmar nómina y horarios a D-5. Completar evaluación antes de habilitar al usuario y, en todo caso, antes de cerrar la marcha blanca de su ola. & Escalar por CT a la jefatura responsable el mismo día del incumplimiento; reprogramar sesiones y disponer suplente sólo si está certificado. Mantener pendiente la certificación y el cierre de ola afectado. \\
R-016 / JP & Contraparte Técnica, responsable de encauzar la decisión hacia Calidad, TI, Operaciones, Comercial o Finanzas según la materia. & Identificar en cada solicitud el área decisora, titular funcional y suplente, pregunta, antecedentes completos, tarea bloqueada y fecha límite; obtener respuesta o acta de decisión de la autoridad competente. Para promoción regional, identificar autorizador titular y suplente con facultad delegada. & Registrar agenda y responsables al cierre del mes 1. Para insumos de coordinación, cursar a D-10 y recibir respuesta a D-5 de la tarea D; para recepción de entregables rigen los diez días hábiles del art. 18.3. Excepción de S4, 4.3.2.5: avisar al suplente a los 15 min sin respuesta del titular y resolver dentro del máximo de 30 min. & JP notifica a CT el vencimiento el mismo día, solicita encauzamiento en un día hábil y, si sigue sin resolución en dos días hábiles, pide instancia ejecutiva extraordinaria. En promoción regional, escalar al suplente a los 15 min y mantener sin autorizar la promoción si vence el máximo de 30 min; registrar el efecto sobre el RTO, sin promover por silencio. Si amenaza seguridad, continuidad o hito, escala inmediatamente y mantiene bloqueada la tarea dependiente. \\
R-019 / ARQ & Contraparte Técnica; Jefe de TI coordina revisión técnica y la Gerenta General encauza autorización con la asesoría competente del CLIENTE. & Resolver por escrito regiones, residencia, transferencia y tratamiento de datos sobre el expediente de arquitectura y seguridad; identificar autoridad firmante y condiciones de habilitación. & Solicitar expediente y autoridad en mes 1, decisión inicial al cierre del mes 2 y siempre antes de habilitar entornos con datos sujetos a autorización; aplicar la fecha que venza primero. & Escalar por CT a la autoridad identificada y al Comité de Arquitectura; preparar alternativa con sus impactos. Mantener deshabilitada la transferencia no aprobada; no elegir unilateralmente otra región ni presumir autorización. \\
R-020 / SEG & Contraparte Técnica; Gerente de Operaciones autoriza la intervención operacional y Jefe de TI valida su ejecución técnica. & Aprobar flujo de autorización urgente, contactos y suplentes; por incidente, decidir por escrito la intervención, ventana, alcance, reversión y responsable con el expediente de seguridad. & Flujo aprobado al cierre del mes 1. Cada intervención se solicita al detectar la necesidad; si no hay respuesta, reiterar y escalar en una hora o antes si el plazo de atención o continuidad exige actuar. Autorización siempre previa al cambio. & Aplicar contención permitida por el procedimiento, notificar inmediatamente a CT y Operaciones y convocar decisión ejecutiva urgente. Mantener estado seguro sin eludir congelamiento, alterar plazos contractuales ni interpretar silencio como permiso. \\
R-024 / IMP & Gerente de Operaciones y representante de la empresa transportista; supervisor de andén verifica al conductor real. & Confirmar por despacho conductor, vehículo y viaje; acreditar identidad, autorización y certificación del reemplazo y registrar la vinculación auditada persona--vehículo--viaje. & Pedir nómina y suplentes a D-5 de incorporar la empresa. En cada despacho, confirmar conductor antes de salida; reemplazo no avisado se verifica y vincula en andén antes de habilitar uso. & Operaciones y el transportista buscan un conductor certificado para el mismo turno. Si no hay reemplazo, Operaciones redistribuye las entregas entre rutas habilitadas con capacidad o reprograma las restantes, avisa a los destinatarios y registra la decisión por entrega. El supervisor verifica al conductor antes de autorizar la salida. \\
R-025 / JP & Contraparte Técnica; gerencia solicitante fundamenta la petición y Comité Ejecutivo decide los cambios dentro de sus atribuciones. & Entregar solicitud escrita con objetivo y criterio de aceptación; decidir alcance, esfuerzo y programación sobre el análisis de impacto antes de incorporar el cambio a la línea base. & Registrar y remitir a CT el mismo día hábil de recibirla; comunicar fecha de decisión necesaria antes de su primera actividad dependiente. No iniciar la implementación hasta aprobación formal. & Escalar al Comité Ejecutivo si la decisión amenaza un hito. Mantener la línea base vigente y continuar el trabajo independiente; el silencio no aprueba el cambio ni permite cargar desarrollo adicional. \\
R-027 / JP & Gerente de Finanzas, con Gerente Comercial para canasta y malla; Contraparte Técnica canaliza conformidad y controversias. & Aprobar canasta, índices y regla estacional de 7.4.1, reproducir tres liquidaciones en sombra y conciliar evidencia de cada liquidación y de las diferencias del período. & Aprobar la fórmula inicial en el mes 17 (fecha nominal) y a más tardar al cierre del mes 20. Validar tres liquidaciones en sombra en meses 21--23 antes del primer efecto de pago del mes 24; revisión posterior antes de cerrar cada liquidación. & Escalar por CT a Finanzas y al Comité Ejecutivo. Sin regla o reproducción válida no activar el tramo; ante disputa, suspender sólo el devengo del período según el mecanismo contractual y mantener servicio y acompañamiento fijo. \\
R-028 / JP & Gerente de Operaciones y Jefe de TI; sindicato de conductores como interlocutor del acuerdo, por CT. & Acta sobre finalidades, acceso y controles de la telemetría y uso de terminales; documentar observaciones y el acuerdo exigible para cualquier control de jornada. & Convocar a D-10 y resolver a D-5 del inicio de la ola de reparto; acuerdo 9.7.2 antes de habilitar el uso que lo requiere. & JP escala por CT a Operaciones; no habilitar cámaras en cabina ni control de jornada sin acuerdo. Mantener el uso operacional autorizado y la cobertura pendiente identificada. \\
R-029 / IMP & Jefa de Bodega y Gerente de Operaciones; comunicación de altas y disponibilidad de tutores por turno. & Nómina de altas, turnos y tutores; registro de tiempo de aprendizaje y evaluación individual del preparador, con cobertura certificada mientras completa formación. & Meses 13--56; informar cada alta antes de asignarla al turno y completar inducción/evaluación antes del trabajo autónomo. Revisar cobertura en cada relevo. & IMP, y SRE desde mes 22, solicita cobertura certificada y refuerzo por CT a Bodega/Operaciones; mantener sin habilitación autónoma al preparador no evaluado. \\
R-030 / ARQ & Jefe de TI; custodios del ERP y WMS y proveedor del legado cuando su acceso sea necesario. & Accesos autorizados, formatos, diccionario y contratos de interfaz/extracción; resultado de ensayo temprano y restricciones del proveedor del ERP/WMS. & Solicitar en los primeros cinco días hábiles del mes 1 y verificar contrato inicial al cierre del mes 1; probar cada interfaz antes de su integración y cada extracción antes de validar la carga. & ARQ escala por CT a TI, prepara alternativa de archivo controlado y mide sus límites; mantener bloqueada la aceptación dependiente sin modificar el ERP ni omitir datos. \\
R-031 / JP & Gerente de Operaciones y Contraparte Técnica; gerencias usuarias para disponibilidad de validación. & Agenda de usuarios, decisiones y continuidad autorizada para el peak; estado diario de las condiciones de cierre y resolución formal de impactos en hitos. & Preparar antes del mes 19 y completar habilitaciones en agosto bajo el supuesto de inicio de S7. Verificación diaria en meses 19--20 y antes del cierre previsto del mes 20. & JP escala inmediatamente por CT al Comité Ejecutivo; mantener continuidad autorizada y no cerrar con condiciones incumplidas ni desplegar en congelamiento sin autorización. \\
R-032 / ARQ & Jefe de TI y Contraparte Técnica; Operaciones para ventanas de actualización y aceptación del servicio. & Calendario de ventanas, política de versiones y conformidad con las pruebas de adecuación; aprobación de cambios que afecten compromisos. & Inventario en mes 1, revisión trimestral hasta mes 56; acordar intervención antes del límite de soporte comprobado y antes de activar la versión. & ARQ escala por CT el impacto; ensayar reversión y reprogramación autorizada de la adecuación, preservando atención de vulnerabilidades y compromisos de servicio. \\
R-033 / ARQ & Jefe de TI y Contraparte Técnica; autoridad contractual del CLIENTE para condiciones de salida y acceso. & Accesos y titularidad de cuentas/claves, evidencia de exportación y restauración y decisión sobre restricciones contractuales de sustitución. & Verificar en recepción inicial y actualización anual de 1.6.3; resolver cada restricción antes del ensayo o transición de salida que dependa de ella. & ARQ eleva por CT la restricción al Comité Ejecutivo y prepara salida gradual comprobable; conservar acceso y continuidad bajo condiciones autorizadas sin presumir equivalencia entre proveedores. \\
R-034 / ARQ & Jefe de TI y Gerente de Finanzas; Operaciones coordina cargas y salidas afectadas. & Procedimiento de restitución del ERP, responsables y evidencias de emisión/anulación, con resultado de ensayo de pérdida de sala y cambio de carga. & Acordar procedimiento antes del ensayo AL-DTE-01/AL-DR-01 de H5/H10. En incidente, comunicar de inmediato a TI/Finanzas y decidir cada salida antes de despachar. & ARQ/SRE escalan por CT a TI/Finanzas; bloquear salida sin guía válida y conservar la carga original sólo si conserva guía vigente. Medir restitución completa frente al RTO exigido. \\
R-035 / IMP & Gerente de Operaciones y representante de cada empresa transportista; conductor como participante de la validación práctica. & Disponibilidad de conductores y demostración de uso en andén, evidencia de inducción y compromiso operativo del transportista ante rechazo individual. & Confirmar participantes a D-5 de la ola de reparto y demostrar uso antes de la primera salida del conductor; actuar en andén ante cada rechazo. & IMP escala por CT a Operaciones y al transportista; solicitar conductor certificado y dispuesto y acompañar. Registrar entregas/evidencias pendientes sin dar cobertura por cumplida. \\
R-036 / ARQ & Gerente de Operaciones y Jefe de TI; SRE gestiona desde mes 22, familia ARQ. & Entregar y validar perfil por ventana, concentración temporal y mezcla de transacciones, incluyendo peak y crecimiento comprometidos; conservar mediciones y aprobación de cada versión. & Solicitar en los primeros cinco días hábiles del mes 1; validar al cierre del mes 1 y revalidar antes de H5/H10, preparación del peak y cambios del perfil. & ARQ, y SRE desde mes 22, registra divergencia y avisa a JP/CT el mismo día hábil; presenta recuperación en dos días hábiles y siempre antes del hito. Conservar bloqueo, corregir y repetir CT-14 en ventana autorizada antes de habilitar. Si la divergencia pertenece al alcance comprometido, LafroX corrige como trabajo ordinario; no presume ampliación contractual. \\
\end{tablalafrox}
\end{landscape}
\endgroup

\clearpage
\section{Anexo 8.B: Contingencia de esfuerzo por riesgo}

Este anexo estima la exposición de esfuerzo de los 36 riesgos de T-16. Para cada riesgo se presentan la probabilidad de ocurrencia $p_0$, los episodios condicionales previstos $n$ y los impactos acumulados optimista (O), más probable (M) y pesimista (P), en HH. Cada escenario suma las estimaciones de los dos perfiles para un episodio y las multiplica por $n$. La estimación PERT acumulada es $\mathrm{PERT}=(O+4M+P)/6$ y el valor esperado es $\mathrm{VE}=p_0\times\mathrm{PERT}$. Son estimaciones de planificación, no frecuencias históricas ni compromisos de esfuerzo asignado; deben validarse con los responsables. La suma de VE se presenta como exposición inicial de esfuerzo conciliada con la proyección y reservas propuestas de 8.H, no como resultado de una simulación ni como distribución temporal o por perfil. Las respuestas no modifican la estimación inicial $p_0$ de esta tabla; la proyección posterior se calcula por separado en 8.F.

\begingroup
\renewcommand{\thetable}{8.B.1}
\renewcommand{\arraystretch}{0.9}
\begin{landscape}
\setlength{\textwidth}{\linewidth}
\begin{tablalafrox}{Estimación acumulada de esfuerzo y valor esperado por riesgo (HH). Fuente: Riesgos de formulario T-16}{tab:sd8-anexo-b-pert}{G{0.07} G{0.08} G{0.08} G{0.10} G{0.10} G{0.10} G{0.12} G{0.12}}{\cab{ID} & \cab{$p_0$} & \cab{Episodios $n$} & \cab{O (HH)} & \cab{M (HH)} & \cab{P (HH)} & \cab{PERT (HH)} & \cab{VE (HH)}}
R-001 & 45\% & 1 & 40,00 & 96,00 & 200,00 & 104,00 & 46,80\\
R-002 & 60\% & 1 & 72,00 & 184,00 & 360,00 & 194,67 & 116,80\\
R-003 & 50\% & 1 & 48,00 & 120,00 & 240,00 & 128,00 & 64,00\\
R-004 & 55\% & 1 & 40,00 & 104,00 & 208,00 & 110,67 & 60,87\\
R-005 & 50\% & 1 & 40,00 & 88,00 & 184,00 & 96,00 & 48,00\\
R-006 & 55\% & 1 & 40,00 & 96,00 & 184,00 & 101,33 & 55,73\\
R-007 & 35\% & 1 & 24,00 & 72,00 & 176,00 & 81,33 & 28,47\\
R-008 & 10\% & 6 & 12,00 & 30,00 & 72,00 & 34,00 & 3,40\\
R-009 & 35\% & 1 & 32,00 & 72,00 & 144,00 & 77,33 & 27,07\\
R-010 & 30\% & 1 & 24,00 & 64,00 & 128,00 & 68,00 & 20,40\\
R-011 & 35\% & 1 & 20,00 & 40,00 & 80,00 & 43,33 & 15,17\\
R-012 & 25\% & 1 & 20,00 & 56,00 & 112,00 & 59,33 & 14,83\\
R-013 & 55\% & 1 & 48,00 & 120,00 & 240,00 & 128,00 & 70,40\\
R-014 & 20\% & 1 & 40,00 & 96,00 & 192,00 & 102,67 & 20,53\\
R-015 & 30\% & 1 & 32,00 & 80,00 & 160,00 & 85,33 & 25,60\\
R-016 & 45\% & 1 & 32,00 & 80,00 & 160,00 & 85,33 & 38,40\\
R-017 & 25\% & 8 & 24,00 & 48,00 & 144,00 & 60,00 & 15,00\\
R-018 & 5\% & 2 & 12,00 & 24,00 & 64,00 & 28,67 & 1,43\\
R-019 & 30\% & 1 & 24,00 & 64,00 & 136,00 & 69,33 & 20,80\\
R-020 & 15\% & 2 & 16,00 & 40,00 & 80,00 & 42,67 & 6,40\\
R-021 & 15\% & 1 & 8,00 & 20,00 & 40,00 & 21,33 & 3,20\\
R-022 & 10\% & 1 & 8,00 & 20,00 & 40,00 & 21,33 & 2,13\\
R-023 & 10\% & 1 & 8,00 & 24,00 & 48,00 & 25,33 & 2,53\\
R-024 & 45\% & 3 & 18,00 & 36,00 & 72,00 & 39,00 & 17,55\\
R-025 & 35\% & 2 & 32,00 & 80,00 & 160,00 & 85,33 & 29,87\\
R-026 & 20\% & 1 & 32,00 & 80,00 & 160,00 & 85,33 & 17,07\\
R-027 & 40\% & 1 & 24,00 & 56,00 & 120,00 & 61,33 & 24,53\\
R-028 & 45\% & 1 & 24,00 & 64,00 & 128,00 & 68,00 & 30,60\\
R-029 & 55\% & 6 & 96,00 & 192,00 & 384,00 & 208,00 & 114,40\\
R-030 & 50\% & 1 & 48,00 & 128,00 & 280,00 & 140,00 & 70,00\\
R-031 & 45\% & 1 & 24,00 & 64,00 & 160,00 & 73,33 & 33,00\\
R-032 & 40\% & 2 & 32,00 & 96,00 & 224,00 & 106,67 & 42,67\\
R-033 & 25\% & 1 & 40,00 & 120,00 & 280,00 & 133,33 & 33,33\\
R-034 & 30\% & 2 & 32,00 & 96,00 & 240,00 & 109,33 & 32,80\\
R-035 & 40\% & 2 & 24,00 & 64,00 & 160,00 & 73,33 & 29,33\\
R-036 & 30\% & 1 & 16,00 & 48,00 & 120,00 & 54,67 & 16,40 \\
 & & & & & & \textbf{Total VE} & 1199,52\\
\end{tablalafrox}
\end{landscape}

\endgroup

\noindent La cuantificación depende de probabilidades, rangos de esfuerzo y episodios condicionales estimados; no demuestra frecuencias observadas ni incorpora dependencias entre riesgos. Los resultados no constituyen una simulación, una asignación por período o perfil, ni una garantía de esfuerzo suficiente. El total se calcula sumando los valores esperados sin redondear y redondeando sólo el resultado final; las cifras VE por riesgo se muestran a dos decimales, por lo que su suma impresa puede diferir en 0,01 HH.



\subsection*{Hipótesis de ocurrencia y esfuerzo}
\addcontentsline{toc}{subsection}{Hipótesis de ocurrencia y esfuerzo}

Las probabilidades y HH de los riesgos R-028--R-036 son juicios explícitos de planificación basados en las condiciones citadas, no frecuencias observadas ni compromisos aceptados. $p_0$ significa al menos un episodio del evento durante la ventana indicada; $n$ estima episodios condicionados a que el evento ocurra. Los escenarios O/M/P de 8.B.1 son acumulados: se suman los dos perfiles por episodio y se multiplican por $n$ una sola vez, antes de aplicar PERT y $p_0$. El 38\% de rotación anual y las 2.600 entregas del peak no se usan como probabilidades.

Las HH representan trabajo extraordinario de LafroX por encima del alcance ordinario; no valorizan pérdidas de mercadería, reconstrucción del ERP del CLIENTE, tareas ordinarias de capacitación ni aumentos de demanda ya incluidos en el dimensionamiento. R-031 recoge sólo coordinación y recuperación de programación; la verificación y remediación de capacidad corresponde al control técnico de escalabilidad CT-14. Un incidente recibe un ID causal principal y sus HH se imputan una sola vez; los efectos en otros riesgos quedan vinculados sin repetir el esfuerzo. R-017 no duplica el mismo corte WAN como caída de ERP, SII, POS o stock central; R-034 cuantifica la dependencia del ERP del CLIENTE; CT-15 establece la corrección y verificación de la consolidación a cargo de LafroX. R-032 y R-033 cubren cambios excepcionales externos de soporte o restricciones efectivas de salida. El mantenimiento, la portabilidad y las actualizaciones ordinarias son obligaciones de LafroX; no se presupuestan otra vez como contingencia. Los controles técnicos se realizan dentro del trabajo ordinario y no se presupuestan nuevamente como contingencia.

R-017 adopta $p_0=25\%$ por la exposición inicial de Concepción con un enlace sin respaldo verificado y la falla del último camino utilizable hasta mes 56. Es una hipótesis, no una frecuencia inferida del caso ni de los cortes de Talca. Se estiman ocho episodios condicionados, escenarios acumulados 24/48/144 HH, PERT de 60 HH y VE de 15 HH. R-014 adopta $p_0=20\%$ y un impacto de nivel 4 por la consecuencia R18-16 y la jubilación prevista dentro de dos años, sin fecha calendario especificada; el evento incierto es la transferencia insuficiente de conocimiento, no el retiro conocido.

\begingroup
\renewcommand{\thetable}{8.B.2}
\begin{landscape}
\setlength{\textwidth}{\linewidth}
\begin{tablalafrox}{Hipótesis de ocurrencia, ventana y esfuerzo extraordinario de R-028--R-036. Fuente: condiciones citadas y estimación de planificación de LafroX.}{tab:sd8-ampliacion-hipotesis}{G{0.07} F{0.19} Y F{0.28} F{0.18}}{\cab{ID / $p_0$ / $n$} & \cab{Ventana de exposición} & \cab{Fundamento del juicio de probabilidad} & \cab{Base O/M/P y exclusiones} & \cab{Fuente de la condición}}
R-028 / 45\% / 1 & Desde levantamiento de finalidades hasta incorporación de cada ola de reparto, meses 1--20. & La objeción ya existe, pero el evento incierto es el desacuerdo sobre el uso propuesto. Se adopta 45\% para una ronda extraordinaria de coordinación previa a la ola; no se supone rechazo seguro. & JP y IMP: cada perfil 12/32/64 HH por episodio; un episodio. Sólo negociación y reprogramación extraordinarias, fuera del levantamiento y acuerdo ordinarios. & Caso, caps. 10 (restricción 10), 16.1 y 19; S3, S-13; S7, acuerdo 9.7.2. \\
R-029 / 55\% / 6 & Meses 13--56; desde primera marcha blanca y durante la operación nocturna. & El 38\% es rotación anual, no p0. Se propone 55\% de al menos un desborde de inducción durante 44 meses y, condicionado a ello, seis cohortes con atención extraordinaria; son hipótesis de exposición recurrente. & IMP y CAL: cada perfil 8/16/32 HH por cohorte; seis cohortes. Excluir las 2.816 HH ordinarias de 9.5.2 y el trabajo autónomo del preparador. & Caso, caps. 7.2, 10 (restricción 11) y 19; S7, pp. 21--22 y 32; T-14, 9.5.2; T-15, meses 13--56. \\
R-030 / 50\% / 1 & Levantamiento del mes 1 hasta integración y migraciones de ambas etapas, antes de H5/H10. & Se adopta 50\% ante ausencia documental confirmada y capacidad de acceso sin confirmar, para un bloqueo conjunto de descubrimiento; no se confunde documentación ausente con interfaz inexistente. & ARQ y DAT: cada perfil 24/64/140 HH; un episodio de descubrimiento y reproceso extraordinarios. Excluir levantamiento y adaptadores ordinarios. & Caso, cap. 5 y cap. 19; S3, S-18, S-14 y V-16; S7, 2.1.5 y 5.1.1. \\
R-031 / 45\% / 1 & Cierre de E2 en meses 19--20; septiembre del calendario supuesto de inicio de S7, sujeto a confirmar inicio contractual. & El peak y el congelamiento son condiciones conocidas. El 45\% representa un bloqueo del cierre en esa coincidencia, no la probabilidad de que septiembre ocurra. & JP y IMP: cada perfil 12/32/80 HH; un episodio. Sólo coordinación, validación y recuperación de calendario; la corrección técnica de una saturación pertenece al control técnico de escalabilidad CT-14. & Caso, caps. 10 (restricción 8), 17.5 y 19; S7, secciones 7.3.1 y 7.3.5; T-18. \\
R-032 / 40\% / 2 & Meses 1--56; ciclos de soporte y compatibilidad efectivamente publicados, sin suponer fechas de fin de soporte. & Se propone 40\% de al menos una adecuación excepcional en 56 meses y dos episodios condicionados. La vigencia del soporte y su fecha de término se verifican con la información publicada por el proveedor. & ARQ y DES: cada perfil 8/24/56 HH por episodio; dos episodios. Excluir las actualizaciones anuales ordinarias 10.2.5 y mantenimiento ya incluido. & S4, plataforma declarada; S7, 10.2.5. \\
R-033 / 25\% / 1 & Meses 1--56 y preparación de la salida contractual; no se asume que la salida extraordinaria ocurrirá. & Se propone 25\% para una restricción extraordinaria de salida o sustitución en el contrato; el uso de AWS por sí solo no demuestra bloqueo. & ARQ y SRE: cada perfil 20/60/140 HH; un episodio. Sólo resolución excepcional de restricciones; excluir actualización anual del plan y salida ordinaria contratada. & S4, servicios gestionados; T-14, 1.6.3 y cuenta 10.5. \\
R-034 / 30\% / 2 & Desde pruebas de emisión y primer despacho hasta mes 56; dos episodios excepcionales condicionados. & Se propone 30\% de al menos un bloqueo de nueva emisión por ERP indisponible en el horizonte; no se deriva de la respuesta del SII ni se duplica pérdida de sala como otro evento. & ARQ y SRE: cada perfil 8/24/60 HH por episodio; dos episodios. Sólo coordinación, recuperación y validación extraordinarias de la integración; no incluye reconstruir el ERP del CLIENTE. & Caso, cap. 10, restricciones 1 y 4; S4, pp. 77 y 126. \\
R-035 / 40\% / 2 & Incorporación de empresas y conductores, y reemplazos durante operación hasta mes 56. & Se propone 40\% de al menos un rechazo individual en incorporación y operación; dos episodios condicionados. El hecho de trabajar para un tercero no demuestra rechazo. & IMP y CAL: cada perfil 6/16/40 HH por episodio; dos episodios. Sólo atención extraordinaria del rechazo; excluir inducción ordinaria de 9.5.3, firma de R-011 y habilitación de R-024. & Caso, caps. 10 (restricción 6), 15, RT-22.04 y 19; S7, 9.7.1 y 9.5.3. \\
R-036 / 30\% / 1 & Meses 1--56; un episodio acumulado en toda la ventana. & Hipótesis no observada de divergencia de concentración temporal o mezcla respecto del perfil validado; el peak conocido no es incierto. p0=p1=30\%, sin descuento por eficacia futura. & ARQ/CAL: 60/40 de O/M/P=16/48/120 HH acumuladas; PERT=54,67 HH. Sólo coordinación y recuperación extraordinarias. CT-14 y correcciones del alcance son ordinarias; A=0 y C=16 HH contenidas en PERT. & Aclaración 2, 8.2; Caso, caps. 17.5 y 19; perfil que se confirma con Operaciones/TI; CT-14. \\
\end{tablalafrox}
\end{landscape}
\endgroup

\clearpage
\section{Anexo 8.C: Respaldo de la valoración inicial}

Se aplican las escalas de 8.1.3 a los 36 eventos del T-16. $I_{HH}$ es el nivel de esfuerzo del PERT condicional acumulado del Anexo 8.B; $I_D$ es el mayor nivel de las dimensiones no HH identificadas en la justificación; $I=\max(I_{HH},I_D)$. La justificación señala la dimensión determinante y distingue una amenaza a aceptación (nivel 4) de una consecuencia grave (nivel 5). Son juicios iniciales de planificación sobre el evento declarado, sin descuento por eficacia futura ni cálculo de riesgo residual. La probabilidad se deriva de $p_0$ del Anexo 8.B; el fundamento de la estimación se explicita en 8.B.2. Las referencias contractuales corresponden a las disposiciones ya descritas en 8.3.3; no se estiman multas ni importes.

\begingroup
\renewcommand{\thetable}{8.C.1}
\renewcommand{\arraystretch}{1.10}
\begin{tablalafrox}{Fundamento del impacto inicial de los 36 riesgos. Fuente: eventos del T-16, PERT del Anexo 8.B y escalas de 8.1.3.}{tab:sd8-impactos-respaldo}{G{0.08} G{0.17} G{0.12} G{0.12} Y}{\cab{ID} & \cab{PERT HH / $I_{HH}$} & \cab{$I_D$ / I} & \cab{$q_I$} & \cab{Dimensión y fundamento}}
R-001 & 104,00 / 3 & 4 / 4 & 0,40 & Aceptación: La certificación EDI incompleta pone en riesgo la evidencia de las cuatro semanas finales de E2. La carga manual controlada mantiene la continuidad y la activación EDI requiere certificación. \\
R-002 & 194,67 / 4 & 5 / 5 & 0,80 & Calidad sanitaria: Los vínculos sanitarios sin evidencia de origen pueden ser irrecuperables; el saneamiento no reconstruye una historia que no existe. \\
R-003 & 128,00 / 3 & 3 / 3 & 0,20 & Alcance: El historial incompleto exige reproceso de la carga y excepciones de reposición; el evento declarado no demuestra pérdida definitiva del alcance. \\
R-004 & 110,67 / 3 & 4 / 4 & 0,40 & Aceptación: Las diferencias sin explicación ni ajuste aprobado bloquean la aceptación de inventario antes del corte, condición de marcha blanca del art. 17.3. \\
R-005 & 96,00 / 3 & 4 / 4 & 0,40 & Aceptación: Sin saldo inicial conciliado no puede demostrarse M8 con carga aceptada; requiere resolver la evidencia de aceptación de los arts. 18.1--18.2. \\
R-006 & 101,33 / 3 & 4 / 4 & 0,40 & Aceptación: La ausencia de antecedentes verificables impide respaldar las comparaciones de OTIF o cuadratura exigidas para aceptar el período de referencia. \\
R-007 & 81,33 / 3 & 3 / 3 & 0,20 & Alcance: La deuda viva adicional obliga a extraer y conciliar nuevamente el entregable; no se declara pérdida de datos ni incumplimiento definitivo. \\
R-008 & 34,00 / 2 & 3 / 3 & 0,20 & Plazo: El cierre tributario queda pendiente de confirmación externa y requiere seguimiento y conciliación del entregable; el silencio no acredita aprobación. \\
R-009 & 77,33 / 2 & 4 / 4 & 0,40 & Aceptación sanitaria: La falta de aprobación bloquea la certificación térmica y el cierre de marcha blanca; el evento es aprobación pendiente, sin afirmar daño sanitario consumado. \\
R-010 & 68,00 / 2 & 4 / 4 & 0,40 & Plazo: La falta de equipos bloquea instalación y pruebas de una ola y amenaza su aceptación y secuencia de implantación. \\
R-011 & 43,33 / 2 & 4 / 4 & 0,40 & Alcance: Sin acuerdo no se acredita la cobertura del transportista ni sus pruebas de la ola, por lo que se amenaza la aceptación del alcance comprometido. \\
R-012 & 59,33 / 2 & 4 / 4 & 0,40 & Aceptación sanitaria: La negativa impide integrar el termógrafo y acreditar temperatura de la ruta; amenaza la aceptación sin demostrar por sí sola daño sanitario grave. \\
R-013 & 128,00 / 3 & 4 / 4 & 0,40 & Aceptación: El registro exclusivo fuera de la solución impide acreditar adopción y puede extender el cierre de la ola y su marcha blanca. \\
R-014 & 102,67 / 3 & 4 / 4 & 0,40 & Aceptación: el conocimiento debe quedar transferido y probado antes de la jubilación para acreditar R18-16; el antecedente de caída de 6 a 9 puntos de cumplimiento amenaza continuidad y aceptación. \\
R-015 & 85,33 / 3 & 4 / 4 & 0,40 & Aceptación: Usuarios sin evaluación completa no acreditan capacitación y certificación, condición acumulativa de cierre de marcha blanca del art. 17.3. \\
R-016 & 85,33 / 3 & 4 / 4 & 0,40 & Plazo: La decisión pendiente bloquea tareas de diseño, prueba o aceptación y consume holgura, amenazando hitos que requieren decisión coordinada. \\
R-017 & 60,00 / 2 & 4 / 4 & 0,40 & Servicio: La interrupción limita servicios centrales y acumula sincronización; amenaza continuidad y niveles de servicio, sin suponer eficacia del respaldo. \\
R-018 & 28,67 / 2 & 4 / 4 & 0,40 & Servicio: Superar la autonomía deja componentes sin servicio hasta restituir energía; amenaza continuidad sin imputar pérdidas ajenas al proyecto. \\
R-019 & 69,33 / 2 & 4 / 4 & 0,40 & Aceptación: Sin aprobación de residencia no puede habilitarse ni acreditarse la recuperación propuesta; requiere decisión antes de aceptar los entornos. \\
R-020 & 42,67 / 2 & 4 / 4 & 0,40 & Contrato: La autorización pendiente enfrenta congelamiento y plazo de remediación, amenazando cumplimiento; no demuestra por sí sola pérdida de continuidad contractual. \\
R-021 & 21,33 / 1 & 2 / 2 & 0,10 & Plazo: El efecto declarado se limita a repetir casos de un recorrido dentro del paquete de validación; no se establece afectación de un hito contractual. \\
R-022 & 21,33 / 1 & 2 / 2 & 0,10 & Plazo: Se evalúa la repetición de casos de prueba interrumpidos; las consecuencias personales del siniestro se atienden por el procedimiento del transportista y no son HH del proyecto. \\
R-023 & 25,33 / 2 & 2 / 2 & 0,10 & Plazo: El efecto del proyecto es reprogramar y repetir el recorrido de validación; la valoración no incluye pérdidas de vehículo o carga. \\
R-024 & 39,00 / 2 & 3 / 3 & 0,20 & Alcance: El reemplazo requiere habilitación, formación extraordinaria y verificación de registros del entregable de implantación. \\
R-025 & 85,33 / 3 & 3 / 3 & 0,20 & Contrato: El efecto estimado es evaluar y tramitar la solicitud; el desarrollo adicional sólo existe tras aprobación y ajuste de línea base. \\
R-026 & 85,33 / 3 & 4 / 4 & 0,40 & Contrato: Una exigencia aplicable obliga a adecuar y validar controles de la solución; amenaza cumplimiento y requiere decisión contractual, sin presumir una reforma o multa concreta. \\
R-027 & 61,33 / 2 & 3 / 3 & 0,20 & Contrato: La discrepancia afecta el cálculo y la liquidación de un período del tramo variable; el evento no establece interrupción ni terminación contractual. \\
R-028 & 68,00 / 2 & 4 / 4 & 0,40 & Aceptación: el desacuerdo puede impedir incorporar a los usuarios y completar pruebas de la ola; no se presupone habilitación de vigilancia excluida. \\
R-029 & 208,00 / 4 & 4 / 4 & 0,40 & Aceptación y calidad: personal no preparado amenaza la certificación y el flujo nocturno; el esfuerzo condicional de 208 HH también alcanza nivel 4. \\
R-030 & 140,00 / 3 & 4 / 4 & 0,40 & Aceptación y plazo: sin intercambio o extracción verificados se bloquean integración y migración exigidas, aunque los datos históricos existan. \\
R-031 & 73,33 / 2 & 4 / 4 & 0,40 & Contrato y aceptación: una condición incumplida impide cerrar marcha blanca y amenaza su hito; no se compensa con otras condiciones cumplidas. \\
R-032 & 106,67 / 3 & 4 / 4 & 0,40 & Seguridad y servicio: una incompatibilidad que excede mantenimiento ordinario puede impedir conservar la versión segura y el servicio aceptado. \\
R-033 & 133,33 / 3 & 4 / 4 & 0,40 & Contrato y continuidad: una restricción no resuelta impide una transición comprobable del servicio o la entrega utilizable de datos. \\
R-034 & 109,33 / 3 & 4 / 4 & 0,40 & Servicio y contrato: impedir una guía nueva bloquea despachos y amenaza la ventana crítica, aun cuando el WMS esté recuperado. \\
R-035 & 73,33 / 2 & 4 / 4 & 0,40 & Aceptación y adopción: un conductor que no usa el flujo deja evidencia y cobertura pendientes, aun con firma de empresa y certificación previa. \\
R-036 & 54,67 / 2 & 4 / 4 & 0,40 & Aceptación y servicio: la divergencia exige revalidar capacidad antes de habilitar el alcance; puede amenazar aceptación o servicio. Se conserva I=4 aunque el esfuerzo sea nivel 2; no se presupone daño sanitario ni cambio de alcance. \\
\end{tablalafrox}
\endgroup


\clearpage
\section{Anexo 8.D: Dependencias operativas y continuidad}
Las condiciones siguientes delimitan la operación, las decisiones y las pruebas habilitantes. Los responsables deben verificar la ejecución y eficacia de los controles antes de autorizar el alcance afectado. Los meses se cuentan desde el inicio contractual. Las referencias R identifican riesgos; CT-01--CT-15 identifican los 15 controles técnicos vigentes; CT-15 conserva stock, antes CT-16, e identidad conectada/local se integra en CT-03; no se renumeran los 15 controles vigentes. CT sin número designa la Contraparte Técnica.

\begingroup
\renewcommand{\thetable}{8.D.1}
\begin{landscape}
\setlength{\textwidth}{\linewidth}
\renewcommand{\arraystretch}{1.08}
\begin{tablalafrox}{Dependencias operativas, decisiones y condiciones de continuidad. Fuente: requisitos de operación del Caso 02; arquitectura, procedimientos de continuidad y programa de implantación.}{tab:sd8-dependencias-operativas}{F{0.14} Y F{0.19} F{0.31}}{\cab{Ámbito} & \cab{Condición operativa} & \cab{Responsable y contraparte} & \cab{Momento, evidencia y decisión}}
Piloto INN-01 / CT-08 & Validar catálogo con Comercial/Calidad en meses 7--9. El paquete abarca 11--15; H5 corresponde a 11--12 y el piloto real a 13--15. & DAT dirige; IMP coordina; CAL verifica. & I-01A en mes 15: al menos 70\% de aciertos sobre al menos 100 avisos evaluables. Sin cumplimiento, no ampliar a otras rutas; mantener avisos informativos y alcance obligatorio. \\
Identidad / CT-03 & Autenticación conectada y local con manifiesto, permisos, turno y dispositivo válidos; expiración y baja sin ampliar vigencia. & SEG responde; ARQ diseña; DES/SRE ejecutan; CAL verifica con TI/Operaciones. & AL-ID-01 y AL-OFF-01 antes de H5/H10: 26 h, corte una hora después, relevos a las 8/16 h y ruta de 14,5 h acotada a 16 h. Denegaciones, cola y conciliación conforme a 8.E y matriz 8.D.3; seguridad conserva G5. \\
Emisión ERP / R-034 & El ERP del CLIENTE es el único emisor y no se modifica. La pérdida de Talca bloquea nuevas guías hasta recuperar el ERP; sólo se usan guías válidas ya emitidas. R-008 trata la confirmación externa del SII. & ARQ coordina; SRE ejecuta con TI y Operaciones. & AL-DTE-01 y AL-DR-01 antes de H5/H10: demostrar frontera de emisión y recuperación; RTO de 4 h y RPO de 15 min. No simular disponibilidad del ERP mediante otro emisor. \\
Stock / CT-15 & Reserva central confirmada sólo con retención local durable; autoridad única por ámbito, exclusión del escritor anterior y conciliación de UUID, sitio y época. & Equipo DAT hasta mes 20; SRE desde mes 21, con apoyo puntual de DAT; ARQ diseña; CAL verifica. & AL-STOCK-01: última unidad, pérdida de acuse y reintento sin doble reserva; preventa E1 y portal/EDI reales E2. AL-DR-01: pérdida, promoción, retorno y conciliación. R-004 trata inventario inicial y R-017 conectividad; no duplicar contingencia. \\
Pagos y servicios externos / CT-01--02 & Consultar estado antes de reintentar pagos y conciliar con clave idempotente. Ante caída de mapas, usar sólo rutas autorizadas disponibles y mantener nuevas consultas pendientes. & ARQ diseña; DES/SRE ejecutan; CAL verifica. & Pruebas de timeout antes y después del cargo, reenvío, caída y recuperación antes de H5/H10. Imputar un mismo corte WAN una sola vez a su causa principal. \\
Conocimiento de rutas / R-014 & Transferir conocimiento en meses 1--3 y verificar fecha V-14 de retiro. La prueba debe prescindir de la intervención del planificador. & IMP con Operaciones y planificador; CAL verifica. & Prueba de planificación en meses 13--15 sin el planificador. Conservar evidencia del resultado R18-16 y escalar conocimiento no transferido antes de habilitar autonomía. \\
Enlaces / R-017 & Verificar respaldo recibido en Concepción antes de aceptar la ola. Talca/Concepción: fibra y LTE, Starlink en espera caliente. Cross-docking: Starlink y LTE Entel/Claro; conmutación menor a 30 s. & SRE con TI/Operaciones; CAL verifica recepción. & Recepción antes de H3 y aceptación ampliada antes de H5. Ensayar pérdida del último camino Starlink de madrugada en Los Ángeles. El diseño conectado no garantiza disponibilidad del enlace. Se conserva el tratamiento del corte conforme a R-017. \\
Energía / R-018 & Distinguir autonomía de UPS de 30 min de continuidad de Talca de 24 h; la autonomía de un componente no demuestra la de todo el servicio. & SRE con TI/Operaciones; CAL verifica. & Ensayar pérdida de suministro, transferencia, continuidad y restitución antes de aceptación. Registrar medición por componente y servicio; no acreditar autonomía sin evidencia. \\
Autorizaciones / R-016, R-019 & La aprobación previa de residencia de datos es distinta de la decisión de promover durante un incidente. El silencio no aprueba una transferencia ni una promoción. & JP y ARQ; titular y suplente autorizados del CLIENTE. & Durante promoción regional, si el titular no responde en 15 min, decide el suplente dentro del máximo de 30 min. Validar residencia antes de habilitar transferencia; conservar actas y evidencia. \\
Cobertura RF en cámaras de frío & Medir señal y comprobar la comunicación y el registro de operaciones desde los puntos de trabajo, con terminales, puertas y ocupación representativos. & ARQ/SRE con TI y Jefa de Bodega; CAL verifica. & Antes de aceptar la instalación, registrar por cámara los puntos medidos, las condiciones del ensayo, los resultados y las zonas sin cobertura. Corregir y repetir las pruebas afectadas hasta verificar los puntos de operación. R-017 cubre los enlaces WAN. \\
Redundancia física & Verificar pares de firewall, nodos, discos y switches, junto con autoridad de escritura, recuperación y capacidad efectiva. & ARQ diseña; SRE ejecuta; CAL verifica. & Antes de aceptar infraestructura, inyectar fallas y registrar conmutación y restitución. La redundancia declarada no reemplaza la prueba del servicio ni permite doble escritor. \\
Rotación nocturna / R-029 & El 38\% es rotación anual, no probabilidad del evento. Mantener inducción continua en meses 13--56, turno 22:00--06:00, aprendizaje máximo de dos horas y evaluación previa al trabajo autónomo. & IMP hasta mes 21; SRE desde 22; Jefa de Bodega y Operaciones aportan tutor y cobertura. & Verificar competencia de cada reemplazo y capacidad de tutoría; impedir trabajo autónomo no certificado. No imputar otra vez las HH ordinarias de capacitación ni confundir rotación con ausencia a sesiones o resistencia individual. \\
Recuperación regional & Distinguir objetivo de recuperación de 75 min (15+60) del máximo sanitario de 2 h; no consumir ese margen como disponibilidad autorizada. & SRE ejecuta; ARQ valida frontera; Calidad del CLIENTE decide aspectos sanitarios. & Pruebas AL-DTE-01 y AL-DR-01, con 96 salidas válidas, antes de H5/H10 y semestralmente. Medir tiempos, integridad y condiciones de reanudación por procedimiento. \\
Cierre bajo peak / R-031, CT-14 & Acordar continuidad, usuarios y congelamiento antes del mes 19; mantener seguimiento diario del cierre E2 en meses 19--20. Separar coordinación del cierre de verificación técnica de capacidad. & JP con Contraparte Técnica/Operaciones; SRE y ARQ verifican capacidad. & Comprobar carga por ventana y dependencia, incluido ERP, antes de H5/H10 y preparación del peak. Preservar las cuatro semanas finales de evidencia y elevar refuerzo o cambio antes del hito. \\
Telemetría / R-028 & Acordar finalidades, acceso y controles con Operaciones, TI y sindicato antes de incorporar cada ola; sin cámaras en cabina ni control de jornada no acordado. & JP e IMP con las contrapartes funcionales; SEG verifica privacidad. & Acta previa a la ola, prueba de controles y restricción de usos no acordados. Diferenciar el desacuerdo institucional de la adopción individual; conservar evidencia de decisión. \\
\end{tablalafrox}
\end{landscape}
\endgroup


\clearpage
\subsection*{Cobertura de condiciones de operación y diseño}
\addcontentsline{toc}{subsection}{Cobertura de condiciones de operación y diseño}
La Tabla 8.D.2 vincula diez temas con sus fuentes en el Caso 02 y la
Aclaración 2, 8.2. Utiliza los identificadores vigentes e indica la evidencia
necesaria para verificar la implementación y eficacia, incluidas las
validaciones que corresponden a S4, S9 y S13.

\begingroup
\renewcommand{\thetable}{8.D.2}
\begin{landscape}
\setlength{\textwidth}{\linewidth}
\renewcommand{\arraystretch}{1.08}
\begin{tablalafrox}{Condiciones de operación y diseño y evidencia requerida. Fuente: Caso 02 y Aclaración 2.}{tab:sd805-cobertura}{F{0.15} F{0.16} F{0.24} F{0.14} Y}{\cab{Tema / referencia vigente} & \cab{Fuente de la condición} & \cab{Evento y respuesta} & \cab{Responsable} & \cab{Evidencia y estado}}
Conocimiento del planificador / R-014 & Caso, cap. 19; retiro y efecto OTIF en T-16. & Transferencia incompleta antes del retiro; recuperar y validar reglas/excepciones. & IMP; CAL verifica. & V-14 y prueba sin planificador en piloto; conservar evidencia del efecto de 6--9 puntos OTIF. Estado: cobertura definida; la eficacia se mide en ejecución. \\
Objeción sindical / R-028 & Caso, caps. 10, 16.1 y 19. & Desacuerdo sobre finalidades/controles antes de la ola; acuerdo 9.7.2, sin vigilancia excluida. & JP/IMP; Operaciones, TI y sindicato. & Acta de finalidades, accesos y controles antes de la ola; cobertura pendiente ante desacuerdo. Estado: cobertura definida; la eficacia se mide en ejecución. \\
Conductores externos / R-035 & Caso, caps. 10, 15 y 19; RT-22.04. & Rechazo individual pese a acuerdo e incorporación; demostración en andén y conductor certificado/dispuesto. & IMP; Operaciones y transportista. & Uso práctico antes de salida y evidencia de adopción; distinto de firma R-011 y reemplazo R-024. Estado: cobertura definida; la eficacia se mide en ejecución. \\
Rotación nocturna / R-029 & Caso, caps. 7.2, 10 y 19. & Altas superan tutoría o llegan sin competencia; inducción y cobertura por turno. & IMP; SRE desde 22; Bodega/Operaciones. & Evaluación previa a autonomía y aprendizaje máximo dos horas; matriz 8.F.2a. El 38\% no es p0. Estado: cobertura definida; la eficacia se mide en ejecución. \\
Interfaces legadas / R-030 & Caso, caps. 5 y 19. & Contrato/extracción no verificables a tiempo; ensayo temprano ERP/WMS sin modificar ERP. & ARQ/DAT; TI; CAL verifica. & Entradas/salidas, tiempos y conciliación contra origen; matriz 8.F.2a. Distinto de calidad de historia. Estado: cobertura definida; la eficacia se mide en ejecución. \\
Enlace de Concepción / R-017 & Caso, caps. 6 y 19. & Interrupción antes de validar respaldos o pérdida del último camino; ensayo y restitución. & SRE; TI y operadores de enlace. & Conmutación y recuperación probadas antes de aceptar la ola; respaldo ofertado no equivale a disponible. Estado: cobertura definida; la eficacia se mide en ejecución. \\
Peak y cierre / R-031 & Caso, caps. 10, 17.5 y 19. & Condición de cierre incumplida bajo peak/congelamiento; continuidad y recuperación autorizadas. & JP; CAL y Operaciones. & Seis condiciones acumulativas y cuatro semanas finales; matriz 8.F.2a. No duplica capacidad CT-14. Estado: cobertura definida; la eficacia se mide en ejecución. \\
Obsolescencia / R-032 & Aclaración 2, 8.2; plataforma declarada. & Cambio excepcional de soporte/compatibilidad que excede mantenimiento ordinario; transición probada. & ARQ; DES/SRE. & Inventario y revisión trimestral de soporte, regresión y reversión; actualización ordinaria excluida de contingencia. Estado: cobertura definida; la eficacia se mide en ejecución. \\
Bloqueo por proveedor / R-033 & Aclaración 2, 8.2; servicios gestionados. & Restricción efectiva de exportación o sustitución; probar restauración y transición utilizable. & ARQ/SRE; TI. & Credenciales, exportaciones e infraestructura como código; prueba de restauración. Elegir AWS no acredita bloqueo. Estado: cobertura definida; la eficacia se mide en ejecución. \\
Escalabilidad / R-036 y CT-14 & Aclaración 2, 8.2; Caso, caps. 17.5 y 19; perfil que se confirma con Operaciones/TI. & Divergencia de concentración/mezcla frente al perfil validado; confirmar, comparar y repetir CT-14. & ARQ/CAL; Operaciones/TI; SRE desde 22. & Perfil versionado e informe de carga/estrés: ERP, colas, errores, p95 y recuperación; habilitar sólo alcance probado. Estado: cobertura definida; la eficacia se mide en ejecución. \\
\end{tablalafrox}
\end{landscape}
\endgroup

\paragraph{Decisión ante el perfil de carga}
Si el perfil es conforme, CAL verifica CT-14 y la autoridad competente habilita
sólo el alcance probado; R-036 conserva su hipótesis de exposición en la ventana.
Si existe variación pendiente de confirmar, ARQ conserva perfil y mediciones,
solicita validación a Operaciones/TI, aplica la respuesta de R-036 y mantiene
la promoción afectada pendiente hasta revalidar capacidad.
Si hay saturación dentro de la carga comprometida, LafroX registra un problema,
corrige el diseño y repite CT-14 con cargo al trabajo ordinario. La concentración
o mezcla no detalladas no convierten por sí solas esa carga en ampliación de
alcance: JP contrasta requisitos y acuerdos antes de cualquier decisión contractual.
Una misma causa recibe un único ID principal; R-030 atiende la interfaz,
R-031 la coordinación del cierre y R-036 la divergencia del perfil, vinculando
efectos sin duplicar HH. El peak conocido y el crecimiento contratado permanecen
en el dimensionamiento obligatorio de CT-14.


\clearpage
\subsection*{Correspondencia de continuidad y evidencia}
\addcontentsline{toc}{subsection}{Correspondencia de continuidad y evidencia}
La matriz 8.D.3 enlaza diseño, riesgo/control y pruebas de S4/S9. Se consultan
las ramas corregidas vigentes de S4 y S9.
Los commits identifican versiones documentales, no implementaciones probadas.
No se dispone aquí de configuraciones de producto, trazas, mediciones ni actas
de ejecución o aceptación de los tres alcances; se registran como evidencia por obtener.

\begingroup
\renewcommand{\thetable}{8.D.3}
\begin{landscape}
\setlength{\textwidth}{\linewidth}
\renewcommand{\arraystretch}{1.08}
\begin{tablalafrox}{Continuidad, diseño, prueba y evidencia requerida. Fuente: S4, anexos 4-M/4-V y 4.3.2; S9, T-13/T-17.}{tab:sd806-correspondencia}{F{0.11} F{0.18} F{0.19} F{0.15} Y}{\cab{Tema / S8} & \cab{Referencia exacta S4} & \cab{Protocolo y campaña S9} & \cab{Responsable} & \cab{Evidencia requerida y estado}}
ERP / R-034 & Anexo 4-M, AL-DTE-01 y AL-DR-01; 4.3.2, recuperación de Talca y ERP; Anexo 4-V, correspondencia de aceptación. & T-13, 7.5: AL-DTE-01 en resiliencia 6.3.1/6.3.2; AL-DR-01 en recuperación 6.4.1/6.4.2, antes de H5/H10. T-17, 3.1. & ARQ/SRE; TI, Finanzas y Operaciones; CAL verifica y JP registra. & Versión de carga, folio, guía/anulación, acuse y tiempos de 96 salidas, con ERP disponible, indisponible y cambio de carga. RTO de 4 h y RPO de 15 min medidos. Bloquear sin guía válida; recuperar WMS no recupera ERP. Estado: diseño definido; implementación, ensayo y acta en la etapa correspondiente. \\
Stock / CT-15 & Anexo 4-V, AL-STOCK-01; Anexo 4-M, AL-DR-01; 4.3.2, conmutación y retorno. & T-13, 7.5: AL-STOCK-01 en integración 6.1.1 (E1) y 6.1.2 (E2); AL-DR-01 en 6.4.1/6.4.2. T-17, 3.1. & Equipo DAT hasta mes 20; SRE desde mes 21, con apoyo puntual de DAT; ARQ diseña; CAL verifica; JP registra. & Estado central/local con retención durable, última unidad, UUID/época, acuses perdidos, exclusión del escritor previo, retorno y conciliación. Preventa real E1; preventa, portal y EDI reales E2. Sin doble reserva ni stock negativo; simulador no acredita E2. Estado: diseño definido; implementación, ensayo y acta en la etapa correspondiente. \\
Identidad / CT-03 & Anexo 4-M, AL-ID-01 y AL-OFF-01; Anexo 4-R, controles de identidad; Anexo 4-V, correspondencia. & T-13, 7.5: AL-ID-01 y AL-OFF-01 en resiliencia 6.3.1/6.3.2 antes de H5/H10. T-17, 3.1; seguridad conserva puerta G5. & SEG responde; ARQ diseña; DES/SRE ejecutan con TI/Operaciones; CAL verifica. & Manifiesto de 26 h, corte a una hora, relevos 8/16 h; ruta de 14,5 h y máximo 16 h. Probar expiración, firma, turno, permisos, reinicio/reloj y baja local/central aislada; cero accesos indebidos, cola conservada y conciliación única. Sin PIN ni secretos en expediente. Estado: diseño definido; implementación, ensayo y acta en la etapa correspondiente. \\
\end{tablalafrox}
\end{landscape}
\endgroup

\paragraph{Decisión por alcance y etapa}
CAL comprueba versión, configuración, entradas, resultado esperado y observado,
incidencias y repetición tras corrección; JP registra fecha y próxima revisión.
Un protocolo previsto o un PDF compilado no sustituyen una ejecución.
Resultado insuficiente o fallido mantiene el bloqueo y exige corregir y repetir.
Sólo evidencia satisfactoria y aprobación competente habilitan el alcance
concreto; T-17, 3.1, integra el expediente H5/H10 y la Contraparte Técnica acepta
con acta. Un bloqueo correcto de salida sin guía no demuestra continuidad de las
96 salidas. Los controles son obligaciones ordinarias del proveedor, sin nuevas
HH de contingencia. CAL verifica la implementación con resultados satisfactorios y las actas exigidas.

\clearpage
\section{Anexo 8.E: Controles técnicos preventivos}
Los 15 controles vigentes (CT-01--CT-15) establecen obligaciones de diseño, innovación y calidad de LafroX. CT-03 reúne identidad conectada y local conforme a S4/S9; CT-15 conserva stock, antes identificado como CT-16. Los 15 identificadores vigentes se mantienen; la ampliación de identidad no crea otra contingencia ni renumera stock. Cada ficha define una solución técnica, su responsable y la evidencia necesaria para verificarla. La habilitación requiere ejecutar las pruebas y aprobar la evidencia de cumplimiento. CT sin número designa la Contraparte Técnica.

CAL verifica la evidencia y JP controla el cierre. SRE asume el trabajo operativo de INN-03 (familia DAT) desde el mes 21 y dirige el despliegue por rutas de INN-05 (familia IMP) desde el mes 22. La Tabla 8.1 identifica los responsables por familia. No se promueve un flujo con defectos críticos o altos abiertos ni con pruebas habilitantes fallidas. Una función opcional que no supere sus pruebas permanece deshabilitada y se conserva el alcance obligatorio.

\begingroup
\renewcommand{\thetable}{8.E.1}
\begin{landscape}
\setlength{\textwidth}{\linewidth}
\renewcommand{\arraystretch}{1.08}
\begin{tablalafrox}{Controles técnicos: solución, responsable y verificación. Fuente: requisitos de arquitectura, seguridad, innovación y calidad de la solución.}{tab:cierre-ci}{G{0.08} F{0.25} Y F{0.34}}{\cab{Control} & \cab{Responsable} & \cab{Solución requerida} & \cab{Evidencia y condición de cierre}}
CT-01 & ARQ; DES implementa; CAL verifica. & Idempotencia persistente, consulta de estado antes de reintentar, conciliación y estado pendiente sin confirmar. & Simular timeout antes/después del cargo y reenvío de la misma clave: una sola imputación y conciliación verificable. Antes de H5/H10 en el alcance aplicable; CAL bloquea promoción del flujo fallido. \\
CT-02 & ARQ; SRE opera; CAL verifica. & Definir circuit breaker, caché de rutas autorizadas, límites y estado degradado ante caída del servicio de mapas elegido, sin dar nuevas rutas por calculadas. & Inyectar caída y recuperación del proveedor: continuidad autorizada de rutas disponibles, nuevas consultas pendientes y restitución sin pérdida. Antes de H5/H10 y tras cambio de proveedor; corregir el diseño si falla. \\
CT-03 & SEG responde; ARQ diseña; DES/SRE ejecutan; CAL verifica. & Autenticación y sesión conectadas y locales, mínimo privilegio, detección, revocación y auditoría. Manifiesto de bodega de 26 h; al aislar una hora después quedan 25 h. Verificar relevos a las 8/16 h, firma, sitio, dispositivo, rol, PIN y turno. Reinicio/reloj no extienden vigencia. Reparto de 14,5 h con autorización máxima de 16 h desde emisión. & AL-ID-01 y AL-OFF-01 en resiliencia 6.3.1/6.3.2 antes de H5/H10; seguridad sin hallazgos críticos/altos antes de G5. Exigir cero accesos indebidos, cola conservada y conciliación única. Distinguir baja local inmediata y baja central durante aislamiento; verificar el detalle de esta ficha y 8.D.3. Corregir y repetir antes de habilitar. \\
CT-04 & ARQ; SEG valida autenticidad; CAL verifica. & Validar origen, integridad, correlación y repetición antes de procesar EDI; poner mensajes inválidos en cuarentena sin producir pedidos o avisos. & Pruebas negativas con firma/origen alterados, mensaje repetido y correlación inválida: rechazo trazable sin efecto de negocio. Antes de H10 para EDI y regresión del alcance aplicable; corregir antes de promover. \\
CT-05 & SRE; SEG define política; CAL verifica. & Proteger almacenamiento y sesiones, aplicar enrolamiento y capacidades de bloqueo/gestión remota disponibles; revocar accesos tras pérdida y delimitar exposición. & Ensayo de terminal extraviado en línea y sin señal: datos protegidos, acceso no autorizado denegado y revocación/bloqueo al recuperar enlace. Antes de H5/H10 y tras cambio de política; no presumir borrado remoto sin conectividad. \\
CT-06 & SEG; SRE ejecuta; CAL verifica. & Aplicar límites, protección perimetral, aislamiento de portales y alertas; conservar capacidad para usuarios legítimos y continuidad de servicios críticos. & Inyectar ráfaga controlada en ambiente autorizado y medir bloqueo, latencia y disponibilidad contra umbrales de S4/T-13. Antes de H5/H10 y tras cambio perimetral; ningún ensayo ofensivo no autorizado en Producción. \\
CT-07 & SEG; SRE ejecuta; CAL verifica. & Restringir privilegios, separar funciones, preservar auditoría y respaldos y comprobar recuperación ante cambios no autorizados de una cuenta administrativa. & Ensayar cambio indebido, detección, contención y restauración con integridad demostrada; recuperar conforme RTO/RPO exigidos. Antes de H5/H10, restauraciones y pruebas de seguridad periódicas de S9; no cerrar con evidencia comprometida. \\
CT-08 & DAT dirige; IMP coordina piloto; CAL verifica. & Validar catálogo con Comercial/Calidad en 7--9, mostrar confianza y confirmar saldo en visita. Mantener alcance obligatorio; si no se valida, usar sólo aviso informativo por vida útil remanente. & Paquete 11--15: H5 en 11--12 y piloto real en 13--15. I-01A en mes 15: al menos 70\% de aciertos sobre al menos 100 avisos evaluables. Sin cumplir, no ampliar INN-01 a otras rutas; corregir y repetir. Registrar los resultados del piloto. \\
CT-09 & CAL; DES implementa; SRE aporta incidentes. & Comparar reproducción con incidente de terreno y limitar captura de batería/datos. Conservar diagnóstico convencional y regresión base si la función opcional no demuestra utilidad. & Conservar casos incidente--reproducción y medición de I-02A y consumo contra lo comprometido en S13, antes de ampliar captura; corregir captura/reproducción fallida y no acreditar el indicador sin evidencia. \\
CT-10 & SEG; DES implementa; CAL verifica. & Sanitizar por reglas fijas antes de sacar paquetes del terminal; minimizar datos y modelar amenazas. Suspender conjuntos afectados y producir muestra protegida ante filtración. & Pruebas con datos personales sembrados y campos no previstos: ningún dato personal sin sanitizar abandona el terminal; informe de revisión de amenazas y regresión. Bloquear la captura/exportación hasta aprobar, antes del primer uso y ante cada cambio de esquema. \\
CT-11 & DAT; Calidad del CLIENTE valida parámetros; CAL verifica. & Validar familias, mantener RNG-04 y no extender vencimiento por inferencia; presentar la estimación sólo en M9. Familia no validada conserva fecha impresa e historia visible. & Evidencia de I-03B por familia conforme a S13 antes de habilitar orden sugerido, con pruebas de no relajación de RNG-04. CAL bloquea familia no validada, DAT corrige y repite; la participación del CLIENTE no traslada responsabilidad de diseño. \\
CT-12 & DAT; DES implementa; CAL verifica. & Vincular lote--sensor mediante eventos de ubicación M2 y declarar tramos sin dato; impedir estimación cuando la historia esté incompleta. & Ensayar movimiento sin evento, sensor equivocado y sincronización tardía: historia incompleta visible y ninguna vida útil inferida sin fundamento. Antes de habilitar INN-03 y en regresión; corregir asociación antes de promover. \\
CT-13 & IMP; DAT/DES implementan; SEG revisa privacidad; CAL verifica. & Codiseñar hoja comprensible; aplicar umbral de agregación y supresión en cada emisión comparativa y revisión legal previa. Conservar versión de un bloque si no se entiende y bloquear bloque 4 si expone terceros. & Prueba I-05B con usuarios y pruebas de inferencia conforme S13 antes de habilitar cada bloque; informe de comprensión y supresión. Corregir y repetir antes de habilitar. \\
CT-14 & SRE; ARQ dimensiona; DES corrige; CAL verifica. & Comparar perfil de carga acordado y observado, concentración por ventana y mezcla de operaciones; demostrar capacidad efectiva por dependencia, incluidos ERP, colas y recursos compartidos, peak conocido y crecimiento comprometido. Ajustar el dimensionamiento ofertado sin degradar requisitos. & Pruebas de carga/estrés y p95 según S4/T-13, incluyendo peak y cuello de botella ERP, antes de H5/H10 y preparación de peak. Medir capacidad, colas, errores y recuperación; bloquear promoción si incumple. Conservar perfil, entradas, resultados y aprobación. Ante prueba fallida, LafroX corrige y repite antes de promover; no presume demanda adicional del CLIENTE. R-036 trata divergencias del perfil; R-031, coordinación del cierre; R-030, interfaces legadas. \\
CT-15 & Equipo DAT hasta mes 21; SRE desde mes 22; ARQ diseña; CAL verifica. & Una reserva central confirmada exige retención local durable de igual UUID, época y cantidad. Mantener una sola autoridad de escritura por ámbito; excluir escritor previo al promover, conciliar sitio/UUID y conservar capturas pendientes sin confirmar desde copia vencida. & AL-STOCK-01: última unidad concurrente, acuse perdido, reintento, cancelación y época obsoleta, sin doble reserva ni stock negativo; preventa real E1 (6.1.1), portal/EDI reales E2 (6.1.2). AL-DR-01: pérdida, promoción, retorno y conciliación (6.4.1/6.4.2), antes de H5/H10. CAL bloquea flujo inconsistente y exige corrección/regresión; simuladores no acreditan producto E2. \\
\end{tablalafrox}
\end{landscape}
\endgroup

\subsection*{Referencias de diseño y pruebas previstas}
\addcontentsline{toc}{subsection}{Referencias de diseño y pruebas previstas}

Las referencias siguientes localizan el sustento técnico de cada control en las versiones corregidas de S4, S9 y S13. Identifican diseño y verificación por ejecutar, no resultados aprobados. La ficha de 8.E.1 conserva la condición de cierre. Cuando sólo existe cobertura general, se declara la dependencia para que la rama responsable complete el protocolo específico y su trazabilidad antes de habilitar el flujo.

\begin{description}
\item[CT-01.] S4, «Módulo M7 Cobranza y rendición», en «Módulos funcionales y límites de contexto»: consulta del pago incierto por la misma clave, sin nuevo cobro ni autorización por silencio. T-13 de S9, 7.2, contempla la dependencia externa sin respuesta, con error o lenta; es cobertura general, no un caso POS específico acreditado. ARQ/DES deben precisar el contrato con la pasarela en S4 y CAL incorporar en S9 los casos de timeout, reenvío y conciliación de la ficha.
\item[CT-02.] S4, Anexo 4-D, «Módulo M4 Rutas», integración con GIS: API externa y caché de mapas por zona. La continuidad sin conexión que allí se describe pertenece al diseño previo y no se reincorpora como solución vigente. T-13, 7.2, aporta sólo el escenario general de dependencia externa; no acredita un protocolo de circuit breaker. ARQ/SRE deben completar en S4 los límites y estados degradados del diseño conectado y CAL su caída/recuperación específica en S9.
\item[CT-03.] S4, Anexo 4-M, AL-ID-01 y AL-OFF-01, y Anexo 4-R, controles 5.15/5.18, 8.5 y 8.15/8.16: identidad conectada y local, permisos, sesión, bajas y auditoría. S9, T-13, 7.5, asigna resiliencia 6.3.1/6.3.2 y define el expediente de AL-ID-01; T-17, 3.1, conserva versión y actas. La ficha incluye esa continuidad local; los controles de seguridad mantienen sus pruebas y puerta G5. La previsión documental no acredita implementación ni resultados.
\item[CT-04.] S4, Anexo 4-R, protección del canal AS2 mediante firma, cifrado S/MIME y MDN. Este sustento de autenticidad no acredita por sí solo cuarentena, correlación ni rechazo de repeticiones; ARQ/SEG deben completar esa trazabilidad de diseño en S4 y CAL los casos negativos EDI en S9, para E2.
\item[CT-05.] S4, Anexo 4-R, controles 8.7/8.8 y 8.24 y «Ciclo de identidad y custodia»: protección del terminal, cifrado y revocación conectada. SRE/SEG deben concretar en S4 las capacidades de gestión remota disponibles; CAL debe vincular en S9 el ensayo de extravío, sin presumir bloqueo remoto durante la pérdida de enlace.
\item[CT-06.] S4, Anexo 4-R, control 8.20/8.22: borde y superficie B2B separada. T-13, 7.1, define los umbrales de carga, no un ensayo específico de denegación de servicio. SEG/SRE deben completar límites y protección en S4 y CAL vincular en S9 la ráfaga autorizada y las mediciones de la ficha.
\item[CT-07.] S4, Anexo 4-R, controles 8.2, 8.13/8.14 y 8.15/8.16: segregación, respaldo y auditoría; T-13, 7.5, AL-DR-01, aporta recuperación del WMS. No equivale a probar abuso administrativo. SEG/SRE deben completar el escenario de contención/restauración en S4 y CAL su verificación específica en S9. Para CT-04--CT-07 no se acredita aún correspondencia individual con casos de S9; SEG y CAL deben registrarla, sin cerrar los controles con estas referencias generales.
\item[CT-08.] S13, 13.1.6, I-01A, y 13.1.7, «Control de adopción»: validación del catálogo, certificación y piloto antes de ampliar rutas. El indicador y el piloto son evidencia prevista, no aciertos ya observados.
\item[CT-09 y CT-10.] S13, 13.2.6, I-02A/I-02B e I-02D, y 13.2.7, «Controles de adopción y seguridad»: CT-09 cubre fidelidad y consumo; CT-10, sanitización previa a exportar. CAL debe conservar la comparación incidente--reproducción y el consumo medido; SEG/CAL, la prueba negativa de datos prohibidos en cada liberación y cambio de esquema. Los indicadores no sustituyen esos registros.
\item[CT-11 y CT-12.] S13, 13.3.6, I-03A/I-03B, y 13.3.7, «Controles de validación por familia y trazabilidad térmica»: CT-11 valida por familia sin relajar RNG-04 ni extender la fecha impresa; CT-12 ensaya movimiento sin evento, sensor equivocado y sincronización tardía, con historia incompleta visible y sin estimación. DAT/CAL deben aportar los resultados antes de habilitar cada familia.
\item[CT-13.] S13, 13.5.6, I-05B/I-05D, y 13.5.7, «Control de comprensión y privacidad»: comprensión del prototipo, revisión legal y umbral/supresión en cada emisión comparativa. IMP/SEG/CAL deben registrar resultados antes de habilitar los bloques respectivos; la prueba de comprensión no acredita por sí sola privacidad.
\item[CT-14.] S4, «Primer cuello de botella bajo la carga de septiembre», tabla «Detección del primer cuello de botella»: ERP por ACL, base de Talca, cola y M2. T-13, 7.1, especifica cargas, ventana sostenida, crecimiento/quiebre, p95 e informe por componente; 7.5, AL-PERF-01, complementa la carga con aislamiento por perfil. SRE/ARQ/CAL deben contrastar la hipótesis con mediciones, no con el cálculo de dimensionamiento.
\item[CT-15 (stock).] S4, Anexo 4-V, AL-STOCK-01, y Anexo 4-M, AL-DR-01, con procedimientos de conmutación/retorno de 4.3.2: retención local durable, exclusión del escritor previo y conciliación. T-13 de S9, 7.5, asigna 6.1.1 a preventa E1, 6.1.2 a preventa/portal/EDI E2 y 6.4.1/6.4.2 a recuperación. T-17, 3.1, conserva el expediente y el resultado por etapa. Un simulador no acredita producto E2. CT-16 histórico de stock equivale a CT-15 vigente; la identidad local antes separada se integra en CT-03, sin reutilizar CT-15 para identidad.
\end{description}

CAL debe incorporar al expediente aplicable la versión ensayada, entradas, resultados, defectos y acta; JP mantiene pendientes las dependencias de diseño/prueba hasta su resolución en las ramas responsables. Esta trazabilidad documental no acredita implementación, ejecución ni aceptación técnica.

\subsection*{Condiciones de identidad local de CT-03}
El manifiesto de bodega dura 26 h y el ensayo corta el enlace una hora después
de emitirlo: quedan 25 h verificables. La caché descriptiva de 24 h no emite
credenciales ni amplía el manifiesto. En los relevos a las 8 y 16 h del corte se
validan firma, emisor, sitio, dispositivo enrolado, rol, PIN y turno vigente;
sin enlace no se habilitan cuentas no preinscritas ni privilegios nuevos.
Manifiesto o turno vencido, permiso inválido, firma alterada o PIN incorrecto
se deniegan y auditan. Reinicio, actividad o retroceso del reloj no extienden
vigencia; sin continuidad demostrable de la hora confiable se deniega sesión nueva.

En reparto, la autorización se emite dentro de la hora previa a salir, cubre
la ruta de 14,5 h y expira al fin programado más 30 min, con máximo de 16 h
desde su emisión. Al vencer, se deniegan nuevas operaciones protegidas y se
conserva la cola. La reautenticación conectada permite conciliar una sola vez
los hechos anteriores al vencimiento; otra ruta exige autorización nueva.

Una baja conocida en el sitio deniega acceso local inmediato. Una baja conocida
sólo por la autoridad central durante el aislamiento deja exposición hasta la
expiración: no se promete revocación remota instantánea. Al reconectar se bloquea
la sesión y se ponen en cuarentena los eventos posteriores a la baja para
conciliación, sin borrarlos. El expediente de AL-ID-01 conserva emisión,
expiración, huella del manifiesto, política, versión del verificador, tiempos,
decisiones y UUID; excluye PIN y secretos. CAL y SEG verifican resultados y
la aprobación competente habilita sólo el alcance probado.

\clearpage
\section{Anexo 8.F: Respuestas, esfuerzo y riesgo residual proyectado}
\label{sec:cierre-respuestas}
Las 36 fichas complementan 8.A y T-16. A es esfuerzo preventivo adicional total en la ventana, no costo de desarrollo ni pruebas ordinarias de la EDT. Se estima una revisión extraordinaria de excepción/coordinación/cierre por riesgo: 60\% del dueño y 40\% de CAL (ARQ cuando CAL es dueño), acumulados en toda la ventana. Para vigilancia ordinaria, aceptación pasiva y confirmación ordinaria del perfil de R-036, A=0. C es esfuerzo de respuesta condicional total, repartido en esos perfiles y contenido en el PERT H0 de 8.B, no sumado de nuevo ni multiplicado otra vez por episodios n. Son hipótesis de dos perfiles, no compromisos ni rendimientos observados. JP debe contrastar estos esfuerzos con los paquetes existentes antes de aprobarlos; si una acción ya está presupuestada, se descuenta aquí y se recalculan sumas.

El plazo preventivo es el de 8.A.2 para las 29 dependencias, antes del hito/fecha allí indicado; para los siete externos se fija en cada ficha. A se ejecuta una vez como presupuesto acumulado y se programa a D-5 del primer hito afectado; si hay revisiones repetidas (p. ej. rotación), se reparte ese total, sin multiplicarlo por 44 meses ni por n. Las HH preventivas mínimas no sustituyen capacitación, mantenimiento, respaldo o controles obligatorios incluidos en el presupuesto ordinario.

La probabilidad posterior p1 es un juicio de planificación de ocurrencia del mismo evento y ventana de p0, condicionado a que se ejecuten los controles y se obtenga la evidencia indicada. No es frecuencia medida ni eficacia certificada; se mantiene la probabilidad de dependencias externas que LafroX no puede modificar. En historia sanitaria, retirada del planificador, energía, SII, siniestros y cambio normativo no se reduce p0 por describir una respuesta. Los otros p1 proponen menor ocurrencia por resolver temprano la dependencia especificada; esa hipótesis se invalida si no existe su acta o prueba. I1 conserva el impacto global inicial, incluyendo consecuencias sanitarias y contractuales: ni el bloqueo de salida ni una mejora de detección borran una pérdida histórica. E1 se obtiene de la banda de p1 y el factor de I1 con los mismos umbrales de 8.1.3.

Para R-002, R-006, R-029, R-030 y R-031, p1 contempla el tratamiento completo de 8.F.1, incluido el plan adicional, sujeto a verificación. La valoración inicial de T-16.1 es la referencia anterior al tratamiento; no se ha estimado una valoración intermedia que aísle el efecto del plan adicional. Por tanto, no corresponde descontar otra reducción a p1 al activarlo. La justificación y las puertas de decisión se precisan después de la Tabla~\ref{tab:cierre-residual}.

El esfuerzo condicional remanente de atención es L=H0-C; no es una nueva estimación de consecuencias sanitarias ni de plazo. El esfuerzo residual esperado es p1 L. El tratamiento esperado es T=A+p1 C. El beneficio bruto en HH es B=p0 H0-p1 L. B/T compara reducción esperada de trabajo de atención con trabajo de tratamiento; el saldo neto B-T=(p0-p1)H0-A evita doble conteo. Así, costo total posterior=T+p1 L=A+p1 H0. Estos indicadores no valorizan alimentos, multas, seguridad de personas ni aceptación contractual; un retorno menor que uno se justifica por obligación o continuidad, sin eliminar el cálculo. Las respuestas contingentes no constituyen prevención adicional gratuita.

La base p0/O/M/P corresponde a las ventanas de exposición del Anexo 8.B. Los intervalos de programación de 8.H son hipótesis de distribución de HH, no cambian la exposición descrita en 8.B.2 ni convierten HH en días. D-5 y D-10 son días hábiles, salvo los minutos y turnos expresamente indicados. Los tiempos propuestos son máximos internos y se acortan cuando el SLA o el procedimiento de incidente exija actuar antes; no amplían compromisos de servicio ni autorizan indisponibilidad en ventanas críticas. Las fechas de respuesta no sustituyen los diez días hábiles contractuales de revisión/subsanación.

\begingroup
\renewcommand{\thetable}{8.F.1}
\begin{landscape}
\setlength{\textwidth}{\linewidth}
\renewcommand{\arraystretch}{1.08}
\begin{tablalafrox}{Fichas ejecutables de los 36 riesgos. Fuente: 8.A y estimaciones adicionales de planificación de S8.}{tab:cierre-fichas}{G{0.06} Y F{0.29} F{0.15} F{0.20}}{\cab{ID} & \cab{Acción preventiva y evidencia} & \cab{Contingencia y plazo desde disparador} & \cab{Perfiles / HH adicionales} & \cab{Momento y regla de programación}}
R-001 & Solicitar al Gerente Comercial, con TI y cada cadena, agenda EDI en mes 13, confirmada en mes 14, y certificación a D-5 del inicio de las cuatro semanas finales de evidencia de E2. Ante falta de respuesta, JP escala por CT; continuar pruebas independientes sin acreditar la cadena pendiente ni sustituir EDI por carga manual. Evidencia de salida: Agenda EDI y certificación de cada cadena, con pruebas de intercambio y registro de pendientes. & Al D-5 sin certificación, JP escala en un día hábil y conserva pendiente la cadena; sólo ampliar tras su certificación. & ARQ / CAL; A: 4,80/3,20; C: 9,60/6,40 HH. & Ventana de programación 13--20. A acumulado a D-5 del primer hito afectado; para repeticiones, repartir el total. Aplicar fechas de 8.A.2 y el plazo contingente de esta fila. \\
R-002 & Solicitar a la Jefa de Calidad, con TI y Bodega, evidencia de origen y decisión sobre lotes no reconstruibles; política V-07 al cierre del mes 2 y excepciones a D-5 de validar cada carga. Si no responde, JP escala por CT, se aísla el conjunto y queda pendiente su aceptación; no se inventa historia sanitaria. Evidencia de salida: Inventario de vínculos sanitarios por lote, evidencia de origen y resolución firmada de Calidad sobre cada excepción. & Plan adicional: DAT aísla el conjunto el mismo día; Calidad decide antes del corte. No cargar ni despachar lotes sin trazabilidad exigida; si no es reconstruible, elevar decisión de alcance contractual, sin inventar historia. Completar con 8.F.2a: DAT/CAL concilian vínculos y prueban bloqueo; Calidad valida V-07. Habilitar sólo lotes trazables; ante fallo, corregir y repetir antes del corte. & DAT / CAL; A: 7,20/4,80; C: 14,40/9,60 HH. & Ventana de programación 1--20. A acumulado a D-5 del primer hito afectado; para repeticiones, repartir el total. Aplicar fechas de 8.A.2 y el plazo contingente de esta fila. \\
R-003 & Solicitar al Jefe de Abastecimiento y a TI el historial de tres años y validación de V-09 al cierre del mes 3; corregir extractos a D-5 de validar la carga. Ante silencio, escalar por CT, separar períodos incompletos y continuar el perfilado disponible sin reducir la historia exigida. Evidencia de salida: Extracto de tres años perfilado y reconciliado con Abastecimiento/TI; acta V-09. & A D-5 sin extracto válido, separar períodos y pedir extracción complementaria; revalidar antes de cargar. & DAT / CAL; A: 3,60/2,40; C: 9,60/6,40 HH. & Ventana de programación 1--18. A acumulado a D-5 del primer hito afectado; para repeticiones, repartir el total. Aplicar fechas de 8.A.2 y el plazo contingente de esta fila. \\
R-004 & Solicitar a Bodega, con Operaciones y Finanzas, protocolo de conciliación en mes 2, validado en mes 3. Confirmar conteo a D-5 del corte y cerrar acta y ajustes antes de D. Si no hay respuesta, escalar por CT, repetir conciliación y bloquear el corte afectado hasta contar con evidencia aprobada. Evidencia de salida: Acta de conteo, causas de diferencias y ajustes autorizados por Bodega/Finanzas. & Bloquear corte con diferencia no explicada; DAT presenta conciliación y nueva prueba antes de la ventana autorizada. & DAT / CAL; A: 4,80/3,20; C: 12,00/8,00 HH. & Ventana de programación 1--20. A acumulado a D-5 del primer hito afectado; para repeticiones, repartir el total. Aplicar fechas de 8.A.2 y el plazo contingente de esta fila. \\
R-005 & Solicitar al Gerente de Operaciones y a Finanzas respaldos y conciliación V-08 de saldos por cliente y transportista; política en mes 2 y respuesta a D-5 de validar la carga M8 de cada ola. Ante silencio, escalar por CT y mantener M8 pendiente para esos saldos; no presumir saldo cero. Evidencia de salida: Saldos de envases aceptados por cliente y transportista, política V-08 y evidencia de M8. & Mantener saldos afectados pendientes; Finanzas resuelve antes de su carga, sin asumir saldo cero. & DAT / CAL; A: 3,60/2,40; C: 9,60/6,40 HH. & Ventana de programación 1--20. A acumulado a D-5 del primer hito afectado; para repeticiones, repartir el total. Aplicar fechas de 8.A.2 y el plazo contingente de esta fila. \\
R-006 & Solicitar a Comercial y Operaciones las horas históricas y a Finanzas los montos de rendición, con validación de V-10 al cierre del mes 1. Si no responden, escalar por CT y mantener pendiente la referencia de aceptación; la captura futura no reconstruye la historia anterior. Evidencia de salida: Fuentes históricas de horas y montos verificadas con Comercial, Operaciones y Finanzas, V-10. & Plan adicional: JP reúne áreas dentro de dos días hábiles desde bloqueo; en cinco días hábiles exige evidencia o eleva al Comité Ejecutivo la condición contractual. No sustituir historia por mediciones futuras ni aceptar resultado sin referencia. Completar con 8.F.2a: IMP/CAL reproducen OTIF y cuadratura; las áreas validan V-10. Sin fuentes conciliadas, mantener comparación pendiente y repetir tras corregir. & IMP / CAL; A: 4,80/3,20; C: 9,60/6,40 HH. & Ventana de programación 1--20. A acumulado a D-5 del primer hito afectado; para repeticiones, repartir el total. Aplicar fechas de 8.A.2 y el plazo contingente de esta fila. \\
R-007 & Solicitar a Finanzas y TI el universo de deuda vigente, con respuesta inicial al cierre del mes 3 y conciliación a D-5 de validar la migración. Ante silencio, escalar por CT y mantener la cartera incompleta pendiente; preparar extracción complementaria sin omitir deuda obligatoria. Evidencia de salida: Universo de deuda vigente conciliado con Finanzas/TI y extracción complementaria comprobada. & No dar por migrada cartera incompleta; escalar a D-5 de cada carga y repetir conciliación. & DAT / CAL; A: 3,60/2,40; C: 9,60/6,40 HH. & Ventana de programación 1--18. A acumulado a D-5 del primer hito afectado; para repeticiones, repartir el total. Aplicar fechas de 8.A.2 y el plazo contingente de esta fila. \\
R-008 & Mantener trazabilidad de trámites pendientes, consulta y reintentos idempotentes hasta recibir confirmación del SII. Ante demora externa, conservar el estado pendiente y escalar según el procedimiento acordado; el silencio no equivale a aprobación. Evidencia de salida: Cola tributaria con correlación, consulta y reintento idempotente; confirmación externa conservada. & SRE registra pendiente al detectar demora, consulta dentro de 15 min y revisa cada 30 min; si afecta salida, avisa de inmediato a Finanzas/CT. No interpretar silencio como aprobación. & SRE / CAL; A: 0,00/0,00; C: 4,80/3,20 HH. & Ventana de programación 13--56. A acumulado a D-5 del primer hito afectado; para repeticiones, repartir el total. Aplicar fechas de 8.A.2 y el plazo contingente de esta fila. \\
R-009 & Solicitar a la Jefa de Calidad la matriz térmica versionada V-11 durante mes 1 y su aprobación al cierre del mes 2, antes de certificar. Si no responde, escalar por CT y mantener el tratamiento preventivo de lecturas fuera de ficha; no cerrar marcha blanca ni aprobar parámetros por el CLIENTE. Evidencia de salida: Matriz térmica V-11 versionada y firmada por Calidad por familia antes de certificación. & Sin parámetros aprobados, bloquear certificación del alcance afectado y escalar a D-5; no relajar bloqueo térmico. & CAL / ARQ; A: 3,60/2,40; C: 7,20/4,80 HH. & Ventana de programación 1--20. A acumulado a D-5 del primer hito afectado; para repeticiones, repartir el total. Aplicar fechas de 8.A.2 y el plazo contingente de esta fila. \\
R-010 & Solicitar al Jefe de TI, con Operaciones, calendario de equipos en mes 1, entrega confirmada a D-5 y recepción técnica antes de la instalación o prueba D. Ante silencio o faltante, escalar por CT, continuar trabajo independiente y bloquear las actividades sin equipo. Evidencia de salida: Recepción técnica por equipo y sitio, con conformidad de TI/Operaciones antes de la ola. & A D-5 sin equipos, JP decide reprogramación de trabajo independiente en un día hábil; no certificar con sustitutos no validados. & IMP / CAL; A: 2,40/1,60; C: 7,20/4,80 HH. & Ventana de programación 1--20. A acumulado a D-5 del primer hito afectado; para repeticiones, repartir el total. Aplicar fechas de 8.A.2 y el plazo contingente de esta fila. \\
R-011 & Solicitar al Gerente de Operaciones y al representante de cada transportista el acuerdo y la nómina desde mes 2, firmados a D-5 de su incorporación a la ola. Ante silencio, escalar por CT y no habilitar a esa empresa; conservar su cobertura pendiente sin retirarla del alcance. Evidencia de salida: Acuerdo operativo firmado y nómina de cada transportista antes de incorporar la ola. & A D-5 sin firma, escalar por CT el mismo día; no habilitar empresa pendiente ni reducir cobertura obligatoria. & IMP / CAL; A: 2,40/1,60; C: 4,80/3,20 HH. & Ventana de programación 2--20. A acumulado a D-5 del primer hito afectado; para repeticiones, repartir el total. Aplicar fechas de 8.A.2 y el plazo contingente de esta fila. \\
R-012 & Solicitar a Operaciones y al propietario permiso de instalación por vehículo con D-10 de la visita y obtenerlo a D-5. Ante silencio, escalar por CT, suspender la intervención y proponer otra visita o vehículo autorizado; mantener pendiente la evidencia térmica. Evidencia de salida: Autorización escrita por vehículo y visita; propietario y Operaciones confirman disponibilidad. & A D-5 sin permiso, suspender visita y acordar alternativa autorizada antes de instalar. & IMP / CAL; A: 2,40/1,60; C: 7,20/4,80 HH. & Ventana de programación 1--20. A acumulado a D-5 del primer hito afectado; para repeticiones, repartir el total. Aplicar fechas de 8.A.2 y el plazo contingente de esta fila. \\
R-013 & Solicitar a Comercial y Operaciones, con Bodega por turno, responsables y agenda de acompañamiento a D-5 del inicio de cada ola, convocando a D-10. Si no responden o persiste el registro externo, escalar por CT, focalizar acompañamiento y mantener pendiente la evidencia de adopción. Evidencia de salida: Transacciones de muestra registradas en el flujo nuevo y cobertura certificada por turno/ola. & Si persiste flujo paralelo, IMP abre refuerzo en el siguiente turno y mide diariamente; sin evidencia de adopción no cerrar ola. & IMP / CAL; A: 4,80/3,20; C: 12,00/8,00 HH. & Ventana de programación 13--56. A acumulado a D-5 del primer hito afectado; para repeticiones, repartir el total. Aplicar fechas de 8.A.2 y el plazo contingente de esta fila. \\
R-014 & Solicitar al Gerente de Operaciones y al planificador confirmación de retiro y disponibilidad V-14 al cierre del mes 1; transferir y validar reglas con el ayudante en meses 1--3. Probar las rutas y excepciones sin el planificador en la marcha blanca de meses 13--15 (S7, 7.2.1). Si falta respuesta o se retira antes, escalar por CT y reconstruir con ayudante e historia verificable, sin dar por validadas reglas pendientes. Evidencia de salida: V-14, reglas y excepciones transferidas al ayudante y prueba en 13--15 sin planificador. & Retiro confirmado no cambia p0 del fallo de transferencia. Si faltan reglas, IMP inicia reconstrucción en dos días hábiles y valida antes de cerrar el piloto. & IMP / CAL; A: 3,60/2,40; C: 9,60/6,40 HH. & Ventana de programación 1--15. A acumulado a D-5 del primer hito afectado; para repeticiones, repartir el total. Aplicar fechas de 8.A.2 y el plazo contingente de esta fila. \\
R-015 & Solicitar a Operaciones, Bodega, Comercial y TI, según perfil, nómina y liberación horaria a D-5 de cada sesión convocada a D-10. Ante silencio o ausencia, escalar por CT y reprogramar evaluación; usar sólo suplentes certificados y no acreditar preparación ni cierre de ola sin evidencia. Evidencia de salida: Lista de convocados, asistencia, evaluación y suplentes certificados por sesión. & Ausencia: avisar a jefatura el mismo día y fijar recuperación dentro de cinco días hábiles, antes de trabajo autónomo. & IMP / CAL; A: 2,40/1,60; C: 7,20/4,80 HH. & Ventana de programación 1--56. A acumulado a D-5 del primer hito afectado; para repeticiones, repartir el total. Aplicar fechas de 8.A.2 y el plazo contingente de esta fila. \\
R-016 & Solicitar a CT área decisora y suplente, con agenda al cierre del mes 1; cursar insumos a D-10 y pedir respuesta a D-5. Las aceptaciones formales conservan diez días hábiles. Para promover la región de contingencia, avisar al suplente a los 15 min sin respuesta del titular y obtener decisión dentro de 30 min (S4, 4.3.2.5); no esperar el plazo ordinario. Mantener bloqueada la acción sin autorización; el silencio no permite cambios. Evidencia de salida: Matriz de decisión con titular y suplente; recepción de insumos y resolución trazables. & D-5 sin decisión: escalar un día hábil después; promoción regional: suplente a 15 min y decisión dentro de 30 min, sin actuar por silencio. & JP / CAL; A: 2,40/1,60; C: 7,20/4,80 HH. & Ventana de programación 1--56. A acumulado a D-5 del primer hito afectado; para repeticiones, repartir el total. Aplicar fechas de 8.A.2 y el plazo contingente de esta fila. \\
R-017 & Mitigar. Verificar respaldos antes de aceptar la ola de Concepción; no considerar disponible un enlace sólo por estar ofertado. Mantener y ensayar los enlaces ya ofertados: en Talca y Concepción, fibra y LTE con Starlink en espera caliente; en cross-docking, Starlink con LTE Entel/Claro y conmutación menor a 30 s. En Los Ángeles de madrugada, probar pérdida del camino Starlink restante y continuidad local de su ventana de 3 h. Coordinar restitución y conciliación sin duplicar el mismo corte como falla de SII, POS o mapas. Referencia: S4, 4.2.5.3. Evidencia de salida: Recepción de respaldos por sitio y ensayo de pérdida del último camino, incluida Concepción y Los Ángeles. & SRE activa continuidad al corte y avisa a TI/CT en 15 min; comprobar conmutación menor a 30 s, ventana local y restitución antes de normalizar. & SRE / CAL; A: 2,40/1,60; C: 7,20/4,80 HH. & Ventana de programación 1--56. A acumulado a D-5 del primer hito afectado; para repeticiones, repartir el total. Aplicar fechas de 8.A.2 y el plazo contingente de esta fila. \\
R-018 & Verificar autonomía, alarmas y procedimiento de recuperación de los componentes de LafroX, coordinando las dependencias de energía con el CLIENTE/proveedor. Si se excede la autonomía, activar recuperación y registrar la indisponibilidad e impacto. Evidencia de salida: Autonomía y transferencia verificadas por componente; coordinación de energía con TI/Operaciones. & Al corte, SRE mide autonomía restante y alerta en 15 min; activar recuperación antes de agotarla y bloquear operación insegura. & SRE / CAL; A: 0,00/0,00; C: 4,80/3,20 HH. & Ventana de programación 13--56. A acumulado a D-5 del primer hito afectado; para repeticiones, repartir el total. Aplicar fechas de 8.A.2 y el plazo contingente de esta fila. \\
R-019 & Solicitar por CT, con TI y la autoridad del CLIENTE encauzada por Gerencia General, aprobación escrita de residencia y transferencia en mes 1, resuelta al cierre del mes 2 o antes de habilitar datos, lo que ocurra primero. Ante silencio, escalar y preparar alternativa de arquitectura; mantener deshabilitada la transferencia no aprobada. Evidencia de salida: Aprobación escrita de residencia/transferencia para regiones propuestas antes de habilitar datos. & Sin aprobación al mes 2, JP escala en un día hábil y ARQ presenta alternativa en cinco días hábiles; no transferir datos sin autorización. & ARQ / CAL; A: 2,40/1,60; C: 7,20/4,80 HH. & Ventana de programación 1--2. A acumulado a D-5 del primer hito afectado; para repeticiones, repartir el total. Aplicar fechas de 8.A.2 y el plazo contingente de esta fila. \\
R-020 & Solicitar a CT, Operaciones y TI flujo y suplentes al cierre del mes 1; ante incidente, pedir autorización escrita antes del cambio. Sin respuesta, escalar en una hora o antes si lo exige el plazo de atención, manteniendo contención autorizada y estado seguro. El silencio no permite intervenir durante congelamiento. Evidencia de salida: Flujo de cambios urgentes y suplentes aprobados; autorización escrita por intervención. & Durante incidente escalar en una hora o antes si exige el servicio; contener sin intervención prohibida y conservar congelamientos. & SEG / CAL; A: 2,40/1,60; C: 7,20/4,80 HH. & Ventana de programación 1--56. A acumulado a D-5 del primer hito afectado; para repeticiones, repartir el total. Aplicar fechas de 8.A.2 y el plazo contingente de esta fila. \\
R-021 & Detener de forma segura el recorrido afectado, registrar los casos de prueba incompletos y coordinar con el CLIENTE/transportista una nueva ventana antes de repetir la validación. Evidencia de salida: Ventana y recorrido de prueba confirmados por Operaciones/transportista antes de salir. & Al aviso de falla, suspender prueba, notificar el mismo día y acordar recorrido dentro de dos días hábiles; repetir sólo con vehículo seguro. & IMP / CAL; A: 0,00/0,00; C: 4,80/3,20 HH. & Ventana de programación 1--20. A acumulado a D-5 del primer hito afectado; para repeticiones, repartir el total. Aplicar fechas de 8.A.2 y el plazo contingente de esta fila. \\
R-022 & Priorizar la atención del accidente conforme a los procedimientos del transportista y del CLIENTE; suspender la prueba, registrar su impacto y reprogramar sólo cuando exista autorización para continuar. Evidencia de salida: Autorización de recorrido y protocolo de seguridad del CLIENTE/transportista. & Suspender de inmediato ante siniestro; coordinar el mismo día y fijar nueva ventana dentro de dos días hábiles, condicionada a autorización segura. & IMP / CAL; A: 0,00/0,00; C: 4,80/3,20 HH. & Ventana de programación 1--20. A acumulado a D-5 del primer hito afectado; para repeticiones, repartir el total. Aplicar fechas de 8.A.2 y el plazo contingente de esta fila. \\
R-023 & Coordinar con el CLIENTE y el transportista la respuesta al asalto, robo o incendio, proteger a las personas y registrar qué evidencia y pruebas se perdieron. Acordar una nueva ejecución segura antes de continuar la validación. Evidencia de salida: Custodia y seguridad del transporte corresponden al CLIENTE/transportista; evidencia de prueba preservada. & Interrumpir ante asalto, avisar el mismo día y acordar reejecución dentro de dos días hábiles; no reanudar sin habilitación segura. & IMP / CAL; A: 0,00/0,00; C: 4,80/3,20 HH. & Ventana de programación 1--20. A acumulado a D-5 del primer hito afectado; para repeticiones, repartir el total. Aplicar fechas de 8.A.2 y el plazo contingente de esta fila. \\
R-024 & Solicitar a Operaciones y al transportista nómina y suplentes a D-5 de incorporar la empresa y confirmar conductor, vehículo y viaje antes de cada salida. Si falta respuesta o llega un sustituto, escalar, verificar en andén y solicitar conductor certificado; no habilitar uso no autorizado ni acreditar cobertura incompleta. Evidencia de salida: Identidad, autorización, capacitación y vínculo persona--vehículo--viaje verificados en andén. & Sustituto no certificado: bloquear acceso al detectar y solicitar otro conductor ese mismo turno; verificar antes de salida. & IMP / CAL; A: 2,40/1,60; C: 7,20/4,80 HH. & Ventana de programación 2--56. A acumulado a D-5 del primer hito afectado; para repeticiones, repartir el total. Aplicar fechas de 8.A.2 y el plazo contingente de esta fila. \\
R-025 & Solicitar por CT a la gerencia peticionaria solicitud escrita y decisión del cambio antes de su primera actividad dependiente; registrar y remitir el mismo día hábil. Ante silencio, mantener línea base y continuar trabajo independiente; si amenaza un hito, escalar al Comité Ejecutivo sin iniciar desarrollo no aprobado. Evidencia de salida: Solicitud y análisis de impacto aprobados antes de iniciar trabajo fuera de alcance. & Registrar solicitud el mismo día; JP responde o escala en dos días hábiles y mantiene línea base hasta decisión formal. & JP / CAL; A: 2,40/1,60; C: 7,20/4,80 HH. & Ventana de programación 1--56. A acumulado a D-5 del primer hito afectado; para repeticiones, repartir el total. Aplicar fechas de 8.A.2 y el plazo contingente de esta fila. \\
R-026 & Revisar cambios normativos aplicables, determinar su efecto con las áreas competentes y tramitar ajustes de controles, retención o eliminación mediante decisión y control de cambios documentados, sin posponer una obligación legal aplicable. Evidencia de salida: Seguimiento de cambios aplicables con asesoría competente; decisión contractual documentada. & ARQ registra aviso el mismo día hábil, solicita interpretación dentro de dos y presenta impacto dentro de cinco tras recibirla; no inventar multas o reformas. & ARQ / CAL; A: 0,00/0,00; C: 9,60/6,40 HH. & Ventana de programación 1--56. A acumulado a D-5 del primer hito afectado; para repeticiones, repartir el total. Aplicar fechas de 8.A.2 y el plazo contingente de esta fila. \\
R-027 & Solicitar a Finanzas, con Comercial y CT, reglas de 5.4.3 aprobadas al cierre del mes 20 y tres liquidaciones en sombra reproducidas en meses 21--23 antes del pago variable del mes 24. Ante silencio, no activar el tramo; una disputa se escala y suspende sólo el devengo correspondiente según contrato, manteniendo el servicio. Evidencia de salida: Reglas de costo de servir aprobadas al mes 20 y tres liquidaciones en sombra reproducidas en 21--23. & Antes del mes 24 sin acuerdo, no activar tramo variable; Finanzas/CT reciben disputa en un día hábil y mantienen servicio mientras resuelven. & JP / CAL; A: 3,60/2,40; C: 7,20/4,80 HH. & Ventana de programación 20--56. A acumulado a D-5 del primer hito afectado; para repeticiones, repartir el total. Aplicar fechas de 8.A.2 y el plazo contingente de esta fila. \\
R-028 & Mitigar y escalar. JP coordina con Operaciones, TI y sindicato el acuerdo 5.4.2 antes de la ola: convocar a D-10 y resolver a D-5. Explicar finalidades, acceso y controles; mantener excluidas las cámaras en cabina y no usar GPS para jornada sin acuerdo previo. Ante desacuerdo, escalar por CT, mantener sólo el uso operacional autorizado y registrar la cobertura de la ola pendiente. Evidencia de salida: Acuerdo 5.4.2 con finalidades, accesos y controles comprendidos por Operaciones/TI/sindicato. & D-5 sin acuerdo: JP escala el mismo día y conserva sólo uso operacional autorizado; no cámaras en cabina ni GPS de jornada no acordado. & JP / CAL; A: 3,60/2,40; C: 7,20/4,80 HH. & Ventana de programación 1--20. A acumulado a D-5 del primer hito afectado; para repeticiones, repartir el total. Aplicar fechas de 8.A.2 y el plazo contingente de esta fila. \\
R-029 & Mitigar. Mantener 7.1.2 en meses 13--56, turno 22:00--06:00, tutor y evaluación antes del trabajo autónomo, con aprendizaje máximo de dos horas. Bodega y Operaciones entregan altas por turno antes de asignarlas. IMP conduce hasta mes 21 y SRE desde 22, familia IMP. Ante sobrecarga, usar personal certificado, reforzar la inducción y escalar por CT; no duplicar las HH ordinarias de 7.1.2. Evidencia de salida: Altas por turno, tutor, evaluación antes de autonomía y aprendizaje dentro de dos horas según 7.1.2. & Plan adicional: Bodega asigna cobertura certificada antes del turno; IMP/SRE revisa cada alta diariamente. Si supera tutoría, no asignar trabajo autónomo; presentar refuerzo en dos días hábiles y revalidar antes de habilitar. Completar con 8.F.2a: IMP/SRE y CAL verifican cobertura y evaluación por alta. Si falta tutor, el aprendizaje excede dos horas o la evaluación no se aprueba, impedir autonomía y repetir formación/verificación. & IMP / CAL; A: 7,20/4,80; C: 19,20/12,80 HH. & Ventana de programación 13--56. A acumulado a D-5 del primer hito afectado; para repeticiones, repartir el total. Aplicar fechas de 8.A.2 y el plazo contingente de esta fila. \\
R-030 & Mitigar y escalar. ARQ solicita a TI acceso, diccionario y condiciones de extracción en mes 1 (1.2.3, 3.3.2 y V-16); DAT verifica la migración del WMS. Probar temprano sin modificar el ERP. Si falta una interfaz, evaluar intercambio de archivos controlado, sujeto a demostrar los requisitos; elevar por CT el bloqueo y sus impactos. La alternativa no acredita por sí sola latencia, integridad ni emisión tributaria. Evidencia de salida: Contrato de extracción e intercambio confirmado con TI; ensayo temprano sin modificar ERP y sustitución progresiva de WMS. & Plan adicional: ARQ eleva brecha al finalizar semana 2 del mes 1; TI entrega acceso o alternativa verificable antes de integración. Ensayar archivo controlado si satisface requisitos; bloquear certificación si persiste interfaz no verificable. Completar con 8.F.2a: ARQ/DAT ensayan con TI; CAL verifica integridad, latencia y emisión por ERP. Ante fallo, corregir y repetir antes de habilitar integración. & ARQ / CAL; A: 7,20/4,80; C: 14,40/9,60 HH. & Ventana de programación 1--18. A acumulado a D-5 del primer hito afectado; para repeticiones, repartir el total. Aplicar fechas de 8.A.2 y el plazo contingente de esta fila. \\
R-031 & Mitigar y escalar. JP acuerda con Operaciones y CT disponibilidad y continuidad antes de septiembre; cerrar habilitaciones en agosto según S7. En meses 19--20 mantener evidencia diaria de las seis condiciones de marcha blanca. Un defecto durante 1--25 de septiembre no autoriza despliegue: activar continuidad aprobada y R-020 si requiere autorización. Si persiste una condición incumplida, no declarar cierre; elevar plan de recuperación sin mover hitos por cuenta propia. Evidencia de salida: Disponibilidad y continuidad acordadas antes de septiembre; seguimiento diario de las seis condiciones de cierre. & Plan adicional: cerrar habilitaciones en mes 19 antes del congelamiento; JP revisa incidentes diariamente en 19--20. Escalar el mismo día una condición incumplida y acordar recuperación en dos días hábiles, sin intervención prohibida ni aceptación ficticia. Completar con 8.F.2a: JP/CAL verifican ensayo de continuidad previo al congelamiento y seis condiciones acumulativas. Ante fallo, mantener cierre pendiente y repetir en ventana autorizada. & JP / CAL; A: 4,80/3,20; C: 9,60/6,40 HH. & Ventana de programación 19--20. A acumulado a D-5 del primer hito afectado; para repeticiones, repartir el total. Aplicar fechas de 8.A.2 y el plazo contingente de esta fila. \\
R-032 & Mitigar. ARQ mantiene inventario de versiones y revisa soporte y compatibilidad trimestralmente desde mes 1, con DES y SRE. Preparar y probar adecuación antes del límite verificado de soporte, dentro de ventanas autorizadas. Si el cambio excede la actualización ordinaria 8.2.5, registrar el esfuerzo extraordinario y elevar impactos por CT; usar reversión ensayada y soporte permitido, sin sostener una versión insegura por conveniencia. Evidencia de salida: Inventario de versiones y soporte, revisión trimestral y compatibilidad probada antes del límite verificado. & Cambio excepcional: ARQ registra aviso en un día hábil y presenta transición en cinco; programar antes de fin de soporte, conservando ventanas y pruebas. & ARQ / CAL; A: 3,60/2,40; C: 12,00/8,00 HH. & Ventana de programación 1--56. A acumulado a D-5 del primer hito afectado; para repeticiones, repartir el total. Aplicar fechas de 8.A.2 y el plazo contingente de esta fila. \\
R-033 & Mitigar. ARQ verifica con TI propiedad de cuentas, acceso a datos, claves, exportaciones e infraestructura como código en la recepción inicial y en cada actualización anual del plan 8.1.4. Probar una restauración utilizable, no sólo descargar archivos. Ante restricción, documentar brecha, mantener continuidad autorizada y escalar por CT alternativas y decisión contractual; no declarar reversibilidad completa ni cambiar de proveedor sin prueba y aprobación. Evidencia de salida: Cuentas, claves, exportación y restauración utilizable verificadas en recepción y revisión anual de salida. & Restricción efectiva: ARQ escala en un día hábil y presenta alternativas en cinco; no sustituir proveedor sin prueba/aprobación ni declarar reversibilidad por una descarga. & ARQ / CAL; A: 2,40/1,60; C: 14,40/9,60 HH. & Ventana de programación 1--56. A acumulado a D-5 del primer hito afectado; para repeticiones, repartir el total. Aplicar fechas de 8.A.2 y el plazo contingente de esta fila. \\
R-034 & Mitigar y escalar. ARQ/SRE acuerdan con TI y Finanzas restitución del ERP y ensayan AL-DTE-01/AL-DR-01 antes de H5/H10, incluida pérdida de sala. Mantener sólo salidas con guía preemitida aún válida y bloquear las que requieren nueva guía. Recuperar WMS en nube no restaura el ERP ni exceptúa el RTO requerido; registrar y escalar por CT cualquier incumplimiento, sin segundo emisor ni aprobación tributaria por silencio. Evidencia de salida: Procedimiento ERP con TI/Finanzas probado en AL-DTE-01/AL-DR-01 antes de H5/H10. & Al fallo ERP avisar de inmediato a TI/Finanzas, mantener sólo guías preemitidas válidas y bloquear nuevas salidas sin guía. Activar restitución con RTO 4 h/RPO 15 min exigidos; escalar a los 30 min sin recuperación encaminada. & ARQ / CAL; A: 3,60/2,40; C: 12,00/8,00 HH. & Ventana de programación 13--56. A acumulado a D-5 del primer hito afectado; para repeticiones, repartir el total. Aplicar fechas de 8.A.2 y el plazo contingente de esta fila. \\
R-035 & Mitigar y escalar. IMP coordina con Operaciones, transportista y conductor demostración y acompañamiento antes de su salida, dentro de 7.1.3. Confirmar disponibilidad a D-5 de la ola y aceptación práctica en andén. Ante rechazo, solicitar a la empresa conductor certificado y dispuesto, mantener acompañamiento y elevar por CT la cobertura pendiente; no atribuir a LafroX facultades laborales ni declarar completa la evidencia faltante. Evidencia de salida: Demostración práctica en andén e incorporación acordada con transportista y conductor habilitado. & Ante rechazo actuar antes de salida: IMP solicita conductor certificado/dispuesto ese turno y registra cobertura pendiente; refuerzo antes de siguiente asignación. & IMP / CAL; A: 3,60/2,40; C: 9,60/6,40 HH. & Ventana de programación 2--56. A acumulado a D-5 del primer hito afectado; para repeticiones, repartir el total. Aplicar fechas de 8.A.2 y el plazo contingente de esta fila. \\
R-036 & Mitigar. ARQ solicita el perfil a Operaciones/TI en los primeros cinco días hábiles del mes 1 y lo valida al cierre de ese mes; revalida antes de H5/H10, preparación del peak y cambios del perfil. SRE asume gestión operativa desde mes 22, familia ARQ. Ante divergencia, conservar mediciones, ejecutar CT-14, bloquear promoción afectada y elevar recuperación por CT antes del hito. El peak conocido, crecimiento comprometido y corrección del diseño dentro del alcance siguen a cargo de LafroX. Evidencia de salida: Perfil versionado y validado por Operaciones/TI, comparación de carga acordada y observada por ventana y operación, informe CT-14 y decisión de habilitación del alcance verificado. & ARQ, y SRE desde mes 22, registra divergencia y avisa a JP/CT el mismo día hábil; presenta recuperación en dos días hábiles y siempre antes del hito. Conservar bloqueo, corregir y repetir CT-14 en ventana autorizada antes de habilitar. Si la divergencia pertenece al alcance comprometido, LafroX corrige como trabajo ordinario; no presume ampliación contractual. & ARQ / CAL; A: 0,00/0,00; C: 9,60/6,40 HH. SRE asume desde mes 22, familia ARQ. & Ventana de programación 1--56; A=0. Confirmación del perfil y CT-14 en trabajo ordinario. C es parte de H0, no se suma de nuevo; no multiplicar por revisiones ni por meses. \\
\end{tablalafrox}
\end{landscape}
\endgroup

\begingroup
\renewcommand{\thetable}{8.F.2}
\begin{landscape}
\setlength{\textwidth}{\linewidth}
\renewcommand{\arraystretch}{1.08}
\begin{tablalafrox}{Riesgo residual proyectado, condicionado a evidencia; no resultado observado. Fuente: hipótesis por riesgo y escalas de S8.}{tab:cierre-residual}{G{0.06} F{0.13} F{0.17} F{0.16} Y}{\cab{ID} & \cab{p1 / nivel / factor} & \cab{I1 / factor; E1 y zona} & \cab{Esfuerzo residual} & \cab{Fundamento, evidencia y decisión posterior}}
R-001 & 25,00\% / 2 / 0,30 & 4 / 0,40; 0,12 (Media) & L=88,00 HH; p1L=22,00 HH. & Hipótesis p1=25,00\% por resolución temprana de la dependencia: Agenda EDI y certificación de cada cadena, con pruebas de intercambio y registro de pendientes. Revisar al verificar evidencia y en comité quincenal; si falla la hipótesis, conservar la valoración inicial y activar 8.F.1. \\
R-002 & 60,00\% / 4 / 0,70 & 5 / 0,80; 0,56 (Crítica) & L=170,67 HH; p1L=102,40 HH. & La contención no reconstruye historia: p1=p0 e I1 conserva la consecuencia sanitaria grave. Inventario por lote y decisión de Calidad son necesarios, no prueba de recuperación. Aplicar el criterio R-002 posterior a esta tabla; bloquear lotes sin trazabilidad exigida. \\
R-003 & 30,00\% / 2 / 0,30 & 3 / 0,20; 0,06 (Baja) & L=112,00 HH; p1L=33,60 HH. & Hipótesis p1=30,00\% por resolución temprana de la dependencia: Extracto de tres años perfilado y reconciliado con Abastecimiento/TI; acta V-09. Revisar al verificar evidencia y en comité quincenal; si falla la hipótesis, conservar la valoración inicial y activar 8.F.1. \\
R-004 & 30,00\% / 2 / 0,30 & 4 / 0,40; 0,12 (Media) & L=90,67 HH; p1L=27,20 HH. & Hipótesis p1=30,00\% por resolución temprana de la dependencia: Acta de conteo, causas de diferencias y ajustes autorizados por Bodega/Finanzas. Revisar al verificar evidencia y en comité quincenal; si falla la hipótesis, conservar la valoración inicial y activar 8.F.1. \\
R-005 & 30,00\% / 2 / 0,30 & 4 / 0,40; 0,12 (Media) & L=80,00 HH; p1L=24,00 HH. & Hipótesis p1=30,00\% por resolución temprana de la dependencia: Saldos de envases aceptados por cliente y transportista, política V-08 y evidencia de M8. Revisar al verificar evidencia y en comité quincenal; si falla la hipótesis, conservar la valoración inicial y activar 8.F.1. \\
R-006 & 35,00\% / 3 / 0,50 & 4 / 0,40; 0,20 (Alta) & L=85,33 HH; p1L=29,87 HH. & La búsqueda temprana y V-10 sustentan p1 como hipótesis; no garantizan que exista historia recuperable. I1 conserva la amenaza a aceptación. Aplicar el criterio R-006 posterior a esta tabla; sin referencia histórica válida, no acreditar la comparación. \\
R-007 & 25,00\% / 2 / 0,30 & 3 / 0,20; 0,06 (Baja) & L=65,33 HH; p1L=16,33 HH. & Hipótesis p1=25,00\% por resolución temprana de la dependencia: Universo de deuda vigente conciliado con Finanzas/TI y extracción complementaria comprobada. Revisar al verificar evidencia y en comité quincenal; si falla la hipótesis, conservar la valoración inicial y activar 8.F.1. \\
R-008 & 10,00\% / 1 / 0,10 & 3 / 0,20; 0,02 (Baja) & L=26,00 HH; p1L=2,60 HH. & No se reduce p0: Cola tributaria con correlación, consulta y reintento idempotente; confirmación externa conservada. Revisar al verificar evidencia y en comité quincenal; si falla la hipótesis, conservar la valoración inicial y activar 8.F.1. \\
R-009 & 20,00\% / 2 / 0,30 & 4 / 0,40; 0,12 (Media) & L=65,33 HH; p1L=13,07 HH. & Hipótesis p1=20,00\% por resolución temprana de la dependencia: Matriz térmica V-11 versionada y firmada por Calidad por familia antes de certificación. Revisar al verificar evidencia y en comité quincenal; si falla la hipótesis, conservar la valoración inicial y activar 8.F.1. \\
R-010 & 20,00\% / 2 / 0,30 & 4 / 0,40; 0,12 (Media) & L=56,00 HH; p1L=11,20 HH. & Hipótesis p1=20,00\% por resolución temprana de la dependencia: Recepción técnica por equipo y sitio, con conformidad de TI/Operaciones antes de la ola. Revisar al verificar evidencia y en comité quincenal; si falla la hipótesis, conservar la valoración inicial y activar 8.F.1. \\
R-011 & 25,00\% / 2 / 0,30 & 4 / 0,40; 0,12 (Media) & L=35,33 HH; p1L=8,83 HH. & Hipótesis p1=25,00\% por resolución temprana de la dependencia: Acuerdo operativo firmado y nómina de cada transportista antes de incorporar la ola. Revisar al verificar evidencia y en comité quincenal; si falla la hipótesis, conservar la valoración inicial y activar 8.F.1. \\
R-012 & 15,00\% / 2 / 0,30 & 4 / 0,40; 0,12 (Media) & L=47,33 HH; p1L=7,10 HH. & Hipótesis p1=15,00\% por resolución temprana de la dependencia: Autorización escrita por vehículo y visita; propietario y Operaciones confirman disponibilidad. Revisar al verificar evidencia y en comité quincenal; si falla la hipótesis, conservar la valoración inicial y activar 8.F.1. \\
R-013 & 30,00\% / 2 / 0,30 & 4 / 0,40; 0,12 (Media) & L=108,00 HH; p1L=32,40 HH. & Hipótesis p1=30,00\% por resolución temprana de la dependencia: Transacciones de muestra registradas en el flujo nuevo y cobertura certificada por turno/ola. Revisar al verificar evidencia y en comité quincenal; si falla la hipótesis, conservar la valoración inicial y activar 8.F.1. \\
R-014 & 20,00\% / 2 / 0,30 & 4 / 0,40; 0,12 (Media) & L=86,67 HH; p1L=17,33 HH. & No se reduce p0: V-14, reglas y excepciones transferidas al ayudante y prueba en 13--15 sin planificador. Revisar al verificar evidencia y en comité quincenal; si falla la hipótesis, conservar la valoración inicial y activar 8.F.1. \\
R-015 & 20,00\% / 2 / 0,30 & 4 / 0,40; 0,12 (Media) & L=73,33 HH; p1L=14,67 HH. & Hipótesis p1=20,00\% por resolución temprana de la dependencia: Lista de convocados, asistencia, evaluación y suplentes certificados por sesión. Revisar al verificar evidencia y en comité quincenal; si falla la hipótesis, conservar la valoración inicial y activar 8.F.1. \\
R-016 & 25,00\% / 2 / 0,30 & 4 / 0,40; 0,12 (Media) & L=73,33 HH; p1L=18,33 HH. & Hipótesis p1=25,00\% por resolución temprana de la dependencia: Matriz de decisión con titular y suplente; recepción de insumos y resolución trazables. Revisar al verificar evidencia y en comité quincenal; si falla la hipótesis, conservar la valoración inicial y activar 8.F.1. \\
R-017 & 20,00\% / 2 / 0,30 & 4 / 0,40; 0,12 (Media) & L=48,00 HH; p1L=9,60 HH. & Hipótesis p1=20,00\% por resolución temprana de la dependencia: Recepción de respaldos por sitio y ensayo de pérdida del último camino, incluida Concepción y Los Ángeles. Revisar al verificar evidencia y en comité quincenal; si falla la hipótesis, conservar la valoración inicial y activar 8.F.1. \\
R-018 & 5,00\% / 1 / 0,10 & 4 / 0,40; 0,04 (Baja) & L=20,67 HH; p1L=1,03 HH. & No se reduce p0: Autonomía y transferencia verificadas por componente; coordinación de energía con TI/Operaciones. Revisar al verificar evidencia y en comité quincenal; si falla la hipótesis, conservar la valoración inicial y activar 8.F.1. \\
R-019 & 20,00\% / 2 / 0,30 & 4 / 0,40; 0,12 (Media) & L=57,33 HH; p1L=11,47 HH. & Hipótesis p1=20,00\% por resolución temprana de la dependencia: Aprobación escrita de residencia/transferencia para regiones propuestas antes de habilitar datos. Revisar al verificar evidencia y en comité quincenal; si falla la hipótesis, conservar la valoración inicial y activar 8.F.1. \\
R-020 & 10,00\% / 1 / 0,10 & 4 / 0,40; 0,04 (Baja) & L=30,67 HH; p1L=3,07 HH. & Hipótesis p1=10,00\% por resolución temprana de la dependencia: Flujo de cambios urgentes y suplentes aprobados; autorización escrita por intervención. Revisar al verificar evidencia y en comité quincenal; si falla la hipótesis, conservar la valoración inicial y activar 8.F.1. \\
R-021 & 15,00\% / 2 / 0,30 & 2 / 0,10; 0,03 (Baja) & L=13,33 HH; p1L=2,00 HH. & No se reduce p0: Ventana y recorrido de prueba confirmados por Operaciones/transportista antes de salir. Revisar al verificar evidencia y en comité quincenal; si falla la hipótesis, conservar la valoración inicial y activar 8.F.1. \\
R-022 & 10,00\% / 1 / 0,10 & 2 / 0,10; 0,01 (Baja) & L=13,33 HH; p1L=1,33 HH. & No se reduce p0: Autorización de recorrido y protocolo de seguridad del CLIENTE/transportista. Revisar al verificar evidencia y en comité quincenal; si falla la hipótesis, conservar la valoración inicial y activar 8.F.1. \\
R-023 & 10,00\% / 1 / 0,10 & 2 / 0,10; 0,01 (Baja) & L=17,33 HH; p1L=1,73 HH. & No se reduce p0: Custodia y seguridad del transporte corresponden al CLIENTE/transportista; evidencia de prueba preservada. Revisar al verificar evidencia y en comité quincenal; si falla la hipótesis, conservar la valoración inicial y activar 8.F.1. \\
R-024 & 30,00\% / 2 / 0,30 & 3 / 0,20; 0,06 (Baja) & L=27,00 HH; p1L=8,10 HH. & Hipótesis p1=30,00\% por resolución temprana de la dependencia: Identidad, autorización, capacitación y vínculo persona--vehículo--viaje verificados en andén. Revisar al verificar evidencia y en comité quincenal; si falla la hipótesis, conservar la valoración inicial y activar 8.F.1. \\
R-025 & 30,00\% / 2 / 0,30 & 3 / 0,20; 0,06 (Baja) & L=73,33 HH; p1L=22,00 HH. & Hipótesis p1=30,00\% por resolución temprana de la dependencia: Solicitud y análisis de impacto aprobados antes de iniciar trabajo fuera de alcance. Revisar al verificar evidencia y en comité quincenal; si falla la hipótesis, conservar la valoración inicial y activar 8.F.1. \\
R-026 & 20,00\% / 2 / 0,30 & 4 / 0,40; 0,12 (Media) & L=69,33 HH; p1L=13,87 HH. & No se reduce p0: Seguimiento de cambios aplicables con asesoría competente; decisión contractual documentada. Revisar al verificar evidencia y en comité quincenal; si falla la hipótesis, conservar la valoración inicial y activar 8.F.1. \\
R-027 & 25,00\% / 2 / 0,30 & 3 / 0,20; 0,06 (Baja) & L=49,33 HH; p1L=12,33 HH. & Hipótesis p1=25,00\% por resolución temprana de la dependencia: Reglas de costo de servir aprobadas al mes 20 y tres liquidaciones en sombra reproducidas en 21--23. Revisar al verificar evidencia y en comité quincenal; si falla la hipótesis, conservar la valoración inicial y activar 8.F.1. \\
R-028 & 25,00\% / 2 / 0,30 & 4 / 0,40; 0,12 (Media) & L=56,00 HH; p1L=14,00 HH. & Hipótesis p1=25,00\% por resolución temprana de la dependencia: Acuerdo 5.4.2 con finalidades, accesos y controles comprendidos por Operaciones/TI/sindicato. Revisar al verificar evidencia y en comité quincenal; si falla la hipótesis, conservar la valoración inicial y activar 8.F.1. \\
R-029 & 35,00\% / 3 / 0,50 & 4 / 0,40; 0,20 (Alta) & L=176,00 HH; p1L=61,60 HH. & Tutoría y cobertura certificada sustentan p1 como hipótesis; la rotación recurrente puede superar la capacidad. I1 conserva la amenaza a calidad y certificación. Aplicar el criterio R-029 posterior a esta tabla; no habilitar autonomía sin evaluación. \\
R-030 & 35,00\% / 3 / 0,50 & 4 / 0,40; 0,20 (Alta) & L=116,00 HH; p1L=40,60 HH. & Acceso y ensayo temprano sustentan p1 como hipótesis; el legado no documentado mantiene incertidumbre de integración. I1 conserva la amenaza a aceptación. Aplicar el criterio R-030 posterior a esta tabla; un archivo disponible no acredita equivalencia. \\
R-031 & 35,00\% / 3 / 0,50 & 4 / 0,40; 0,20 (Alta) & L=57,33 HH; p1L=20,07 HH. & Habilitación anticipada y continuidad sustentan p1 como hipótesis; peak y congelamiento limitan la recuperación. I1 conserva la amenaza al hito. Aplicar el criterio R-031 posterior a esta tabla; una condición incumplida impide declarar cierre. \\
R-032 & 25,00\% / 2 / 0,30 & 4 / 0,40; 0,12 (Media) & L=86,67 HH; p1L=21,67 HH. & Hipótesis p1=25,00\% por resolución temprana de la dependencia: Inventario de versiones y soporte, revisión trimestral y compatibilidad probada antes del límite verificado. Revisar al verificar evidencia y en comité quincenal; si falla la hipótesis, conservar la valoración inicial y activar 8.F.1. \\
R-033 & 20,00\% / 2 / 0,30 & 4 / 0,40; 0,12 (Media) & L=109,33 HH; p1L=21,87 HH. & Hipótesis p1=20,00\% por resolución temprana de la dependencia: Cuentas, claves, exportación y restauración utilizable verificadas en recepción y revisión anual de salida. Revisar al verificar evidencia y en comité quincenal; si falla la hipótesis, conservar la valoración inicial y activar 8.F.1. \\
R-034 & 25,00\% / 2 / 0,30 & 4 / 0,40; 0,12 (Media) & L=89,33 HH; p1L=22,33 HH. & Hipótesis p1=25,00\% por resolución temprana de la dependencia: Procedimiento ERP con TI/Finanzas probado en AL-DTE-01/AL-DR-01 antes de H5/H10. Revisar al verificar evidencia y en comité quincenal; si falla la hipótesis, conservar la valoración inicial y activar 8.F.1. \\
R-035 & 25,00\% / 2 / 0,30 & 4 / 0,40; 0,12 (Media) & L=57,33 HH; p1L=14,33 HH. & Hipótesis p1=25,00\% por resolución temprana de la dependencia: Demostración práctica en andén e incorporación acordada con transportista y conductor habilitado. Revisar al verificar evidencia y en comité quincenal; si falla la hipótesis, conservar la valoración inicial y activar 8.F.1. \\
R-036 & 30,00\% / 2 / 0,30 & 4 / 0,40; 0,12 (Media) & L=38,67 HH; p1L=11,60 HH. & Sin reducción de p0: Perfil versionado y validado por Operaciones/TI, comparación de carga acordada y observada por ventana y operación, informe CT-14 y decisión de habilitación del alcance verificado. Conservar I=4 y E=0,12; la prueba permite habilitar sólo el alcance verificado. Si falta evidencia o falla CT-14, mantener bloqueo y corregir/repetir; no acreditar eficacia ni alcance ampliado por el solo resultado de carga. \\
\end{tablalafrox}
\end{landscape}
\endgroup

\subsection*{Verificación de los cinco riesgos prioritarios (8.F.2)}
\label{sec:sd804-verificacion}

La Tabla 8.F.2 conserva la proyección del tratamiento completo; sus valores son
objetivos condicionados de planificación, no una valoración residual acreditada.
Para decidir mientras falta evidencia se utiliza la valoración inicial: R-002,
0,56 Crítica; R-006 y R-029, 0,28 Alta; R-030 y R-031, 0,20 Alta.
La proyección es respectivamente 0,56 Crítica y 0,20 Alta para los otros cuatro.
Ninguna de esas proyecciones demuestra un nivel manejable.

Cada ficha de 8.F.1 se completa con la siguiente matriz de aprobación. El dueño
prepara el expediente; CAL verifica su integridad y reproducibilidad; las áreas
del CLIENTE validan las decisiones de su competencia y JP registra el resultado.
Cada expediente identifica versión, fecha, universo o período cubierto, fuentes,
entradas, resultado esperado y obtenido, excepciones, firmas y próxima revisión.
Las pruebas aquí descritas están propuestas, pendientes de ejecución.

\begingroup
\renewcommand{\thetable}{8.F.2a}
\begin{landscape}
\setlength{\textwidth}{\linewidth}
\renewcommand{\arraystretch}{1.08}
\begin{tablalafrox}{Matriz de verificación y habilitación de SD8-04. Fuente: fichas 8.F.1 y condiciones de aceptación de S8.}{tab:sd804-verificacion}{G{0.06} F{0.24} F{0.17} F{0.18} Y}{\cab{ID} & \cab{Prueba y evidencia requerida} & \cab{Dueño y verificación} & \cab{Disparador y plazo} & \cab{Aprobación y respuesta ante fallo}}
R-002 & Conciliar el inventario completo de vínculos sanitarios exigidos por lote contra ERP/WMS y fuentes de origen. Registrar cada excepción V-07. Ensayar exclusión de un lote sin vínculo en carga y despacho, conservando resultado y registro de bloqueo. & DAT prepara; CAL verifica conciliación y ensayo; Jefa de Calidad del CLIENTE valida V-07 y disposición de excepciones. & Faltante al perfilar: aislar el mismo día. Política al cierre de mes 2; excepciones a D-5 de cada validación; comprobar bloqueo antes del corte. & Habilitar sólo lotes con todos los vínculos exigidos verificables y prueba de bloqueo satisfactoria. Una firma no reemplaza la historia. Ante excepción irrecuperable, mantener lote bloqueado y elevar por CT el efecto contractual antes del corte; continuar sólo alcance independiente. \\
R-006 & Versionar fuentes, período, horas de entrega y montos esperados; reproducir OTIF y cuadratura con las reglas de aceptación vigentes. Conservar entradas, cálculo, conciliación y acta V-10 con faltantes. & IMP prepara; CAL reproduce y contrasta; Comercial/Operaciones validan horas y Finanzas montos y referencia. & V-10 al cierre del mes 1, antes de comparar aceptación. Ante bloqueo, JP reúne áreas en dos días hábiles y exige evidencia o eleva la brecha en cinco. & Aprobar sólo una referencia reproducible, respaldada por fuentes y validada por las áreas responsables. Ante faltantes o divergencias sin explicar, bloquear comparación y aceptación afectadas; corregir fuentes y repetir. La captura futura no sustituye historia obligatoria. \\
R-029 & Conciliar altas por turno con tutores, tiempo disponible y cobertura certificada. Registrar inicio/fin de aprendizaje, evaluación individual y autorización previa a autonomía; contrastar diariamente demanda de inducción y recuperación con capacidad efectiva. & IMP hasta mes 21; SRE desde 22. CAL verifica registros; Bodega y Operaciones confirman tutores y cobertura por turno. & Meses 13--56: antes de cada turno y de cada alta autónoma; revisar diariamente. Ante sobrecarga, cobertura certificada antes del turno y propuesta de refuerzo en dos días hábiles. & Cada alta debe tener tutor, aprendizaje máximo de dos horas y evaluación aprobada antes de autonomía; demanda cubierta por capacidad comprobada. Si falla, impedir autonomía y mantener certificación afectada pendiente; repetir formación/evaluación y verificar refuerzo antes de habilitar. \\
R-030 & Ensayo reproducible con TI sobre ERP/WMS real y accesos autorizados: conservar contratos, entradas/salidas, marcas de tiempo y conciliación contra origen. Probar integridad, latencia requerida y emisión tributaria por ERP sin modificarlo. Aplicar los mismos requisitos al archivo alternativo. & ARQ dirige, DAT concilia extracción/carga; CAL verifica resultados; TI del CLIENTE confirma acceso y comportamiento del legado. & Brecha de acceso o contrato: elevar al finalizar semana 2 del mes 1; verificar contrato inicial al cierre de mes 1. Ensayo aprobado antes de integración y tras corregir cualquier fallo. & Aprobar sólo intercambios y extracción conciliados, sin diferencias inexplicadas y con requisitos aplicables demostrados. Ante fallo o acceso insuficiente, bloquear integración y certificación dependientes, corregir y repetir ensayo; JP eleva alternativas por CT antes del hito. \\
R-031 & Conservar lista de habilitaciones de agosto, ensayo de continuidad con resultados contrastados con compromisos vigentes y registro diario de las seis condiciones de 8.3.3 en meses 19--20, incluidas las cuatro semanas finales y acta de aceptación. & JP coordina; CAL verifica expediente completo; Operaciones valida continuidad y disponibilidad de usuarios; CLIENTE formaliza aceptación. & Habilitaciones y ensayo antes del congelamiento de 1--25 de septiembre. Revisar cada día en 19--20; ante condición incumplida, escalar el mismo día y acordar recuperación en dos días hábiles. & Cierre sólo con las seis condiciones acumulativas y período completo de evidencia. Ante prueba fallida o jornada sin respaldo, mantener cierre pendiente y activar continuidad aprobada; corregir y repetir lo afectado en ventana autorizada, con R-020 cuando corresponda y sin mover unilateralmente hitos. \\
\end{tablalafrox}
\end{landscape}
\endgroup

\paragraph{Resultado de verificación y reevaluación}
Se registran tres resultados: \textbf{satisfactorio}, cuando la evidencia cubre
todo el alcance y período exigidos y cumple el criterio; \textbf{insuficiente},
cuando faltan fuentes, cobertura, firmas o reproducibilidad; y \textbf{fallido},
cuando la prueba contradice el resultado requerido. Sólo el satisfactorio permite
proponer el levantamiento del bloqueo concreto y reevaluar el riesgo; requiere
además la aprobación competente. Los otros dos mantienen el bloqueo y la referencia
inicial, abren corrección con dueño y fecha y exigen repetir la verificación.
Una prueba fallida se registra también como problema cuando el evento ya ocurrió.

Para reevaluar, el dueño y CAL documentan p, nivel P y factor según las bandas
de 8.1.3 en la misma ventana de exposición; revisan esfuerzo, plazo, alcance,
calidad y contrato, toman el mayor impacto y recalculan E. JP registra fundamento,
fuentes y aprobación. El objetivo de seguimiento es una zona inferior a Alta
($E<0,20$), sin rebajar requisitos ni consecuencias graves. Con I=4 y factor 0,40,
ello exigiría justificar P no mayor que 2 ($p\leq30\%$, E=0,12); con I=5 y factor
0,80, P=1 ($p\leq10\%$, E=0,08). Estos límites describen la escala, no nuevas
probabilidades estimadas ni metas de eficacia ya alcanzables. R-002 conserva
p=60\%, I=5 y zona Crítica: aislar no restaura historia. Los otros cuatro conservan
sus proyecciones Altas hasta que evidencia de la ventana sustente otra valoración.
La habilitación puntual y un NPR menor no acreditan por sí solos ese objetivo.

\paragraph{Suficiencia del esfuerzo y condición de programación}
Los presupuestos A/C de las cinco fichas son, respectivamente, R-002: 12/24 HH;
R-006: 8/16 HH; R-029: 12/32 HH; R-030: 12/24 HH y R-031: 8/16 HH.
Cubren la hipótesis de coordinación y respuesta definida al inicio de 8.F;
su suficiencia para conciliar universos completos, repetir ensayos, cubrir altas
reales y sostener evidencia diaria queda pendiente de contraste con volumen,
duración y capacidad disponibles. Antes de la primera verificación, cada dueño
estima por actividad y perfil el trabajo de preparación, ejecución, corrección y
repetición; CAL revisa y JP concilia con los paquetes ordinarios y 8.H.
Las HH ordinarias ya asignadas se identifican y no se imputan de nuevo. Todo
faltante se registra como estimación adicional pendiente de conciliación y
autorización, con efecto en capacidad y plazo; no se considera absorbido por A/C
ni por una reserva sin asignación. Si falta capacidad, JP escala antes del hito y
mantiene la habilitación afectada pendiente. Esta revisión documental conserva
los importes HH y cálculos de 8.F.3 y 8.H hasta disponer de una estimación aprobada.

\subsection*{Justificación y decisión de los residuales altos y crítico}
\label{sec:residuales-prioritarios}

\paragraph{Lectura antes/después y estado de la evidencia}
Antes del tratamiento, R-002 es Crítica y los otros cuatro son Altos en T-16.1.
Después del tratamiento completo proyectado, R-002 permanece Crítica y los demás
permanecen Altos en la Tabla~\ref{tab:cierre-residual}; no se declaran aceptables.
La reducción propuesta de ocurrencia no reduce I1: detectar o bloquear antes
limita exposición operativa, pero no elimina la consecuencia si el evento ocurre.
Los criterios siguientes son evidencia por obtener, no resultados operacionales.
CAL contrasta fuentes, pruebas y pendientes; JP registra fecha, alcance afectado,
resultado, responsable y próxima revisión. Sin prueba satisfactoria no se acredita
la hipótesis p1 y se conserva la valoración inicial para decidir.

\paragraph{R-002: historia sanitaria no reconstruible}
El plan adicional contiene los lotes afectados; no disminuye la probabilidad de
que la historia recibida sea irrecuperable ni la gravedad sanitaria de esa pérdida.
DAT debe contrastar el inventario de vínculos por lote con evidencia de origen;
Calidad del CLIENTE resuelve cada excepción conforme a V-07. La eficacia de la
contención exige comprobar que el conjunto aislado no entra en la carga ni en
los despachos; una resolución firmada sin vínculos exigidos no habilita el lote.
Se mantienen aislamiento el mismo día y decisión antes del corte según 8.F.1.
Si la historia no es reconstruible, JP eleva por CT al Comité Ejecutivo la
imposibilidad y sus efectos de alcance contractual; el lote sigue bloqueado.
Sólo continúa trabajo independiente con datos verificables. Ni esa decisión ni
una mejora de las HH de atención restituyen trazabilidad perdida.

\paragraph{R-006: referencia histórica de aceptación}
La coordinación temprana puede recuperar fuentes existentes, pero no crear horas
de entrega ni montos esperados que nunca se registraron; por ello subsiste la
amenaza a las comparaciones de OTIF y cuadratura. IMP y CAL deben verificar que
las fuentes de Comercial, Operaciones y Finanzas permiten reproducir esas
comparaciones para el período de referencia y documentar V-10 y sus faltantes.
Reunir áreas o empezar a medir hacia adelante no demuestra esa eficacia.
JP mantiene la reunión dentro de dos días hábiles y la exigencia de evidencia o
elevación contractual dentro de cinco según 8.F.1. Sin referencia válida quedan
bloqueadas las comparaciones y su aceptación, no las pruebas independientes.
Si persiste la brecha, el Comité Ejecutivo recibe por CT alternativas e impacto
contractual documentados, sin sustituir la historia exigida ni aceptar el resultado.

\paragraph{R-029: capacidad sostenida ante rotación nocturna}
La tutoría reduce la exposición sólo si cubre las altas reales del turno; una
sesión satisfactoria no acredita capacidad para toda la ventana de operación.
Bodega y Operaciones deben aportar altas y cobertura certificada por turno;
IMP/SRE y CAL contrastan tutor disponible, duración de aprendizaje máxima de dos
horas y evaluación aprobada antes de cada asignación autónoma, según 7.1.2.
La revisión diaria debe mostrar que altas y recuperaciones pendientes caben en
la capacidad de tutoría y evaluación disponible, incluida la cobertura de sus
tareas por personal certificado. Si no caben o falla la evaluación, se bloquea
la autonomía de esa persona y queda pendiente la certificación afectada.
Se mantiene cobertura certificada antes del turno y propuesta de refuerzo en dos
días hábiles; el refuerzo se revalida antes de habilitar. Si persiste la sobrecarga,
JP eleva por CT capacidad y recuperación al Comité Ejecutivo; no habilita por
urgencia ni confunde el refuerzo propuesto con dotación ya disponible.

\paragraph{R-030: verificación efectiva del ERP/WMS}
Acceso, diccionario y contrato de intercambio reducen incertidumbre, pero la
falta de documentación puede impedir demostrar la integración a tiempo.
ARQ y DAT deben ejecutar con TI un ensayo sobre el ERP/WMS real, sin modificar
el ERP, y CAL conservar entradas, salidas, marcas de tiempo y conciliación de
extracción/carga contra origen. Debe probar latencia e integridad exigidas,
sustitución progresiva del WMS y emisión tributaria por el ERP, no sólo lectura
de archivos. Un intercambio controlado se admite únicamente si demuestra los
mismos requisitos; no se presume su equivalencia. ARQ eleva la brecha al finalizar
semana 2 del mes 1; sin acceso o alternativa verificada antes de integración,
se bloquean las pruebas dependientes y la certificación afectada. Continúan
pruebas independientes, sin presentarlas como validación del legado. Si persiste
la brecha, JP eleva por CT alternativas y efectos contractuales al Comité
Ejecutivo, sin exceptuar integridad, latencia ni obligaciones tributarias.

\paragraph{R-031: cierre bajo peak y congelamiento}
Cerrar habilitaciones en agosto y acordar continuidad antes de septiembre no
suprimen la coincidencia del peak con el cierre ni crean una ventana de cambio.
JP y CAL deben contrastar el registro diario de meses 19--20 con las seis
condiciones acumulativas de 8.3.3: ausencia de incidentes críticos/altos atribuibles
abiertos, volumen real comprometido en las cuatro últimas semanas, disponibilidad
y respuesta comprometidas, conciliación sin diferencias inexplicadas, personal
capacitado y certificado y acta de aceptación. Una condición o una jornada aislada
no reemplaza el conjunto ni las cuatro semanas de evidencia.
Durante 1--25 de septiembre se conserva el congelamiento y se activa continuidad
aprobada; toda intervención requiere la autorización de R-020. Ante incumplimiento,
JP escala el mismo día y acuerda recuperación en dos días hábiles según 8.F.1,
sin declarar cierre ni extender unilateralmente el calendario. Se aplica 8.3.3
si debe extenderse la marcha blanca: a costo del adjudicatario y sin desplazar
fases posteriores. Sólo continúa trabajo independiente y operación autorizada.

\paragraph{Tratamiento mientras persista la prioridad}
Los cinco riesgos permanecen abiertos con revisión diaria de evidencia y bloqueos.
JP eleva al Comité Ejecutivo por CT las brechas, alternativas, capacidad y efecto
sobre aceptación antes del hito; la escalada no es una solución ni una aceptación
del residual. La autoridad competente decide dentro de sus atribuciones y deja
acta, sin dispensar requisitos sanitarios, tributarios, de servicio o contractuales.
Levantar un bloqueo requiere evidencia validada del alcance concreto y las actas
exigidas; no rebaja automáticamente la categoría de toda la ventana del riesgo.
Cualquier reevaluación posterior necesita fundamento verificable, no una reducción
numérica atribuida al solo hecho de aprobar el plan adicional.

\begingroup
\renewcommand{\thetable}{8.F.3}
\begin{landscape}
\setlength{\textwidth}{\linewidth}
\renewcommand{\arraystretch}{1.08}
\begin{tablalafrox}{Costo-beneficio y sensibilidad de las 15 respuestas altas/críticas iniciales (HH). Fuente: modelo explícito de 8.F.}{tab:cierre-cb}{G{0.06} F{0.22} F{0.22} F{0.25} Y}{\cab{ID} & \cab{Entradas acumuladas} & \cab{Resultado esperado} & \cab{Sensibilidad / umbral} & \cab{Decisión y justificación}}
R-001 & A=8,00; C=16,00; H0=104,00; p0=45,00\%; p1=25,00\%. & B=24,80; T=12,00; B/T=2,07; neto=12,80 HH. & Sin mejora p1=p0: B/T=0,47; neto=-8,00 HH. Equilibrio neto: p1<=37,31\%. & Eficiencia de esfuerzo favorable bajo la hipótesis p1; mantener el tratamiento por obligación y no usar B/T como autorización de aceptación. \\
R-002 & A=12,00; C=24,00; H0=194,67; p0=60,00\%; p1=60,00\%. & B=14,40; T=26,40; B/T=0,55; neto=-12,00 HH. & Sin mejora p1=p0: B/T=0,55; neto=-12,00 HH. Equilibrio neto: p1<=53,84\%. & Prioridad obligatoria: proteger trazabilidad, condiciones de aceptación y continuidad; se mantiene aun con retorno HH menor que uno. Validar capacidad y elevar su esfuerzo antes del hito. \\
R-004 & A=8,00; C=20,00; H0=110,67; p0=55,00\%; p1=30,00\%. & B=33,67; T=14,00; B/T=2,40; neto=19,67 HH. & Sin mejora p1=p0: B/T=0,58; neto=-8,00 HH. Equilibrio neto: p1<=47,77\%. & Eficiencia de esfuerzo favorable bajo la hipótesis p1; mantener el tratamiento por obligación y no usar B/T como autorización de aceptación. \\
R-005 & A=6,00; C=16,00; H0=96,00; p0=50,00\%; p1=30,00\%. & B=24,00; T=10,80; B/T=2,22; neto=13,20 HH. & Sin mejora p1=p0: B/T=0,57; neto=-6,00 HH. Equilibrio neto: p1<=43,75\%. & Eficiencia de esfuerzo favorable bajo la hipótesis p1; mantener el tratamiento por obligación y no usar B/T como autorización de aceptación. \\
R-006 & A=8,00; C=16,00; H0=101,33; p0=55,00\%; p1=35,00\%. & B=25,87; T=13,60; B/T=1,90; neto=12,27 HH. & Sin mejora p1=p0: B/T=0,52; neto=-8,00 HH. Equilibrio neto: p1<=47,11\%. & Eficiencia de esfuerzo favorable bajo la hipótesis p1; mantener el tratamiento por obligación y no usar B/T como autorización de aceptación. \\
R-009 & A=6,00; C=12,00; H0=77,33; p0=35,00\%; p1=20,00\%. & B=14,00; T=8,40; B/T=1,67; neto=5,60 HH. & Sin mejora p1=p0: B/T=0,41; neto=-6,00 HH. Equilibrio neto: p1<=27,24\%. & Eficiencia de esfuerzo favorable bajo la hipótesis p1; mantener el tratamiento por obligación y no usar B/T como autorización de aceptación. \\
R-011 & A=4,00; C=8,00; H0=43,33; p0=35,00\%; p1=25,00\%. & B=6,33; T=6,00; B/T=1,06; neto=0,33 HH. & Sin mejora p1=p0: B/T=0,41; neto=-4,00 HH. Equilibrio neto: p1<=25,77\%. & Eficiencia de esfuerzo favorable bajo la hipótesis p1; mantener el tratamiento por obligación y no usar B/T como autorización de aceptación. \\
R-013 & A=8,00; C=20,00; H0=128,00; p0=55,00\%; p1=30,00\%. & B=38,00; T=14,00; B/T=2,71; neto=24,00 HH. & Sin mejora p1=p0: B/T=0,58; neto=-8,00 HH. Equilibrio neto: p1<=48,75\%. & Eficiencia de esfuerzo favorable bajo la hipótesis p1; mantener el tratamiento por obligación y no usar B/T como autorización de aceptación. \\
R-016 & A=4,00; C=12,00; H0=85,33; p0=45,00\%; p1=25,00\%. & B=20,07; T=7,00; B/T=2,87; neto=13,07 HH. & Sin mejora p1=p0: B/T=0,57; neto=-4,00 HH. Equilibrio neto: p1<=40,31\%. & Eficiencia de esfuerzo favorable bajo la hipótesis p1; mantener el tratamiento por obligación y no usar B/T como autorización de aceptación. \\
R-028 & A=6,00; C=12,00; H0=68,00; p0=45,00\%; p1=25,00\%. & B=16,60; T=9,00; B/T=1,84; neto=7,60 HH. & Sin mejora p1=p0: B/T=0,47; neto=-6,00 HH. Equilibrio neto: p1<=36,18\%. & Eficiencia de esfuerzo favorable bajo la hipótesis p1; mantener el tratamiento por obligación y no usar B/T como autorización de aceptación. \\
R-029 & A=12,00; C=32,00; H0=208,00; p0=55,00\%; p1=35,00\%. & B=52,80; T=23,20; B/T=2,28; neto=29,60 HH. & Sin mejora p1=p0: B/T=0,59; neto=-12,00 HH. Equilibrio neto: p1<=49,23\%. & Eficiencia de esfuerzo favorable bajo la hipótesis p1; mantener el tratamiento por obligación y no usar B/T como autorización de aceptación. \\
R-030 & A=12,00; C=24,00; H0=140,00; p0=50,00\%; p1=35,00\%. & B=29,40; T=20,40; B/T=1,44; neto=9,00 HH. & Sin mejora p1=p0: B/T=0,50; neto=-12,00 HH. Equilibrio neto: p1<=41,43\%. & Eficiencia de esfuerzo favorable bajo la hipótesis p1; mantener el tratamiento por obligación y no usar B/T como autorización de aceptación. \\
R-031 & A=8,00; C=16,00; H0=73,33; p0=45,00\%; p1=35,00\%. & B=12,93; T=13,60; B/T=0,95; neto=-0,67 HH. & Sin mejora p1=p0: B/T=0,47; neto=-8,00 HH. Equilibrio neto: p1<=34,09\%. & Prioridad obligatoria: proteger trazabilidad, condiciones de aceptación y continuidad; se mantiene aun con retorno HH menor que uno. Validar capacidad y elevar su esfuerzo antes del hito. \\
R-032 & A=6,00; C=20,00; H0=106,67; p0=40,00\%; p1=25,00\%. & B=21,00; T=11,00; B/T=1,91; neto=10,00 HH. & Sin mejora p1=p0: B/T=0,57; neto=-6,00 HH. Equilibrio neto: p1<=34,38\%. & Eficiencia de esfuerzo favorable bajo la hipótesis p1; mantener el tratamiento por obligación y no usar B/T como autorización de aceptación. \\
R-035 & A=6,00; C=16,00; H0=73,33; p0=40,00\%; p1=25,00\%. & B=15,00; T=10,00; B/T=1,50; neto=5,00 HH. & Sin mejora p1=p0: B/T=0,52; neto=-6,00 HH. Equilibrio neto: p1<=31,82\%. & Eficiencia de esfuerzo favorable bajo la hipótesis p1; mantener el tratamiento por obligación y no usar B/T como autorización de aceptación. \\
\end{tablalafrox}
\end{landscape}
\endgroup

\clearpage
\section{Anexo 8.G: Análisis de modos de falla y efectos (FMEA)}
FMEA aplicado a los 15 riesgos con E0 mayor o igual a 0,20. El modo de falla es el evento de T-16 y su efecto es la consecuencia allí declarada. Se utilizan cinco niveles definidos para este proyecto: severidad S=I global (1 corrección localizada, 2 ajuste de paquete, 3 reproceso sin afectar hitos, 4 amenaza a aceptación/servicio/hito, 5 consecuencia sanitaria grave o trazabilidad irrecuperable). O usa el nivel 1--5 de la misma banda de p0/p1, dentro de la ventana del evento, no una frecuencia diaria. D es dificultad para detectar antes del efecto: 1 comprobación automática integral bloqueante; 2 control documental o ensayo completo antes de liberar; 3 muestreo/reconciliación con cobertura parcial; 4 evidencia tardía o validación dispersa; 5 dato histórico no reconstruible o defecto oculto hasta el efecto.

NPR=S por O por D prioriza modos dentro de esta escala (1--125); no equivale a E, a probabilidad ni a HH. Se actúa por severidad 5 aunque NPR sea bajo. La propuesta obliga a tratar S=5, D mayor o igual a 4 y cualquier E alta/crítica. CAL conserva expediente de modo, causa, efecto, prueba y corrección. Las puntuaciones posteriores son proyecciones sujetas a la verificación de la evidencia de cada control. No se usa menor NPR para dispensar requisitos o aceptar un residual alto.

Para R-002, R-006, R-029, R-030 y R-031 se aplica la matriz 8.F.2a. Los valores
posteriores S1/O1/D1 y NPR1 son objetivos condicionados, pendientes de evidencia;
para decidir se conservan los iniciales si ésta resulta insuficiente o la prueba
falla. Sólo una verificación satisfactoria permite proponer una reevaluación,
registrando fuente, cobertura, fecha y aprobación. Mejorar D o NPR no reduce
automáticamente S ni E: se mantiene S1=5 para R-002 y S1=4 para los otros cuatro.
Sus zonas proyectadas siguen Crítica/Altas. Las condiciones de habilitación de
8.F.2a y la revisión de suficiencia de HH de 8.F complementan las acciones siguientes.

\begingroup
\renewcommand{\thetable}{8.G.1}
\begin{landscape}
\setlength{\textwidth}{\linewidth}
\renewcommand{\arraystretch}{1.08}
\begin{tablalafrox}{FMEA de 15 riesgos altos/críticos y evaluación posterior proyectada. Fuente: eventos T-16 y juicios de planificación de S8.}{tab:cierre-fmea}{G{0.06} F{0.12} F{0.22} F{0.23} F{0.18} Y}{\cab{ID} & \cab{FMEA inicial} & \cab{Modo, disparador y efecto} & \cab{Detección inicial y control} & \cab{FMEA posterior proyectada} & \cab{Acción, responsable y condición de habilitación}}
R-001 & S=4; O=3; D=4; NPR=48. & Cadena del alcance sin certificación antes del límite de 28 días. Dado que la habilitación EDI de cada cadena depende de su participación y certificación externa antes de las cuatro semanas finales de evidencia de E2, entonces puede que una cadena incluida en el alcance no complete su participación o certificación EDI antes de la fecha requerida, teniendo como consecuencia que no se acredite completo el intercambio electrónico exigido en E2 y deban reprogramarse pruebas y coordinación. & D0=4: la validación distribuida puede llegar después de la fecha necesaria. Control/acción: Agenda EDI y certificación de cada cadena, con pruebas de intercambio y registro de pendientes. & S1=4; O1=2; D1=2; NPR1=16. D1=2: acta o ensayo específico previo a habilitación; se invalida la mejora de detección si no existe evidencia. & Al D-5 sin certificación, JP escala en un día hábil y conserva pendiente la cadena; sólo ampliar tras su certificación. Dueño: ARQ; verifica CAL. \\
R-002 & S=5; O=4; D=5; NPR=100. & Lotes o vínculos históricos sin evidencia de origen al perfilar la carga. Dado que la migración sanitaria utiliza antecedentes del ERP, WMS y planillas y requiere vínculos históricos verificables con evidencia de origen, entonces puede que los datos recibidos contengan faltantes de lote o vínculos históricos que no puedan resolverse con esa evidencia, teniendo como consecuencia que la carga sanitaria afectada quede pendiente de conciliación y decisión documentada, y aumente el trabajo de saneamiento. & D0=5: faltante de origen o interfaz no verificable puede permanecer oculto hasta conciliar/aceptar. Control/acción: Inventario de vínculos sanitarios por lote, evidencia de origen y resolución firmada de Calidad sobre cada excepción. & S1=5; O1=4; D1=3; NPR1=60. D1=3: aun conciliando, la cobertura histórica puede ser parcial; no se reconstruye lo inexistente. & Plan adicional: DAT aísla el conjunto el mismo día; Calidad decide antes del corte. No cargar ni despachar lotes sin trazabilidad exigida; si no es reconstruible, elevar decisión de alcance contractual, sin inventar historia. Dueño: DAT; verifica CAL. \\
R-004 & S=4; O=4; D=4; NPR=64. & Diferencias sin causa ni ajuste aprobado en el acta previa al corte. Dado que la aceptación del inventario migrado requiere conciliar los saldos legados con el conteo físico y aprobar las explicaciones y ajustes antes del corte, entonces puede que se detecten diferencias sin explicación y ajuste aprobados a tiempo, teniendo como consecuencia que deban repetirse el saneamiento y la validación, y la carga afectada no pueda ser aceptada hasta resolverlas. & D0=4: la validación distribuida puede llegar después de la fecha necesaria. Control/acción: Acta de conteo, causas de diferencias y ajustes autorizados por Bodega/Finanzas. & S1=4; O1=2; D1=2; NPR1=16. D1=2: acta o ensayo específico previo a habilitación; se invalida la mejora de detección si no existe evidencia. & Bloquear corte con diferencia no explicada; DAT presenta conciliación y nueva prueba antes de la ventana autorizada. Dueño: DAT; verifica CAL. \\
R-005 & S=4; O=3; D=4; NPR=48. & No existe saldo inicial aceptado por cliente y transportista. Dado que la carga inicial de M8 requiere saldos de envases por cliente y transportista sustentados en antecedentes conciliables del cliente, entonces puede que los antecedentes entregados no permitan determinar un saldo conciliado por cliente y transportista, teniendo como consecuencia que deban repetirse la preparación y validación de la carga y no pueda demostrarse M8 con saldos iniciales aceptados. & D0=4: la validación distribuida puede llegar después de la fecha necesaria. Control/acción: Saldos de envases aceptados por cliente y transportista, política V-08 y evidencia de M8. & S1=4; O1=2; D1=2; NPR1=16. D1=2: acta o ensayo específico previo a habilitación; se invalida la mejora de detección si no existe evidencia. & Mantener saldos afectados pendientes; Finanzas resuelve antes de su carga, sin asumir saldo cero. Dueño: DAT; verifica CAL. \\
R-006 & S=4; O=4; D=5; NPR=80. & V-10 no obtiene horas o montos históricos verificables. Dado que las comparaciones históricas de aceptación dependen de horas verificables de entrega y montos esperados de rendición que la captura futura no reconstruye, entonces puede que falten horas verificables de entrega o montos esperados de rendición en los antecedentes históricos, teniendo como consecuencia que no puedan respaldarse con esos registros las comparaciones de OTIF o cuadratura y deba repetirse la validación del período de referencia. & D0=5: faltante de origen o interfaz no verificable puede permanecer oculto hasta conciliar/aceptar. Control/acción: Fuentes históricas de horas y montos verificadas con Comercial, Operaciones y Finanzas, V-10. & S1=4; O1=3; D1=3; NPR1=36. D1=3: aun conciliando, la cobertura histórica puede ser parcial; no se reconstruye lo inexistente. & Plan adicional: JP reúne áreas dentro de dos días hábiles desde bloqueo; en cinco días hábiles exige evidencia o eleva al Comité Ejecutivo la condición contractual. No sustituir historia por mediciones futuras ni aceptar resultado sin referencia. Dueño: IMP; verifica CAL. \\
R-009 & S=4; O=3; D=4; NPR=48. & Calidad no aprueba parámetros por producto al preparar certificación. Dado que la certificación de la parametrización térmica requiere una matriz de umbrales y duraciones por producto aprobada por Calidad del cliente, entonces puede que Calidad no apruebe esa matriz antes del cierre de las pruebas y la marcha blanca de los productos afectados, teniendo como consecuencia que no pueda completarse la evidencia de parametrización y aceptación del alcance térmico y deban reprogramarse validaciones. & D0=4: la validación distribuida puede llegar después de la fecha necesaria. Control/acción: Matriz térmica V-11 versionada y firmada por Calidad por familia antes de certificación. & S1=4; O1=2; D1=2; NPR1=16. D1=2: acta o ensayo específico previo a habilitación; se invalida la mejora de detección si no existe evidencia. & Sin parámetros aprobados, bloquear certificación del alcance afectado y escalar a D-5; no relajar bloqueo térmico. Dueño: CAL; verifica CAL. \\
R-011 & S=4; O=3; D=4; NPR=48. & Acuerdo operativo de un transportista pendiente antes de su ola. Dado que la incorporación de cada empresa transportista al reparto requiere que su acuerdo operativo esté firmado antes de la ola correspondiente, entonces puede que una empresa transportista no firme el acuerdo a tiempo, teniendo como consecuencia que sus conductores no puedan incorporarse ni completar las pruebas de esa ola y quede pendiente demostrar la cobertura del transportista. & D0=4: la validación distribuida puede llegar después de la fecha necesaria. Control/acción: Acuerdo operativo firmado y nómina de cada transportista antes de incorporar la ola. & S1=4; O1=2; D1=2; NPR1=16. D1=2: acta o ensayo específico previo a habilitación; se invalida la mejora de detección si no existe evidencia. & A D-5 sin firma, escalar por CT el mismo día; no habilitar empresa pendiente ni reducir cobertura obligatoria. Dueño: IMP; verifica CAL. \\
R-013 & S=4; O=4; D=4; NPR=64. & Transacciones se siguen registrando solo por el procedimiento anterior. Dado que la implantación exige que los usuarios incorporen sus transacciones al registro oficial de la solución, entonces puede que parte de los usuarios continúe registrando exclusivamente en el procedimiento anterior durante una ola, teniendo como consecuencia que no pueda acreditarse la cobertura de adopción requerida para cerrar esa ola y deba extenderse la intervención de implantación. & D0=4: la validación distribuida puede llegar después de la fecha necesaria. Control/acción: Transacciones de muestra registradas en el flujo nuevo y cobertura certificada por turno/ola. & S1=4; O1=2; D1=2; NPR1=16. D1=2: acta o ensayo específico previo a habilitación; se invalida la mejora de detección si no existe evidencia. & Si persiste flujo paralelo, IMP abre refuerzo en el siguiente turno y mide diariamente; sin evidencia de adopción no cerrar ola. Dueño: IMP; verifica CAL. \\
R-016 & S=4; O=3; D=4; NPR=48. & Decisión requerida no llega en la fecha acordada y bloquea la tarea. Dado que los diseños, pruebas y aceptaciones incluyen actividades que dependen de validaciones o decisiones de áreas del cliente, entonces puede que el área requerida no entregue una validación o decisión en la fecha acordada, teniendo como consecuencia que se detengan o repitan tareas dependientes, se consuma holgura y deban reprogramarse recursos. & D0=4: la validación distribuida puede llegar después de la fecha necesaria. Control/acción: Matriz de decisión con titular y suplente; recepción de insumos y resolución trazables. & S1=4; O1=2; D1=2; NPR1=16. D1=2: acta o ensayo específico previo a habilitación; se invalida la mejora de detección si no existe evidencia. & D-5 sin decisión: escalar un día hábil después; promoción regional: suplente a 15 min y decisión dentro de 30 min, sin actuar por silencio. Dueño: JP; verifica CAL. \\
R-028 & S=4; O=3; D=4; NPR=48. & Acuerdo de finalidades y controles pendiente a D-5 de la ola, u objeción formal al uso operacional propuesto.. Dado que existe objeción sindical a cámaras en cabina y al control de jornada por GPS, aunque la propuesta excluye cámaras en cabina y separa la telemetría operacional, entonces puede que no se acuerden o comprendan las finalidades y controles del uso propuesto antes de la ola de reparto, teniendo como consecuencia resistencia a la incorporación y pruebas o aceptaciones pendientes. & D0=4: la validación distribuida puede llegar después de la fecha necesaria. Control/acción: Acuerdo 9.7.2 con finalidades, accesos y controles comprendidos por Operaciones/TI/sindicato. & S1=4; O1=2; D1=2; NPR1=16. D1=2: acta o ensayo específico previo a habilitación; se invalida la mejora de detección si no existe evidencia. & D-5 sin acuerdo: JP escala el mismo día y conserva sólo uso operacional autorizado; no cámaras en cabina ni GPS de jornada no acordado. Dueño: JP; verifica CAL. \\
R-029 & S=4; O=4; D=4; NPR=64. & Preparador nuevo asignado sin tutor o evaluación, o aprendizaje que excede dos horas en el programa 9.5.2.. Dado que preparación nocturna tiene rotación anual del 38\% y requiere capacitación continua, entonces puede que las incorporaciones superen la capacidad de tutoría y evaluación o lleguen al flujo sin competencias verificadas, teniendo como consecuencia errores de preparación, refuerzo extraordinario y certificación de la ola pendiente. & D0=4: la validación distribuida puede llegar después de la fecha necesaria. Control/acción: Altas por turno, tutor, evaluación antes de autonomía y aprendizaje dentro de dos horas según 9.5.2. & S1=4; O1=3; D1=2; NPR1=24. D1=2: acta o ensayo específico previo a habilitación; se invalida la mejora de detección si no existe evidencia. & Plan adicional: Bodega asigna cobertura certificada antes del turno; IMP/SRE revisa cada alta diariamente. Si supera tutoría, no asignar trabajo autónomo; presentar refuerzo en dos días hábiles y revalidar antes de habilitar. Dueño: IMP; verifica CAL. \\
R-030 & S=4; O=3; D=5; NPR=60. & No hay contrato de interfaz o extracción verificado al cierre del mes 1, o el ensayo temprano contradice formato, acceso o comportamiento declarados.. Dado que el CLIENTE no tiene documentación técnica de las interfaces del ERP ni del WMS de 2013, que se reemplaza progresivamente, entonces puede que el levantamiento no permita verificar a tiempo contratos de intercambio y extracción del legado, teniendo como consecuencia reproceso de adaptadores o cargas y bloqueo de pruebas de integración o migración. & D0=5: faltante de origen o interfaz no verificable puede permanecer oculto hasta conciliar/aceptar. Control/acción: Contrato de extracción e intercambio confirmado con TI; ensayo temprano sin modificar ERP y sustitución progresiva de WMS. & S1=4; O1=3; D1=2; NPR1=24. D1=2: acta o ensayo específico previo a habilitación; se invalida la mejora de detección si no existe evidencia. & Plan adicional: ARQ eleva brecha al finalizar semana 2 del mes 1; TI entrega acceso o alternativa verificable antes de integración. Ensayar archivo controlado si satisface requisitos; bloquear certificación si persiste interfaz no verificable. Dueño: ARQ; verifica CAL. \\
R-031 & S=4; O=3; D=4; NPR=48. & Incidente alto/crítico, evidencia o certificación pendiente en las semanas finales del peak, sin ventana de intervención o validación disponible.. Dado que el peak de septiembre casi duplica entregas durante tres semanas y coincide con el cierre de la marcha blanca E2 en el calendario supuesto de S7, entonces puede que un incidente, una validación pendiente o la falta de disponibilidad de usuarios no pueda resolverse dentro de las ventanas autorizadas, teniendo como consecuencia incumplimiento de condiciones de cierre y necesidad de recuperación de la programación. & D0=4: la validación distribuida puede llegar después de la fecha necesaria. Control/acción: Disponibilidad y continuidad acordadas antes de septiembre; seguimiento diario de las seis condiciones de cierre. & S1=4; O1=3; D1=2; NPR1=24. D1=2: acta o ensayo específico previo a habilitación; se invalida la mejora de detección si no existe evidencia. & Plan adicional: cerrar habilitaciones en mes 19 antes del congelamiento; JP revisa incidentes diariamente en 19--20. Escalar el mismo día una condición incumplida y acordar recuperación en dos días hábiles, sin intervención prohibida ni aceptación ficticia. Dueño: JP; verifica CAL. \\
R-032 & S=4; O=3; D=4; NPR=48. & Aviso de soporte/compatibilidad que deja una versión necesaria fuera de la política aprobada o hace fallar la regresión de la actualización prevista.. Dado que la propuesta opera durante 56 meses con componentes como Laravel 13, PHP 8.5 y Proxmox, entonces puede que cambien su soporte o compatibilidad de forma que las actualizaciones ordinarias no basten, teniendo como consecuencia una adecuación extraordinaria y revalidación para mantener seguridad y continuidad. & D0=4: la validación distribuida puede llegar después de la fecha necesaria. Control/acción: Inventario de versiones y soporte, revisión trimestral y compatibilidad probada antes del límite verificado. & S1=4; O1=2; D1=2; NPR1=16. D1=2: acta o ensayo específico previo a habilitación; se invalida la mejora de detección si no existe evidencia. & Cambio excepcional: ARQ registra aviso en un día hábil y presenta transición en cinco; programar antes de fin de soporte, conservando ventanas y pruebas. Dueño: ARQ; verifica CAL. \\
R-035 & S=4; O=3; D=4; NPR=48. & Conductor asignado y habilitado rechaza el flujo digital o entrega sin la evidencia requerida, sin que la causa sea firma pendiente o reemplazo sin aviso.. Dado que los conductores de terceros no son trabajadores de Puelche y requieren incorporación acordada con su transportista, entonces puede que un conductor autorizado y convocado rechace utilizar el terminal o el flujo digital aun existiendo acuerdo de la empresa, teniendo como consecuencia entregas sin evidencia digital completa y cobertura de adopción de la ola pendiente. & D0=4: la validación distribuida puede llegar después de la fecha necesaria. Control/acción: Demostración práctica en andén e incorporación acordada con transportista y conductor habilitado. & S1=4; O1=2; D1=2; NPR1=16. D1=2: acta o ensayo específico previo a habilitación; se invalida la mejora de detección si no existe evidencia. & Ante rechazo actuar antes de salida: IMP solicita conductor certificado/dispuesto ese turno y registra cobertura pendiente; refuerzo antes de siguiente asignación. Dueño: IMP; verifica CAL. \\
\end{tablalafrox}
\end{landscape}
\endgroup

\clearpage
\section{Anexo 8.H: Reservas, perfiles y programación del esfuerzo}
\label{sec:cierre-reservas}
La exposición inicial de los 36 riesgos es de 1.199,52 HH, calculada con las probabilidades y los esfuerzos condicionales sin redondear del Anexo 8.B y redondeada sólo al final. Los controles técnicos de diseño, innovación y calidad forman parte del trabajo ordinario; una diferencia de esfuerzo requiere estimación y aprobación específica.

La propuesta de contingencia específica es la suma p1 H0 de los 36 eventos: incluye atención contingente C y remanente L=H0-C. La prevención A es trabajo planificado adicional y se solicita como tal, separado de reserva. No se suman otra vez la exposición inicial, p1L y T=A+p1C a p1H0. La estimación corresponde al valor esperado; no determina percentiles ni suficiencia estadística de la reserva. Las hipótesis p1 dependen de evidencia; si no se verifica, se retiene la exposición inicial para replanificar y escalar, sin gastar una reserva proyectada como si ya estuviera aprobada.

La reserva protegida de S7, Tabla 7.4, pp. 45--46, es 3.072 HH: 1.152 en marcha blanca E1 (13--15) y 1.920 en producción/estabilización E1 (16--20). Ya integra 203.548,62 HH programadas sobre 200.476,62 HH base/cobertura. No se suma nuevamente a esos totales ni se presta a E2. El crédito inicial contra la contingencia específica es cero y sólo podrá reconocerse tras verificar la coincidencia de capacidades, perfiles, causas y períodos. JP podrá reconocer solapamiento sólo con la triple correspondencia riesgo--paquete--perfil/mes y acta de capacidad; al hacerlo, se resta exactamente una vez de la propuesta adicional, conservando la protección E1. Las fichas preventivas también se descuentan si ya están incluidas en paquetes.

La Propuesta Final aprueba 160 HH de reserva de gestión para revisar, diagnosticar, coordinar y replanificar asuntos emergentes. Se asignan cuatro ciclos de 40 HH, uno por bloque, en los meses 12, 18, 24 y 40. Si el asunto corresponde a un riesgo del T-16, se imputa a la contingencia específica. JP replanifica cualquier esfuerzo que exceda el cupo del ciclo. Dos ciclos por bloque representarían 320 HH en el análisis de sensibilidad.

\paragraph{Asignación por titular y turno.}
En cada mes indicado, los turnos son diurnos, de ocho HH, en los primeros dos días hábiles del mes. La asignación se repite en los cuatro ciclos.

\begin{tablalafrox}{Capacidad de gestión asignada por ciclo.}{tab:gestion-turnos}{Y C{0.22\linewidth} G{0.16}}{\cab{Titular y rol} & \cab{Turno diurno} & \cab{HH por ciclo}}
Alex Aravena (JP) & Día hábil 1 & 8 \\
Bastián Trejo (ARQ) & Días hábiles 1 y 2 & 16 \\
Maximiliano Miño (CAL) & Día hábil 1 & 8 \\
Guillermo Castillo (SRE) & Día hábil 1 & 8 \\
\textbf{Total} & & \textbf{40} \\
\end{tablalafrox}

En los meses 12 y 18, PMO cubre 8 HH de trabajo delegable de JP en el primer turno. En los meses 24 y 40, un profesional DevOps habilitado cubre 8 HH operativas de SRE en ese turno. Para los demás titulares y turnos, la programación reserva primero el cupo indicado y distribuye la carga ordinaria dentro del saldo mensual. El tope de 128 HH programables deja hasta 112 HH ordinarias a ARQ y hasta 120 HH ordinarias a cada uno de los otros titulares en el mes del ciclo. Las coberturas reasignan trabajo de la línea base dentro de sus ventanas y las 160 HH se contabilizan como reserva de gestión.

La tabla mensual distribuye p1H0 uniformemente entre los meses de atención propuestos de cada ficha, 60\% al dueño y 40\% al apoyo (CAL, o ARQ si CAL es dueño). No predice fechas de ocurrencia ni cambia la ventana probabilística. Las cantidades de cálculo conservan precisión; el redondeo de impresión a 0,01 HH puede crear diferencias de centésimas al sumar celdas; los totales se calculan sobre los valores sin redondear. Prevención y gestión tienen sus propios calendarios y no están incluidas en esa tabla. La reserva distribuida es exposición esperada, no dotación comprometida: JP debe recalcular la ruta crítica y nivelar perfiles antes de solicitar su aprobación. No se convierten HH directamente a días.

\begingroup
\renewcommand{\thetable}{8.H.1}
\begin{landscape}
\setlength{\textwidth}{\linewidth}
\renewcommand{\arraystretch}{1.08}
\begin{tablalafrox}{Conciliación de esfuerzo y reservas, HH. Fuente: 8.B/8.F, S7 Tabla 7.4 y escenario explícito de gestión.}{tab:cierre-reserva}{F{0.28} G{0.14} Y}{\cab{Concepto} & \cab{HH} & \cab{Tratamiento y no duplicación}}
Exposición inicial vigente & 1.199,52 & Valoración previa al tratamiento; no sumar a la proyección posterior. \\
Trabajo preventivo adicional A & 180,00 & Trabajo planificado, 60/40 por ficha; descontar coincidencia con EDT antes de aprobar. \\
Contingencia específica propuesta p1H0 & 837,53 & Incluye C y L; condición: evidencia que sustenta p1. \\
Dentro de contingencia: p1C & 142,00 & Subconjunto; no reserva adicional. \\
Dentro de contingencia: p1L & 695,53 & Subconjunto de contingencia, no sumar otra vez. \\
Esfuerzo esperado posterior A+p1H0 & 1017,53 & Prevención más contingencia específica; sin gestión ni protección S7. \\
Gestión aprobada & 160,00 & Cuatro ciclos de 40 HH asignados por titular y turno. \\
Protección ya programada en S7 & 3.072,00 & Sólo E1; 1.152+1.920. Crédito demostrado inicialmente cero. \\
\end{tablalafrox}
\end{landscape}
\endgroup

\begingroup
\renewcommand{\thetable}{8.H.2}
\begin{landscape}
\setlength{\textwidth}{\linewidth}
\renewcommand{\arraystretch}{1.08}
\begin{tablalafrox}{Trabajo preventivo adicional propuesto por mes de primera atención y perfil (HH). Fuente: fichas 8.F.1; ejecutar antes del primer hito dependiente.}{tab:cierre-preventiva}{G{0.07} G{0.10} G{0.10} G{0.10} G{0.10} G{0.10} G{0.10} G{0.10}}{\cab{Mes} & \cab{JP} & \cab{ARQ} & \cab{DAT} & \cab{CAL} & \cab{SRE} & \cab{IMP} & \cab{SEG}}
1 & 8,40 & 18,00 & 22,80 & 48,40 & 2,40 & 15,60 & 2,40 \\
2 & 0,00 & 0,00 & 0,00 & 5,60 & 0,00 & 8,40 & 0,00 \\
13 & 0,00 & 8,40 & 0,00 & 13,60 & 0,00 & 12,00 & 0,00 \\
19 & 4,80 & 0,00 & 0,00 & 3,20 & 0,00 & 0,00 & 0,00 \\
20 & 3,60 & 0,00 & 0,00 & 2,40 & 0,00 & 0,00 & 0,00 \\
\end{tablalafrox}
\end{landscape}
\endgroup

\begingroup
\renewcommand{\thetable}{8.H.3}
\begin{landscape}
\setlength{\textwidth}{\linewidth}
\renewcommand{\arraystretch}{1.08}
\begin{tablalafrox}{Distribución mensual propuesta de contingencia específica (HH). Fuente: p1H0 y ventanas de 8.F.}{tab:cierre-mensual}{G{0.06} G{0.09} G{0.09} G{0.09} G{0.09} G{0.09} G{0.09} G{0.09} G{0.10}}{\cab{Mes} & \cab{JP} & \cab{ARQ} & \cab{DAT} & \cab{CAL} & \cab{SRE} & \cab{IMP} & \cab{SEG} & \cab{Total}}
1 & 1,01 & 7,03 & 7,32 & 12,60 & 0,13 & 2,98 & 0,05 & 31,13 \\
2 & 1,01 & 7,03 & 7,32 & 13,05 & 0,13 & 3,65 & 0,05 & 32,24 \\
3 & 1,01 & 2,87 & 7,32 & 10,28 & 0,13 & 3,65 & 0,05 & 25,31 \\
4 & 1,01 & 2,87 & 7,32 & 10,28 & 0,13 & 3,65 & 0,05 & 25,31 \\
5 & 1,01 & 2,87 & 7,32 & 10,28 & 0,13 & 3,65 & 0,05 & 25,31 \\
6 & 1,01 & 2,87 & 7,32 & 10,28 & 0,13 & 3,65 & 0,05 & 25,31 \\
7 & 1,01 & 2,87 & 7,32 & 10,28 & 0,13 & 3,65 & 0,05 & 25,31 \\
8 & 1,01 & 2,87 & 7,32 & 10,28 & 0,13 & 3,65 & 0,05 & 25,31 \\
9 & 1,01 & 2,87 & 7,32 & 10,28 & 0,13 & 3,65 & 0,05 & 25,31 \\
10 & 1,01 & 2,87 & 7,32 & 10,28 & 0,13 & 3,65 & 0,05 & 25,31 \\
11 & 1,01 & 2,87 & 7,32 & 10,28 & 0,13 & 3,65 & 0,05 & 25,31 \\
12 & 1,01 & 2,87 & 7,32 & 10,28 & 0,13 & 3,65 & 0,05 & 25,31 \\
13 & 1,01 & 5,20 & 7,32 & 12,88 & 0,19 & 5,17 & 0,05 & 31,82 \\
14 & 1,01 & 5,20 & 7,32 & 12,88 & 0,19 & 5,17 & 0,05 & 31,82 \\
15 & 1,01 & 5,20 & 7,32 & 12,88 & 0,19 & 5,17 & 0,05 & 31,82 \\
16 & 1,01 & 5,20 & 7,32 & 12,33 & 0,19 & 4,34 & 0,05 & 30,45 \\
17 & 1,01 & 5,20 & 7,32 & 12,33 & 0,19 & 4,34 & 0,05 & 30,45 \\
18 & 1,01 & 5,20 & 7,32 & 12,33 & 0,19 & 4,34 & 0,05 & 30,45 \\
19 & 8,71 & 3,56 & 5,36 & 15,07 & 0,19 & 4,34 & 0,05 & 37,30 \\
20 & 8,96 & 3,56 & 5,36 & 15,24 & 0,19 & 4,34 & 0,05 & 37,71 \\
21 & 0,75 & 1,30 & 0,00 & 2,88 & 0,19 & 2,03 & 0,05 & 7,20 \\
22 & 0,75 & 1,30 & 0,00 & 2,88 & 0,19 & 2,03 & 0,05 & 7,20 \\
23 & 0,75 & 1,30 & 0,00 & 2,88 & 0,19 & 2,03 & 0,05 & 7,20 \\
24 & 0,75 & 1,30 & 0,00 & 2,88 & 0,19 & 2,03 & 0,05 & 7,20 \\
25 & 0,75 & 1,30 & 0,00 & 2,88 & 0,19 & 2,03 & 0,05 & 7,20 \\
26 & 0,75 & 1,30 & 0,00 & 2,88 & 0,19 & 2,03 & 0,05 & 7,20 \\
27 & 0,75 & 1,30 & 0,00 & 2,88 & 0,19 & 2,03 & 0,05 & 7,20 \\
28 & 0,75 & 1,30 & 0,00 & 2,88 & 0,19 & 2,03 & 0,05 & 7,20 \\
29 & 0,75 & 1,30 & 0,00 & 2,88 & 0,19 & 2,03 & 0,05 & 7,20 \\
30 & 0,75 & 1,30 & 0,00 & 2,88 & 0,19 & 2,03 & 0,05 & 7,20 \\
31 & 0,75 & 1,30 & 0,00 & 2,88 & 0,19 & 2,03 & 0,05 & 7,20 \\
32 & 0,75 & 1,30 & 0,00 & 2,88 & 0,19 & 2,03 & 0,05 & 7,20 \\
33 & 0,75 & 1,30 & 0,00 & 2,88 & 0,19 & 2,03 & 0,05 & 7,20 \\
34 & 0,75 & 1,30 & 0,00 & 2,88 & 0,19 & 2,03 & 0,05 & 7,20 \\
35 & 0,75 & 1,30 & 0,00 & 2,88 & 0,19 & 2,03 & 0,05 & 7,20 \\
36 & 0,75 & 1,30 & 0,00 & 2,88 & 0,19 & 2,03 & 0,05 & 7,20 \\
37 & 0,75 & 1,30 & 0,00 & 2,88 & 0,19 & 2,03 & 0,05 & 7,20 \\
38 & 0,75 & 1,30 & 0,00 & 2,88 & 0,19 & 2,03 & 0,05 & 7,20 \\
39 & 0,75 & 1,30 & 0,00 & 2,88 & 0,19 & 2,03 & 0,05 & 7,20 \\
40 & 0,75 & 1,30 & 0,00 & 2,88 & 0,19 & 2,03 & 0,05 & 7,20 \\
41 & 0,75 & 1,30 & 0,00 & 2,88 & 0,19 & 2,03 & 0,05 & 7,20 \\
42 & 0,75 & 1,30 & 0,00 & 2,88 & 0,19 & 2,03 & 0,05 & 7,20 \\
43 & 0,75 & 1,30 & 0,00 & 2,88 & 0,19 & 2,03 & 0,05 & 7,20 \\
44 & 0,75 & 1,30 & 0,00 & 2,88 & 0,19 & 2,03 & 0,05 & 7,20 \\
45 & 0,75 & 1,30 & 0,00 & 2,88 & 0,19 & 2,03 & 0,05 & 7,20 \\
46 & 0,75 & 1,30 & 0,00 & 2,88 & 0,19 & 2,03 & 0,05 & 7,20 \\
47 & 0,75 & 1,30 & 0,00 & 2,88 & 0,19 & 2,03 & 0,05 & 7,20 \\
48 & 0,75 & 1,30 & 0,00 & 2,88 & 0,19 & 2,03 & 0,05 & 7,20 \\
49 & 0,75 & 1,30 & 0,00 & 2,88 & 0,19 & 2,03 & 0,05 & 7,20 \\
50 & 0,75 & 1,30 & 0,00 & 2,88 & 0,19 & 2,03 & 0,05 & 7,20 \\
51 & 0,75 & 1,30 & 0,00 & 2,88 & 0,19 & 2,03 & 0,05 & 7,20 \\
52 & 0,75 & 1,30 & 0,00 & 2,88 & 0,19 & 2,03 & 0,05 & 7,20 \\
53 & 0,75 & 1,30 & 0,00 & 2,88 & 0,19 & 2,03 & 0,05 & 7,20 \\
54 & 0,75 & 1,30 & 0,00 & 2,88 & 0,19 & 2,03 & 0,05 & 7,20 \\
55 & 0,75 & 1,30 & 0,00 & 2,88 & 0,19 & 2,03 & 0,05 & 7,20 \\
56 & 0,75 & 1,30 & 0,00 & 2,88 & 0,19 & 2,03 & 0,05 & 7,20 \\
\end{tablalafrox}
\end{landscape}
\endgroup
\paragraph{Programación y autorización} JP asigna las fichas preventivas y la contingencia a los paquetes de migración/datos, implantación, integración, soporte, gobierno o salida que atienden cada evento. Antes de autorizar capacidad adicional debe verificar paquete, perfil, mes y capacidad existente, descontando toda coincidencia. La prevención se programa en el mes de primera atención y antes del hito dependiente; la contingencia se nivela con la distribución mensual propuesta. Las 3.072 HH protegidas mantienen su uso exclusivo en E1. Los ciclos de gestión se sitúan en los meses 1, 13, 21 y 39 como capacidad de decisión; sólo se consumen ante un asunto emergente autorizado. El trabajo operativo de INN-03 a cargo de SRE desde el mes 21 conserva la imputación a DAT; el despliegue por rutas de INN-05 a su cargo desde el mes 22 conserva la imputación a IMP. El relevo mantiene una única imputación.

La programación está sujeta a nivelación y autorización de la línea base. JP verifica precedencias, holguras y ruta crítica, preservando congelamientos y las cuatro semanas finales de evidencia. Si la capacidad no permite cumplir, eleva refuerzo o cambio antes del hito, sin desplazarlo unilateralmente. La valoración económica utiliza la misma diferencia aprobada por mes y perfil, multiplicada por la tarifa neta validada. Gestión y protección no se registran como costo consumido por el solo hecho de reservarlas. JP presenta la conciliación al arranque, la actualiza en cada comité quincenal y antes de H5/H10; el Comité Ejecutivo autoriza los cambios y CAL verifica la ausencia de doble imputación. Una propuesta de reserva no acredita financiamiento aprobado.

R-036 añade 16,40 HH a la contingencia específica, distribuidas uniformemente entre meses 1--56: 9,84 HH de familia ARQ y 6,56 HH de CAL en toda la ventana. SRE asume la gestión operativa desde mes 22 sin crear una segunda imputación ni cambiar la familia ARQ. A=0; C=16 HH forma parte de H0. El dimensionamiento, las pruebas CT-14 y la corrección del alcance comprometido conservan su presupuesto ordinario.


\section*{Referencias}
\addcontentsline{toc}{section}{Referencias}
\begingroup
\small\setlength{\parindent}{-0.5in}\setlength{\leftskip}{0.5in}\setlength{\parskip}{3pt}
Distribuidora Puelche S.A. (2026a). \emph{Bases Administrativas de Licitación N.º TFEP-01/2026: Contratación de Solución Integral de Software y Servicios de Operación} (Formulario T-16).


LafroX SpA. (2026a). \emph{Propuesta técnica: Subdocumento 4, Arquitectura lógica, arquitectura física y data center}.

LafroX SpA. (2026b). \emph{Propuesta técnica: Subdocumento 7, Plan de trabajo}, y formularios T-14 y T-15.

LafroX SpA. (2026d). \emph{Propuesta técnica: Subdocumento 13, Innovaciones}.

LafroX SpA. (2026c). \emph{Propuesta técnica: Subdocumento 9, Plan de calidad}, y formularios T-13 y T-17.

Distribuidora Puelche S.A. (2026b). \emph{Bases Técnicas del Caso 02: Logística} (caps. 5--7, 10, 14--19).


\endgroup

\end{document}
